\chapter{Flat base change for the pushforward of a quasi-coherent sheaf ($i=0$)}
\label{chap:Cohomology_FlatBaseChange}
% archon:covers AlgebraicJacobian/Cohomology/FlatBaseChange.lean AlgebraicJacobian/Cohomology/FlatBaseChangeGlobal.lean
% SOURCE: [Stacks Project], Cohomology of Schemes, Section "Cohomology and
% base change, I", and Tag 02KH (read from references/stacks-coherent.tex and
% references/stacks-coherent.md).

\section{Setup and motivation}

This chapter establishes the $i=0$ (i.e.\ \(H^0\) / direct-image) case of flat
base change: the formation of the pushforward \(f_*\mathcal{F}\) of a
quasi-coherent sheaf commutes with flat base change on the target. It is the
foundational, fully self-contained slice of the cohomology-and-base-change
package that ultimately underlies relative representability statements; the
higher derived cases \(R^i f_*\) for \(i \geq 1\) are treated separately and are
out of scope here.

Throughout, \(f : X \to S\) is a morphism of schemes that is quasi-compact and
quasi-separated, and \(\mathcal{F}\) is a quasi-coherent \(\mathcal{O}_X\)-module.
A base change of \(f\) along an arbitrary morphism \(g : S' \to S\) is recorded
by the cartesian square obtained from the fibre product
\(X' = X \times_S S'\):
\[
  \begin{array}{ccc}
    X' & \xrightarrow{\;g'\;} & X \\[2pt]
    {\scriptstyle f'}\big\downarrow & & \big\downarrow{\scriptstyle f} \\[2pt]
    S' & \xrightarrow{\;g\;} & S
  \end{array}
\]
Here \(f' : X' \to S'\) is the base change of \(f\), and \(g' : X' \to X\) is the
first projection. We write \(\mathcal{F}' = (g')^*\mathcal{F}\) for the pullback
of \(\mathcal{F}\) to \(X'\).

\section{The canonical base-change map}

\begin{definition}\leanok
[Base-change map for the pushforward]
  \label{def:pushforward_base_change_map}
  \lean{AlgebraicGeometry.pushforwardBaseChangeMap}
  % SOURCE: [Stacks Project], Cohomology of Schemes, Section "Cohomology and
  % base change, I", base-change diagram (\ref{equation-base-change-diagram})
  % (read from references/stacks-coherent.tex, L877--904).
  % SOURCE QUOTE: "Let $f : X \to S$ be a morphism of schemes.
  % Let $\mathcal{F}$ be a quasi-coherent sheaf on $X$.
  % Suppose further that $g : S' \to S$ is any morphism of schemes. Denote
  % $X' = X_{S'} = S' \times_S X$ the base change of $X$ and denote
  % $f' : X' \to S'$ the base change of $f$.
  % Also write $g' : X' \to X$ the projection,
  % and set $\mathcal{F}' = (g')^*\mathcal{F}$."
  \textit{Source: Stacks Project, Cohomology of Schemes, \S\,Cohomology and base
  change, I.}
  Let \(f : X \to S\) be a morphism of schemes, \(\mathcal{F}\) a quasi-coherent
  \(\mathcal{O}_X\)-module, and \(g : S' \to S\) an arbitrary morphism, giving the
  cartesian square above with projection \(g' : X' \to X\) and
  \(\mathcal{F}' = (g')^*\mathcal{F}\). The \emph{base-change map for the
  pushforward} is the canonical morphism of \(\mathcal{O}_{S'}\)-modules
  \[
    g^*\bigl(f_*\mathcal{F}\bigr) \longrightarrow f'_*\bigl((g')^*\mathcal{F}\bigr)
    = f'_*\mathcal{F}' .
  \]
  It is adjoint to the morphism
  \(f_*\mathcal{F} \to g_*\, f'_*\mathcal{F}' = f_*\,(g')_*\mathcal{F}'\) obtained
  by applying \(f_*\) to the unit
  \(\mathcal{F} \to (g')_*(g')^*\mathcal{F}\) of the
  \(((g')^*, (g')_*)\)-adjunction, using the commutativity
  \(g \circ f' = f \circ g'\) of the square.
\end{definition}

\section{The affine case}

\subsection{Locality of isomorphisms for quasi-coherent sheaf morphisms}

Before treating the affine computation we record three project-local supplements
to Mathlib. Mathlib provides the per-open criterion
(\(\mathrm{IsIso}\,\varphi \iff \forall U,\ \mathrm{IsIso}(\varphi_U)\)) and a
stalkwise isomorphism criterion for \(\mathrm{TopCat.Sheaf}\)-valued morphisms, but
it does not package the stalk-local or basis-local criterion at the level of
morphisms of \(\mathcal{O}_X\)-modules. The following lemmas bridge that gap; they
are exactly the locality tools used in the first reduction of the affine
base-change lemma below, where one checks the base-change map after restricting to
the affine opens of the base. As bespoke infrastructure they carry no external
citation and stand on their proof sketches.

\begin{lemma}\leanok
[Stalk-local criterion for isomorphisms of \(\mathcal{O}_X\)-modules]
  \label{lem:modules_isIso_iff_stalk}
  \lean{AlgebraicGeometry.Modules.isIso_iff_isIso_stalkFunctor_map}
  Let \(X\) be a scheme and \(\varphi : M \to N\) a morphism of sheaves of
  \(\mathcal{O}_X\)-modules. Then \(\varphi\) is an isomorphism if and only if the
  underlying morphism of abelian-group presheaves induces an isomorphism on the
  stalk at every point \(x \in X\).
\end{lemma}

\begin{proof}



  If \(\varphi\) is an isomorphism, then so is its image under the forgetful
  functor to abelian-group presheaves, and the stalk functor at each \(x\),
  being a functor, carries isomorphisms to isomorphisms. Conversely, assume the
  stalk map is an isomorphism at every \(x\). Package the underlying abelian-group
  presheaves of \(M\) and \(N\) as sheaves of abelian groups, and \(\varphi\) as a
  morphism of these sheaves; the stalkwise-isomorphism hypothesis lets us apply the
  \(\mathrm{TopCat.Sheaf}\)-level stalk-local isomorphism criterion to conclude that
  this morphism of sheaves of abelian groups is an isomorphism. Pushing it through
  the forgetful functor to presheaves yields an isomorphism of the underlying
  abelian-group presheaf morphisms, and since the forgetful functor from
  \(\mathcal{O}_X\)-modules to abelian-group presheaves reflects isomorphisms,
  \(\varphi\) itself is an isomorphism.
\end{proof}

\begin{lemma}\leanok
[Basis-local criterion for isomorphisms of \(\mathcal{O}_X\)-modules]
  \label{lem:modules_isIso_of_isBasis}
  \lean{AlgebraicGeometry.Modules.isIso_of_isIso_app_of_isBasis}
  \uses{lem:modules_isIso_iff_stalk}
  Let \(X\) be a scheme, let \(\{B_i\}_{i \in \iota}\) be a family of opens whose
  range is a basis of the topology of \(X\), and let \(\varphi : M \to N\) be a
  morphism of sheaves of \(\mathcal{O}_X\)-modules. If \(\varphi\) restricts to an
  isomorphism on the sections over every basic open \(B_i\), then \(\varphi\) is an
  isomorphism.
\end{lemma}

\begin{proof}
  \uses{lem:modules_isIso_iff_stalk}
  By Lemma~\ref{lem:modules_isIso_iff_stalk} it suffices to prove that the induced
  stalk map is bijective at each point \(x \in X\). For injectivity, a germ that
  vanishes at \(x\) already vanishes on some open neighbourhood, which may be
  shrunk to a basic open \(B_i\); on \(B_i\) the section map \(\varphi_{B_i}\) is
  injective, so the original section is zero, giving injectivity of the stalk map
  from injectivity over the basis. For surjectivity, an arbitrary germ at \(x\) is
  represented by a section of \(N\) over some basic open \(B_i\) containing \(x\);
  since \(\varphi_{B_i}\) is surjective, this section lifts to a section of \(M\)
  over \(B_i\), whose germ at \(x\) maps to the given germ. Thus the stalk map is
  bijective, and \(\varphi\) is an isomorphism.
\end{proof}

\begin{lemma}\leanok
[Affine-open locality criterion for isomorphisms of \(\mathcal{O}_X\)-modules]
  \label{lem:modules_isIso_iff_affineOpens}
  \lean{AlgebraicGeometry.Modules.isIso_iff_isIso_app_affineOpens}
  \uses{lem:modules_isIso_of_isBasis}
  Let \(X\) be a scheme and \(\varphi : M \to N\) a morphism of sheaves of
  \(\mathcal{O}_X\)-modules. Then \(\varphi\) is an isomorphism if and only if it
  restricts to an isomorphism on the sections over every affine open of \(X\).
\end{lemma}

\begin{proof}
  \uses{lem:modules_isIso_of_isBasis}
  The forward implication is immediate, since the section functor over a fixed open
  preserves isomorphisms. For the converse, the affine opens of a scheme form a
  basis of its topology, so the claim is the special case of
  Lemma~\ref{lem:modules_isIso_of_isBasis} in which the basis is the family of
  affine opens.
\end{proof}

\subsection{The affine computation}

The whole result rests on the following elementary computation in the affine
setting, where the base-change map is visibly an isomorphism because tensoring is
associative. Before stating it we isolate the one ingredient not yet available as
a ready-made result on which the affine reduction turns: the description of the
pushforward of a tilde-module along a morphism of affine schemes.

We first record three bespoke \(\Gamma\)-fragment supports that compute the global
sections of an affine pushforward; they carry no external citation and stand on
their proof sketches.

\begin{lemma}

\leanok
[Naturality of the global-sections comparison square]
  \label{lem:globalSectionsIso_hom_comp_specMap_appTop}
  \lean{AlgebraicGeometry.globalSectionsIso_hom_comp_specMap_appTop}
  Let \(\varphi : R \to R'\) be a homomorphism of commutative rings, and write
  \(\mathrm{gs}_R : \Gamma(\operatorname{Spec} R, \top) \cong R\) and
  \(\mathrm{gs}_{R'} : \Gamma(\operatorname{Spec} R', \top) \cong R'\) for the
  global-sections isomorphisms of the structure sheaves. Then the ring square
  \[
    \mathrm{gs}_R \;\circ\; (\operatorname{Spec}\varphi)^{\sharp}_{\top}
    \;=\; \varphi \;\circ\; \mathrm{gs}_{R'}
  \]
  commutes, where \((\operatorname{Spec}\varphi)^{\sharp}_{\top}\) is the
  top-open component of the structure-sheaf comparison map underlying
  \(\operatorname{Spec}\varphi\).
\end{lemma}

\begin{proof}



  This is the naturality of the unit/counit of the
  \(\Gamma \dashv \operatorname{Spec}\) adjunction read off at the top open: the
  global-sections map of \(\operatorname{Spec}\varphi\) is, under the
  global-sections isomorphisms on each side, conjugate to \(\varphi\) itself.
  Transposing the naturality square of the \(\operatorname{Spec}\) unit gives
  exactly the asserted ring equation.
\end{proof}

\begin{lemma}

\leanok
[Global sections of an affine pushforward]
  \label{lem:gammaPushforwardIso}
  \lean{AlgebraicGeometry.gammaPushforwardIso}
  \uses{lem:globalSectionsIso_hom_comp_specMap_appTop}
  Let \(\varphi : R \to R'\) be a homomorphism of commutative rings and let
  \(N\) be an \(\operatorname{Spec}(R')\)-module. Then there is a canonical
  isomorphism of \(R\)-modules
  \[
    \Gamma\bigl((\operatorname{Spec}\varphi)_* N\bigr)
    \;\cong\;
    \operatorname{restr}_\varphi\bigl(\Gamma N\bigr),
  \]
  identifying the \(R\)-module of global sections of the pushforward with the
  restriction of scalars along \(\varphi\) of the \(R'\)-module of global sections
  of \(N\).
\end{lemma}

\begin{proof}
  \uses{lem:globalSectionsIso_hom_comp_specMap_appTop}
  Because \((\operatorname{Spec}\varphi)^{-1}(\top) = \top\), the global sections of
  the pushforward \((\operatorname{Spec}\varphi)_* N\) are, on underlying abelian
  groups, the global sections of \(N\) with no transport along the open-set
  identification. The two sides therefore carry the same underlying abelian group;
  both unwind to nested restriction-of-scalars towers along the structure-sheaf
  global-sections maps. These towers are reconciled by applying the
  restriction-of-scalars composition identity twice and substituting the ring
  equation of
  Lemma~\ref{lem:globalSectionsIso_hom_comp_specMap_appTop}, which shows the
  resulting \(R\)-action is exactly restriction of scalars along \(\varphi\).
\end{proof}

\begin{lemma}

\leanok
[Global sections of an affine pushforward of a tilde-module]
  \label{lem:gammaPushforwardTildeIso}
  \lean{AlgebraicGeometry.gammaPushforwardTildeIso}
  \uses{lem:gammaPushforwardIso}
  Let \(\varphi : R \to R'\) be a homomorphism of commutative rings and let \(M\)
  be an \(R'\)-module. Then there is a canonical isomorphism of \(R\)-modules
  \[
    \Gamma\bigl((\operatorname{Spec}\varphi)_* \widetilde{M}\bigr)
    \;\cong\;
    \operatorname{restr}_\varphi M .
  \]
\end{lemma}

\begin{proof}
  \uses{lem:gammaPushforwardIso}
  Specialise Lemma~\ref{lem:gammaPushforwardIso} to \(N = \widetilde{M}\) and
  compose with the tilde-unit isomorphism
  \(\Gamma(\widetilde{M}) \cong M\) of \(R'\)-modules (which restricts to an
  isomorphism of \(R\)-modules after applying \(\operatorname{restr}_\varphi\)).
\end{proof}

\begin{lemma}

\leanok
[Sections of an affine pushforward over an arbitrary open]
  \label{lem:gammaPushforwardIsoAt}
  \lean{AlgebraicGeometry.gammaPushforwardIsoAt}
  \uses{lem:globalSectionsIso_hom_comp_specMap_appTop}
  Let \(\varphi : R \to R'\) be a homomorphism of commutative rings, let \(N\) be
  an \(\operatorname{Spec}(R')\)-module, and let \(U\) be an open of
  \(\operatorname{Spec}(R)\), with preimage
  \(V = (\operatorname{Spec}\varphi)^{-1}(U)\) in \(\operatorname{Spec}(R')\). Then
  there is a canonical isomorphism of \(R\)-modules
  \[
    \Gamma\bigl((\operatorname{Spec}\varphi)_* N,\, U\bigr)
    \;\cong\;
    \operatorname{restr}_\varphi\bigl(\Gamma(N,\, V)\bigr) .
  \]
  This is the open-indexed generalization of
  Lemma~\ref{lem:gammaPushforwardIso}, which is the case \(U = \top\) (so that
  \(V = \top\) as well). The family \(\{e_U\}_U\) of these isomorphisms is moreover
  natural in the open argument: the structure-sheaf restriction maps on the two
  sides commute with the isomorphisms, so the family is a natural isomorphism of the
  two presheaves of \(R\)-modules in the open variable (this naturality is
  established in the proof below). For \(U = D(a)\) it is exactly the linear
  equivalence \(e_{D(a)}\) of movement~(1) in the proof of
  Lemma~\ref{lem:pushforward_spec_tilde_iso}.
\end{lemma}

\begin{proof}
  \uses{lem:globalSectionsIso_hom_comp_specMap_appTop}
  The construction copies that of Lemma~\ref{lem:gammaPushforwardIso} verbatim,
  with the top open \(\top\) replaced throughout by the arbitrary open \(U\) (and
  its preimage \(V\)). The point that makes this replacement free is that the
  \(R\)-module structure on the sections is supplied \emph{uniformly} by
  restriction of scalars along the global-sections ring map at the top open,
  independently of the evaluation open; consequently both sides unwind to the same
  nested restriction-of-scalars towers, reconciled by applying the
  restriction-of-scalars composition identity twice and substituting the
  \(\top\)-level ring equation of
  Lemma~\ref{lem:globalSectionsIso_hom_comp_specMap_appTop}. Naturality in the open
  follows because every constituent of the construction is either a structure-sheaf
  restriction map or conjugation by a fixed ring map, and all of these commute with
  further restriction along an inclusion of opens.

  \medskip
  \emph{Naturality of the family \(\{e_U\}_U\) in the open (remark).} The family is
  in fact a natural isomorphism between two presheaves of \(R\)-modules on
  \(\operatorname{Spec}(R)\): the source presheaf
  \(U \mapsto \Gamma\bigl((\operatorname{Spec}\varphi)_* N,\, U\bigr)\), with the
  structure-sheaf restriction maps of the pushforward, and the target presheaf
  \(U \mapsto
    \operatorname{restr}_\varphi\bigl(\Gamma(N,\, (\operatorname{Spec}\varphi)^{-1}U)\bigr)\),
  with the structure-sheaf restriction maps of \(N\) transported through the preimage
  operation \((\operatorname{Spec}\varphi)^{-1}\) and read as \(R\)-linear maps by
  restriction of scalars along \(\varphi\). Concretely, for any inclusion of opens
  \(U' \subseteq U\), with \(V = (\operatorname{Spec}\varphi)^{-1}U\) and
  \(V' = (\operatorname{Spec}\varphi)^{-1}U'\), the square with horizontal maps
  \(e_U, e_{U'}\) and vertical maps the respective structure-sheaf restrictions
  commutes. The reason is structural: each \(e_U\) is a composite whose constituents
  are of exactly two kinds — structure-sheaf restriction maps (of \(N\),
  respectively of \((\operatorname{Spec}\varphi)_* N\)) and identity-on-carrier
  \(\operatorname{restrictScalars}\) repackagings (conjugation by the \emph{fixed},
  open-independent ring map supplied by the restriction-of-scalars reconciliation and
  by the \(\top\)-level ring equation of
  Lemma~\ref{lem:globalSectionsIso_hom_comp_specMap_appTop}). Restriction maps commute
  with further restriction by the presheaf functoriality axiom (using
  \(V' = (\operatorname{Spec}\varphi)^{-1}U' \subseteq
  (\operatorname{Spec}\varphi)^{-1}U = V\), as \((\operatorname{Spec}\varphi)^{-1}\) is
  monotone), and conjugation by a fixed ring map commutes with restriction because the
  ring map does not depend on \(U\); the naturality square is the paste of these
  component squares, and each \(e_U\) being an isomorphism makes the family a natural
  isomorphism. This open-naturality is exactly the compatibility invoked, at the
  inclusion \(D(a) \subseteq \top\), in the proof of
  Lemma~\ref{lem:pushforward_spec_tilde_iso}.
\end{proof}

\begin{lemma}

\leanok
[The tilde restriction from the top open to a basic open is a localization]
  \label{lem:tildeRestriction_isLocalizedModule}
  \lean{AlgebraicGeometry.tildeRestriction_isLocalizedModule}
  Let \(M\) be an \(R'\)-module and \(b \in R'\). Then the structure-sheaf
  restriction map
  \[
    \Gamma(\widetilde M,\, \top) \longrightarrow \Gamma(\widetilde M,\, D(b)),
  \]
  read as a map of \(R'\)-modules, exhibits \(\Gamma(\widetilde M, D(b))\) as the
  localization of \(\Gamma(\widetilde M, \top)\) at \(\mathrm{powers}(b)\). This is
  the \(R'\)-side localization fact transported in movement~(3) of the proof of
  Lemma~\ref{lem:pushforward_spec_tilde_iso}.
\end{lemma}

\begin{proof}


  By the defining property of \(\widetilde{(-)}\), the tilde-localization map
  \(M \to \Gamma(\widetilde M, D(b))\) exhibits its target as the localization of
  \(M\) at \(\mathrm{powers}(b)\). The global-sections map
  \(M \to \Gamma(\widetilde M, \top)\) is the localization of \(M\) at the trivial
  submonoid \(\mathrm{powers}(1)\) — since \(D(1) = \top\) — and is therefore
  bijective. The triangle identity relating these two localization maps to the
  structure-sheaf restriction \(\Gamma(\widetilde M, \top) \to
  \Gamma(\widetilde M, D(b))\) then exhibits the restriction itself, post-composed
  with the inverse of the bijection \(M \cong \Gamma(\widetilde M, \top)\), as the
  localization of \(\Gamma(\widetilde M, \top)\) at \(\mathrm{powers}(b)\).
\end{proof}

We record three further bespoke supports that drive the basic-open verification of
the affine-pushforward isomorphism below. They carry no external citation and stand
on their proof sketches.

\begin{lemma}

\leanok
[Restriction of scalars of a localized module]
  \label{lem:powers_restrictScalars}
  \lean{AlgebraicGeometry.IsLocalizedModule.powers_restrictScalars}
  Let \(A\) be a commutative \(R\)-algebra, \(S \subseteq R\) a submonoid, and
  \(f : M \to N\) an \(A\)-linear map that exhibits \(N\) as the localization of
  \(M\) at the image submonoid
  \(\operatorname{algMap}_A(S) \subseteq A\). Then the underlying \(R\)-linear map
  \(\operatorname{restr}_R f\) exhibits \(N\) as the localization of \(M\) at \(S\)
  itself.
\end{lemma}

\begin{proof}



  This is the exact converse of the standard fact that a localization of \(M\) at a
  submonoid \(S\) of \(R\) remains a localization after enlarging the base to an
  \(R\)-algebra \(A\) along the image submonoid \(\operatorname{algMap}_A(S)\). One
  checks the three defining conditions of a localized module directly. Each
  \(s \in S\) acts invertibly on \(N\) because its image
  \(\operatorname{algMap}_A(s)\) does and the two actions coincide by the scalar
  tower \(R \to A \to N\). Surjectivity of \(f\) up to denominators and the
  equality-witness condition transport along the map
  \(s \mapsto \operatorname{algMap}_A(s)\) from \(S\) onto
  \(\operatorname{algMap}_A(S)\), again using
  \(\operatorname{algMap}_A(s)\cdot{} = s\cdot{}\). In the application
  \(A = R'\), \(S = \mathrm{powers}(a)\), one has
  \(\operatorname{algMap}_{R'}(\mathrm{powers}(a)) = \mathrm{powers}(\varphi a)\), so
  this identifies the localization of \(\operatorname{restr}_\varphi M\) at \(a\)
  with the localization of \(M\) at \(\varphi a\); it is the ring-change core of the
  basic-open computation below.
\end{proof}

\begin{lemma}

\leanok
[Section-level counit isomorphism from a localization hypothesis]
  \label{lem:fromTildeGamma_app_isIso_of_localized}
  \lean{AlgebraicGeometry.fromTildeΓ_app_isIso_of_isLocalizedModule}
  Let \(N\) be an \(\operatorname{Spec}(R)\)-module and \(a \in R\). Suppose the
  structure-sheaf restriction map \(\Gamma(N, \top) \to \Gamma(N, D(a))\), read as
  an \(R\)-linear map of the global-section modules, exhibits \(\Gamma(N, D(a))\) as
  the localization of \(\Gamma(N, \top)\) at \(\mathrm{powers}(a)\). Then the
  tilde--\(\Gamma\) counit
  \(\mathrm{fromTilde}\Gamma\,N : \widetilde{\Gamma(N,\top)} \to N\) restricts to an
  isomorphism on the sections over the basic open \(D(a)\).
\end{lemma}

\begin{proof}



  Over \(D(a)\) the section map \(L\) of the counit sits in the triangle identity
  \(L \circ j = \rho\), where
  \(j : \Gamma(N,\top) \to \Gamma\bigl(\widetilde{\Gamma(N,\top)}, D(a)\bigr)\) is
  the tilde-localization map and \(\rho : \Gamma(N,\top) \to \Gamma(N, D(a))\) is
  the structure-sheaf restriction. Both \(j\) and \(\rho\) exhibit their targets as
  the localization of \(\Gamma(N,\top)\) at \(\mathrm{powers}(a)\) — \(j\) by the
  defining property of \(\widetilde{(-)}\), and \(\rho\) by hypothesis. By the
  uniqueness of localized modules there is a unique \(R\)-linear isomorphism \(e\)
  with \(e \circ j = \rho\); the triangle identity then forces \(L = e\), so \(L\) is
  an isomorphism. (This is the per-basic-open input feeding the basis-local
  criterion Lemma~\ref{lem:modules_isIso_of_isBasis}.)
\end{proof}

\begin{lemma}

\leanok
[Affine pushforward of a tilde-module, conditional form]
  \label{lem:pushforward_spec_tilde_iso_conditional}
  \lean{AlgebraicGeometry.pushforward_spec_tilde_iso_of_isLocalizedModule}
  \uses{lem:modules_isIso_of_isBasis, lem:gammaPushforwardTildeIso,
  lem:fromTildeGamma_app_isIso_of_localized}
  Let \(\varphi : R \to R'\) be a homomorphism of commutative rings and \(M\) an
  \(R'\)-module. Suppose that for every \(a \in R\) the structure-sheaf restriction
  of the pushforward \(N := (\operatorname{Spec}\varphi)_*\widetilde M\) exhibits its
  sections over \(D(a)\) as the localization of its global sections at
  \(\mathrm{powers}(a)\). Then there is a canonical isomorphism of
  \(\operatorname{Spec}(R)\)-modules
  \[
    (\operatorname{Spec}\varphi)_*\widetilde M \;\cong\;
    \widetilde{\bigl(\operatorname{restr}_\varphi M\bigr)} .
  \]
\end{lemma}

\begin{proof}
  \uses{lem:modules_isIso_of_isBasis, lem:gammaPushforwardTildeIso, lem:fromTildeGamma_app_isIso_of_localized}
  Write \(N := (\operatorname{Spec}\varphi)_*\widetilde M\). Under the hypothesis,
  Lemma~\ref{lem:fromTildeGamma_app_isIso_of_localized} applies for each \(a \in R\),
  so the counit \(\mathrm{fromTilde}\Gamma\,N\) restricts to an isomorphism on the
  sections over every basic open \(D(a)\). Since the basic opens form a basis of
  \(\operatorname{Spec}(R)\), the basis-local criterion
  Lemma~\ref{lem:modules_isIso_of_isBasis} upgrades this to: the counit
  \(\mathrm{fromTilde}\Gamma\,N\) is an isomorphism, i.e.\ \(N\) lies in the
  essential image of \(\widetilde{(-)}\). Composing the inverse counit with the
  tilde of the global-sections comparison
  Lemma~\ref{lem:gammaPushforwardTildeIso} yields the asserted object isomorphism.
\end{proof}

\begin{lemma}\leanok
[Affine pushforward of a tilde-module]
  \label{lem:pushforward_spec_tilde_iso}
  \lean{AlgebraicGeometry.pushforward_spec_tilde_iso}
  \uses{lem:modules_isIso_of_isBasis, lem:gammaPushforwardTildeIso,
        lem:pushforward_spec_tilde_iso_conditional, lem:powers_restrictScalars,
        lem:gammaPushforwardIso, lem:gammaPushforwardIsoAt,
        lem:tildeRestriction_isLocalizedModule}
% NOTE: the carrier-wall obstruction to `hloc` is resolved (restriction-of-scalars /
% scalar-tower instances feed lem:powers_restrictScalars; lem:tildeRestriction_isLocalizedModule
% supplies the R'-side localization). The single remaining obligation is the open-naturality of
% the section comparison e_{D(a)}, established as a prose remark in the proof of
% lem:gammaPushforwardIsoAt; the proof sketch below derives each per-a hloc(a) from the top-level
% localization plus that naturality square for the inclusion D(a) ⊆ ⊤, so no section-level square
% is hand-proved per a.
  Let \(\varphi : R \to R'\) be a homomorphism of commutative rings and let \(M\)
  be an \(R'\)-module. Then there is a canonical isomorphism of
  \(\operatorname{Spec}(R)\)-modules
  \[
    \bigl(\operatorname{Spec}\varphi\bigr)_*\,\widetilde{M}
    \;\cong\;
    \widetilde{\bigl(\operatorname{restr}_\varphi M\bigr)},
  \]
  where \(\operatorname{Spec}\varphi : \operatorname{Spec}(R') \to
  \operatorname{Spec}(R)\) is the induced morphism of affine schemes,
  \(\widetilde{(-)}\) denotes the quasi-coherent sheaf associated to a module, and
  \(\operatorname{restr}_\varphi M\) is the \(R\)-module obtained from \(M\) by
  restriction of scalars along \(\varphi\). That is, pushing the quasi-coherent
  sheaf \(\widetilde{M}\) forward along \(\operatorname{Spec}\varphi\) is, up to
  this canonical identification, the tilde of the restriction of scalars of \(M\).

  This is the classical fact that the direct image of a quasi-coherent sheaf along
  a morphism of affine schemes is computed on modules by restriction of scalars; it
  is the affine companion of the pullback formula (the affine case of the
  pushforward analogue of Stacks ``$\widetilde{M}$ and pullback''). It is recorded
  here as bespoke project infrastructure because the affine-pushforward
  identification does not appear as a ready-made result in the existing
  quasi-coherent sheaf theory.
\end{lemma}

\begin{proof}
  \uses{lem:modules_isIso_of_isBasis, lem:gammaPushforwardTildeIso, lem:pushforward_spec_tilde_iso_conditional, lem:powers_restrictScalars, lem:gammaPushforwardIso, lem:gammaPushforwardIsoAt, lem:tildeRestriction_isLocalizedModule}
  We build the isomorphism \emph{directly on a basis of basic opens}. This is the
  non-circular route: checking the comparison map on basic opens yields the object
  isomorphism, and quasi-coherence of the pushforward is then a corollary, not a
  prerequisite. (One must not reverse this order: the global-sections comparison
  \(\Gamma\bigl((\operatorname{Spec}\varphi)_*\widetilde{M}\bigr) \cong
  \operatorname{restr}_\varphi M\) of
  Lemma~\ref{lem:gammaPushforwardTildeIso} alone would only upgrade to an object
  isomorphism once the pushforward is already known to be quasi-coherent, so
  quasi-coherence cannot be deduced from an object isomorphism built that way.)

  \emph{The comparison morphism.} There is a canonical morphism of
  \(\operatorname{Spec}(R)\)-modules
  \[
    \alpha : \widetilde{\bigl(\operatorname{restr}_\varphi M\bigr)}
    \longrightarrow \bigl(\operatorname{Spec}\varphi\bigr)_*\widetilde{M},
  \]
  namely the tilde-adjoint of the global-sections identification of
  Lemma~\ref{lem:gammaPushforwardTildeIso}. We show \(\alpha\) is an isomorphism.

  \medskip
  \emph{Reduction to basic opens.} The basic opens \(D(a)\), for \(a \in R\), form a
  basis of the topology of \(\operatorname{Spec}(R)\). By the basis-local criterion
  Lemma~\ref{lem:modules_isIso_of_isBasis} it suffices to check that \(\alpha\)
  restricts to an isomorphism on the sections over each basic open \(D(a)\).

  \medskip
  \emph{The computation over \(D(a)\).} On the source, the sections of
  \(\widetilde{(\operatorname{restr}_\varphi M)}\) over \(D(a)\) are the
  localization \(\bigl(\operatorname{restr}_\varphi M\bigr)[1/a]\) as an
  \(R[1/a]\)-module. On the target, since \((\operatorname{Spec}\varphi)^{-1}D(a) =
  D(\varphi a)\), the sections of \(\bigl(\operatorname{Spec}\varphi\bigr)_*
  \widetilde{M}\) over \(D(a)\) are the sections of \(\widetilde{M}\) over
  \(D(\varphi a)\), i.e.\ the localization \(M[1/\varphi a]\), regarded as an
  \(R[1/a]\)-module by restriction of scalars along the induced map
  \(R[1/a] \to R'[1/\varphi a]\). These two localizations agree: the element \(a\)
  acts on \(\operatorname{restr}_\varphi M\) through \(\varphi a\), so localizing
  \(\operatorname{restr}_\varphi M\) at the multiplicative set generated by \(a\) is
  the same module as localizing \(M\) at the multiplicative set generated by
  \(\varphi a\). Both are the universal localization of
  \(M\) inverting \(\varphi a\), so the comparison map \(\alpha\) is an isomorphism
  on the sections over \(D(a)\). As this holds for every basic open, \(\alpha\) is
  an isomorphism.

  \medskip
  \emph{The explicit formalization route over \(D(a)\).} The agreement of the two
  localizations just described is the entire content, but it must be realized
  through the global ring \(R\), because the \(R\)-action on the pushforward
  sections over \(D(a)\) is built from the global-sections module of
  \(N := (\operatorname{Spec}\varphi)_*\widetilde M\) rather than directly from
  \(R[1/a]\). The reduction therefore proceeds exactly as the conditional
  assembly Lemma~\ref{lem:pushforward_spec_tilde_iso_conditional} prescribes: it
  suffices to supply, for each \(a \in R\), the hypothesis \(\mathrm{hloc}(a)\) that
  the structure-sheaf restriction \(\Gamma(N, \top) \to \Gamma(N, D(a))\) exhibits
  \(\Gamma(N, D(a))\) as the localization of \(\Gamma(N, \top)\) at
  \(\mathrm{powers}(a)\). The \(R\)-action on \(\Gamma(N, D(a))\) is the honest
  restriction-of-scalars action carried by the scalar tower \(R \to R' \to {-}\)
  determined by \(\varphi\); against that structure the ring-change converse
  Lemma~\ref{lem:powers_restrictScalars} is applicable, so the carrier of the
  \(R\)-module is exactly the one to which the localization claim refers, and no
  recursively-defined or ad-hoc section-level scalar action intervenes.
  We obtain \(\mathrm{hloc}(a)\) by transporting the already-known \(R'\)-side
  localization (Lemma~\ref{lem:tildeRestriction_isLocalizedModule}), in
  element-free fashion, across the open-indexed section comparison of
  Lemma~\ref{lem:gammaPushforwardIsoAt} \emph{using its naturality in the open}
  (the naturality remark in the proof of
  Lemma~\ref{lem:gammaPushforwardIsoAt}). This proceeds in three
  movements.

  \emph{(1) The \(D(a)\)-level linear equivalence.} This is the component at
  \(U = D(a)\) of the open-indexed family of
  Lemma~\ref{lem:gammaPushforwardIsoAt}: the \(R\)-linear isomorphism
  \[
    e_{D(a)} : \Gamma\bigl((\operatorname{Spec}\varphi)_*\widetilde M,\, D(a)\bigr)
      \;\cong\;
      \operatorname{restr}_\varphi\bigl(\Gamma(\widetilde M,\, D(\varphi a))\bigr) ,
  \]
  using \((\operatorname{Spec}\varphi)^{-1}D(a) = D(\varphi a)\), which is
  definitional. No fresh \(D(a)\)-level construction is needed: this is precisely the
  open-indexed isomorphism already supplied by
  Lemma~\ref{lem:gammaPushforwardIsoAt} evaluated at the basic open \(D(a)\).

  \emph{(2) The naturality square for \(D(a) \subseteq \top\).} The compatibility
  that relates \(e_{D(a)}\) to the \(\top\)-level isomorphism
  \(e_\top\) (the global-sections comparison of
  Lemma~\ref{lem:gammaPushforwardIso}) is precisely the open-naturality square of the
  family \(\{e_U\}_U\) (established in the proof of
  Lemma~\ref{lem:gammaPushforwardIsoAt}), instantiated at the inclusion
  \(D(a) \subseteq \top\): the square
  \[
    \begin{array}{ccc}
      \Gamma\bigl((\operatorname{Spec}\varphi)_*\widetilde M,\, \top\bigr)
        & \xrightarrow{\;e_\top\;}
        & \operatorname{restr}_\varphi\bigl(\Gamma(\widetilde M,\, \top)\bigr) \\[4pt]
      {\scriptstyle \mathrm{res}_{\top \to D(a)}}\big\downarrow & &
        \big\downarrow{\scriptstyle \mathrm{res}_{\top \to D(\varphi a)}} \\[4pt]
      \Gamma\bigl((\operatorname{Spec}\varphi)_*\widetilde M,\, D(a)\bigr)
        & \xrightarrow{\;e_{D(a)}\;}
        & \operatorname{restr}_\varphi\bigl(\Gamma(\widetilde M,\, D(\varphi a))\bigr)
    \end{array}
  \]
  commutes, using \((\operatorname{Spec}\varphi)^{-1}D(a) = D(\varphi a)\), which is
  definitional. Because the family \(\{e_U\}_U\) is a natural isomorphism, this single
  square is available uniformly in \(a\); there is no per-\(a\) section-level identity
  to re-prove by hand. It intertwines the source restriction
  \(\mathrm{res}_{\top \to D(a)}\) on the pushforward with the (restriction-of-scalars
  reading of the) target restriction \(\mathrm{res}_{\top \to D(\varphi a)}\) on
  \(\widetilde M\).

  \emph{(3) Discharging \(\mathrm{hloc}(a)\) by transport.} The target restriction
  \(\mathrm{res}_{\top \to D(\varphi a)} : \Gamma(\widetilde M, \top) \to
  \Gamma(\widetilde M, D(\varphi a))\) exhibits \(\Gamma(\widetilde M, D(\varphi a)) =
  M[1/\varphi a]\) as the \(\mathrm{powers}(\varphi a)\)-localization of
  \(M = \Gamma(\widetilde M, \top)\) over \(R'\); this is exactly
  Lemma~\ref{lem:tildeRestriction_isLocalizedModule}. Apply the ring-change converse
  Lemma~\ref{lem:powers_restrictScalars} (with \(R\)-algebra \(A = R'\) and submonoid
  \(\mathrm{powers}(a)\), so that
  \(\operatorname{algMap}_{R'}(\mathrm{powers}(a)) = \mathrm{powers}(\varphi a)\)) —
  its \(\operatorname{Algebra} R\, R'\) and scalar-tower hypotheses being supplied by
  the restriction-of-scalars structure along \(\varphi\) — to read this same map as
  the \(\mathrm{powers}(a)\)-localization over \(R\). The commuting naturality square
  of movement (2) intertwines this target restriction with the source restriction
  \(\mathrm{res}_{\top \to D(a)} : \Gamma(N, \top) \to \Gamma(N, D(a))\) by way of the
  isomorphisms \(e_\top\) and \(e_{D(a)}\); transporting the localization property
  across that square — via the transport-along-an-equivalence principle for localized
  modules — yields that \(\mathrm{res}_{\top \to D(a)}\) is the
  \(\mathrm{powers}(a)\)-localization on \(\Gamma(N, D(a))\) over \(R\), i.e.\
  \(\mathrm{hloc}(a)\). No section-level computation specific to the individual \(a\)
  is performed: the only \(a\)-dependent input is the instantiation of the single
  natural-isomorphism square at \(D(a)\). Feeding the family
  \(\{\mathrm{hloc}(a)\}_{a \in R}\) to the conditional assembly
  Lemma~\ref{lem:pushforward_spec_tilde_iso_conditional} produces the unconditional
  object isomorphism \(\alpha\).

  \medskip
  \emph{Naturality in the open drives the transport.} The comparison assembled
  here, \((\operatorname{Spec}\varphi)_* \circ \widetilde{(-)} \cong
  \widetilde{(-)} \circ \operatorname{restr}_\varphi\), is natural in \emph{both}
  arguments — in the module \(M\) and in the open \(U\). The naturality in the
  module is the usual functoriality; the naturality in the open is the naturality
  remark in the proof of Lemma~\ref{lem:gammaPushforwardIsoAt}: the family
  \(\{e_U\}_U\) of Lemma~\ref{lem:gammaPushforwardIsoAt} commutes with the
  structure-sheaf restriction maps on both sides. It is precisely this
  open-naturality that drives the localization transport \emph{uniformly}: the
  per-\(a\) localization fact \(\mathrm{hloc}(a)\) follows from the \(\top\)-level
  localization Lemma~\ref{lem:tildeRestriction_isLocalizedModule} together with that
  naturality square for the inclusion
  \(D(a) \hookrightarrow \top\), rather than from a section-level naturality square
  re-proved by hand for each individual \(a\). This is the single remaining
  obligation of the construction; the carrier of the \(R\)-action on the pushforward
  sections is the restriction-of-scalars structure along \(\varphi\), against which
  Lemma~\ref{lem:powers_restrictScalars} applies directly.

  The \(R\)-module structure on the \(R'\)-side sections to which the ring-change
  converse Lemma~\ref{lem:powers_restrictScalars} is applied is the honest
  restriction-of-scalars structure: \(R\) acts through \(\varphi\), i.e.\ via the
  \(R\)-algebra structure on \(R'\) determined by \(\varphi\) and the associated
  scalar tower \(R \to R' \to M[1/\varphi a]\). It is \emph{not} an ad-hoc,
  recursively-defined scalar action. Lemma~\ref{lem:powers_restrictScalars} is
  stated against exactly this restriction-of-scalars / scalar-tower structure, and
  it is under that structure that localizing \(\operatorname{restr}_\varphi M\) at
  \(a\) coincides with localizing \(M\) at \(\varphi a\).

  % NOTE: This lemma is the affine case of the general fact that the pushforward of
  % a quasi-coherent sheaf along an affine morphism is quasi-coherent. In current
  % Mathlib master that fact is the IsLocalizing / fromTildeGamma characterization
  % (AlgebraicGeometry.isIso_fromTildeΓ_pushforward); it post-dates the project's
  % pinned Mathlib, so the project proves it in-tree by the basic-open route above.
  % The in-tree proof remains valid and is not made redundant by a future Mathlib
  % bump.
  \medskip
  \emph{Quasi-coherence as a corollary.} The object isomorphism just constructed
  exhibits \(\bigl(\operatorname{Spec}\varphi\bigr)_*\widetilde{M} \cong
  \widetilde{(\operatorname{restr}_\varphi M)}\) as the tilde of a module. Tilde
  modules are quasi-coherent and quasi-coherence of \(\operatorname{Spec}(R)\)-modules
  is closed under isomorphism, so the pushforward
  \(\bigl(\operatorname{Spec}\varphi\bigr)_*\widetilde{M}\) is quasi-coherent. This
  conclusion is stated \emph{after} the object isomorphism, of which it is a
  consequence — no separate ``pushforward preserves quasi-coherent'' input is
  needed for the affine lane.
\end{proof}

\begin{remark}
  This single isomorphism discharges both inputs that the affine reduction below
  requires of the pushforward of a tilde-module.
  \begin{enumerate}
    \item \emph{The global-sections comparison.} Applying the global-sections
      (module-of-sections) functor to the isomorphism identifies the
      \(R\)-module of sections of
      \(\bigl(\operatorname{Spec}\varphi\bigr)_*\widetilde{M}\) with
      \(\operatorname{restr}_\varphi M\); this is exactly the \(\Gamma\)-fragment
      comparison used to recognise the section-level base-change map.
    \item \emph{Quasi-coherence of the pushforward.} The tilde of any module is
      quasi-coherent, and quasi-coherence of \(\operatorname{Spec}(R)\)-modules is
      closed under isomorphism; hence
      \(\bigl(\operatorname{Spec}\varphi\bigr)_*\widetilde{M}\) is quasi-coherent.
      In particular no separate ``pushforward preserves quasi-coherent'' theorem is
      needed for the affine lane.
  \end{enumerate}
  Because both inputs the affine reduction asks of the pushforward are corollaries
  of this single object isomorphism — the section-level module identification and
  the quasi-coherence of the pushforward — this lemma is the sole bespoke
  ingredient the affine lane requires.
\end{remark}

\subsection{The pullback companion}

The affine reduction below requires, alongside the pushforward dictionary, the
\emph{pullback} companion: the affine description of the inverse image of a
tilde-module. It is the affine case of the classical fact that the pullback of a
quasi-coherent sheaf is computed on modules by base change of scalars.

The pullback companion is built by uniqueness of left adjoints from the
per-object \(\Gamma\)-fragment comparison of
Lemma~\ref{lem:gammaPushforwardIso}, packaged as a natural isomorphism of
right-adjoint functors (Lemma~\ref{lem:gammaPushforwardNatIso}).

\begin{lemma}

\leanok
[Naturality packaging of the pushforward \(\Gamma\)-comparison]
  \label{lem:gammaPushforwardNatIso}
  \lean{AlgebraicGeometry.gammaPushforwardNatIso}
  \uses{lem:gammaPushforwardIso}
  The per-object comparison of Lemma~\ref{lem:gammaPushforwardIso} assembles into a
  natural isomorphism of the two right-adjoint functors
  \((\operatorname{Spec}\varphi)_* \circ \widetilde{(-)}\) and
  \(\widetilde{(-)} \circ \operatorname{restr}_\varphi\) on modules, exhibiting the
  pushforward dictionary as a natural transformation rather than merely a family of
  object isomorphisms. This is the input to the uniqueness-of-left-adjoints argument
  that produces the pullback companion Lemma~\ref{lem:pullback_spec_tilde_iso}.
\end{lemma}
\begin{proof}
  \uses{lem:gammaPushforwardIso}
  The components are the object isomorphisms of Lemma~\ref{lem:gammaPushforwardIso};
  naturality is the corresponding naturality square, which holds because every
  constituent of the comparison is identity-on-elements.
\end{proof}

\begin{lemma}

\leanok
[Affine pullback of a tilde-module]
  \label{lem:pullback_spec_tilde_iso}
  \lean{AlgebraicGeometry.pullback_spec_tilde_iso}
  \uses{lem:pushforward_spec_tilde_iso, lem:gammaPushforwardNatIso}
  % SOURCE: [Stacks Project], Schemes, Lemma `lemma-widetilde-pullback', Tag 01I9
  % (\S\,Quasi-coherent sheaves on affines) (read from references/stacks-schemes.tex,
  % L1241--1256).
  % SOURCE QUOTE: "Let
  % $(X, \mathcal{O}_X) = (\Spec(S), \mathcal{O}_{\Spec(S)})$,
  % $(Y, \mathcal{O}_Y) = (\Spec(R), \mathcal{O}_{\Spec(R)})$
  % be affine schemes.
  % Let $\psi : (X, \mathcal{O}_X) \to (Y, \mathcal{O}_Y)$ be a
  % morphism of affine schemes, corresponding to the ring map
  % $\psi^\sharp : R \to S$ (see Lemma \ref{lemma-category-affine-schemes}).
  % \begin{enumerate}
  % \item We have $\psi^* \widetilde M = \widetilde{S \otimes_R M}$
  % functorially in the $R$-module $M$.
  % \item We have $\psi_* \widetilde N = \widetilde{N_R}$ functorially
  % in the $S$-module $N$.
  % \end{enumerate}"
  \textit{Source: Stacks Project, Schemes, Lemma ``\(\widetilde M\) and pullback''
  (Tag 01I9), part (1).}
  Let \(\varphi : R \to R'\) be a homomorphism of commutative rings and let \(M\)
  be an \(R\)-module. Then there is a canonical isomorphism of
  \(\operatorname{Spec}(R')\)-modules
  \[
    \bigl(\operatorname{Spec}\varphi\bigr)^{*}\,\widetilde{M}
    \;\cong\;
    \widetilde{\bigl(R' \otimes_R M\bigr)},
  \]
  where \(\operatorname{Spec}\varphi : \operatorname{Spec}(R') \to
  \operatorname{Spec}(R)\) is the induced morphism of affine schemes,
  \(\widetilde{(-)}\) denotes the quasi-coherent sheaf associated to a module, and
  \(R' \otimes_R M\) is the base change of \(M\) along \(\varphi\) — the
  \(R'\)-module obtained by extension of scalars. The isomorphism is natural in
  \(M\). That is, pulling the quasi-coherent sheaf \(\widetilde{M}\) back along
  \(\operatorname{Spec}\varphi\) is, up to this canonical identification, the tilde
  of the base change of \(M\). It is the pullback companion of the affine
  pushforward dictionary Lemma~\ref{lem:pushforward_spec_tilde_iso}, and is part~(1)
  of the same Stacks lemma of which the pushforward formula is part~(2).
\end{lemma}

% SOURCE QUOTE PROOF: "The first assertion follows from the identification in
% Lemma \ref{lemma-compare-constructions}
% and the result of Modules, Lemma \ref{modules-lemma-restrict-quasi-coherent}."
\begin{proof}
  \uses{lem:pushforward_spec_tilde_iso}
  Pullback along \(\operatorname{Spec}\varphi\) is, on quasi-coherent sheaves, the
  extension-of-scalars operation \(- \otimes_R R'\) on modules transported through
  the \(\widetilde{(-)}\) / \(\operatorname{moduleSpec}\Gamma\) equivalence between
  modules and quasi-coherent sheaves. Concretely, the global sections of
  \(\bigl(\operatorname{Spec}\varphi\bigr)^{*}\widetilde{M}\) are the base change
  \(R' \otimes_R M\): the pullback functor \((\operatorname{Spec}\varphi)^{*}\) is
  left adjoint to the pushforward \((\operatorname{Spec}\varphi)_{*}\), and under
  the global-sections identification the latter is restriction of scalars along
  \(\varphi\) (Lemma~\ref{lem:pushforward_spec_tilde_iso}); the left adjoint of
  restriction of scalars is extension of scalars \(M \mapsto R' \otimes_R M\). The
  comparison map is the unit of the base-change adjunction, and the universal
  property of base change identifies it with the canonical isomorphism
  \(\bigl(\operatorname{Spec}\varphi\bigr)^{*}\widetilde{M} \cong
  \widetilde{(R' \otimes_R M)}\). Naturality in \(M\) is the naturality of that
  unit. (One may equally read this off the standard-open description: on
  \(D(\varphi a)\) the localization \((R' \otimes_R M)[1/a]\) is
  \(R'[1/\varphi a] \otimes_{R[1/a]} M[1/a]\), which is the base change of the
  sections of \(\widetilde M\) over \(D(a)\).)
\end{proof}

\subsection{The affine base-change lemma and its remaining obligations}

The affine base-change lemma below proceeds \emph{directly on global sections}: it
reduces, via the locality criterion
Lemma~\ref{lem:modules_isIso_iff_affineOpens} and the two affine dictionaries
Lemmas~\ref{lem:pushforward_spec_tilde_iso} and~\ref{lem:pullback_spec_tilde_iso},
to two obligations that are isolated here as named blocks. Both are extracted from
the proof of Stacks ``Affine base change'', and both are absent from the pinned
Mathlib. The first is the \emph{locality reduction}: restricting to the affine opens
of \(S'\) reduces the isomorphism claim to the affine--affine model
\(S = \operatorname{Spec}(R)\), \(S' = \operatorname{Spec}(R')\),
\(X = \operatorname{Spec}(A)\), \(\mathcal{F} = \widetilde{M}\). The second is the
\emph{section-level identification}: in that affine--affine model, the global-sections
incarnation \(\Gamma(\alpha)\) of the base-change map, transported through the two
proved tilde dictionaries, is the canonical cancellation isomorphism
\(\operatorname{cancelBaseChange}\) — an isomorphism with no flatness hypothesis. This
is the honest cut of the affine close: locality on \(S'\), then a pure computation on
the module of global sections. No separate abstract adjoint-mate identification is
performed at the level of sheaves; the mate is unwound \emph{on sections}, where it
becomes \(\operatorname{cancelBaseChange}\) directly.

\begin{lemma}

\leanok
[Affine-local compatibility of the base-change map]
  \label{lem:base_change_map_affine_local}
  \lean{AlgebraicGeometry.base_change_map_affine_local}
  \uses{def:pushforward_base_change_map, lem:modules_isIso_iff_affineOpens}
  % SOURCE: [Stacks Project], Cohomology of Schemes, Lemma
  % `lemma-affine-base-change' (\S\,Cohomology and base change, I), proof, the
  % "local on $S$ and $S'$" step (read from references/stacks-coherent.tex,
  % L920--926).
  % SOURCE QUOTE: "The statement is local on $S$ and $S'$. Hence we may
  % assume $X = \Spec(A)$, $S = \Spec(R)$,
  % $S' = \Spec(R')$ and $\mathcal{F} = \widetilde{M}$
  % for some $A$-module $M$."
  \textit{Source: Stacks Project, Cohomology of Schemes, Lemma ``Affine base
  change'' (proof, local-on-\(S\)-and-\(S'\) step).}
  Let \(f : X \to S\) be an affine morphism and \(\mathcal{F}\) a quasi-coherent
  \(\mathcal{O}_X\)-module, with a base-change square as above. The formation of the
  base-change map (Definition~\ref{def:pushforward_base_change_map}) is compatible
  with restriction to affine opens of \(S'\): for each affine open \(U \subseteq S'\), the
  restriction of the base-change map to sections over \(U\) agrees with the
  base-change map of the square obtained by restricting along \(U\). Consequently,
  by the affine-open locality criterion
  Lemma~\ref{lem:modules_isIso_iff_affineOpens}, the assertion that the base-change
  map is an isomorphism over an arbitrary base-change square \((S, S', X, X')\)
  reduces to the affine--affine case \(S = \operatorname{Spec}(R)\),
  \(S' = \operatorname{Spec}(R')\), \(X = \operatorname{Spec}(A)\),
  \(\mathcal{F} = \widetilde{M}\) for an \(A\)-module \(M\).

\end{lemma}

\begin{proof}
  \uses{def:pushforward_base_change_map, lem:modules_isIso_iff_affineOpens}
  By the affine-open locality criterion
  Lemma~\ref{lem:modules_isIso_iff_affineOpens}, the base-change map is an isomorphism
  of \(\mathcal{O}_{S'}\)-modules as soon as its restriction to the sections over each
  affine open \(U \subseteq S'\) is an isomorphism. The only content beyond the bare criterion is to identify that
  restriction with the base-change map of the affine--affine square cut out over
  \(U\). We obtain the identification by unfolding the definition of the map and
  pushing the restriction functor through each of its building blocks. The
  compatibility this needs is \emph{not} a new coherence theorem: it is the naturality
  of the adjunction transpose, together with the standard fact that pushforward
  commutes with restriction to an open of the base. We spell this out in three steps.

  \medskip
  \emph{Step 1: unfold the map.} By
  Definition~\ref{def:pushforward_base_change_map}, the base-change map is the image,
  under the \((g^*, g_*)\)-adjunction transpose — the natural bijection
  \(\operatorname{Hom}(g^*A, B) \cong \operatorname{Hom}(A, g_*B)\) — of the morphism
  \[
    f_*\mathcal{F} \xrightarrow{\,f_*(\eta_{\mathcal{F}})\,}
    f_*(g')_*(g')^*\mathcal{F} = g_*\,f'_*(g')^*\mathcal{F},
  \]
  where \(\eta_{\mathcal{F}} : \mathcal{F} \to (g')_*(g')^*\mathcal{F}\) is the unit of
  the \(((g')^*, (g')_*)\)-adjunction and the identification
  \(f_*(g')_* = g_*\,f'_*\) is the pseudofunctoriality of pushforward applied to the
  square's commutativity \(g \circ f' = f \circ g'\). Thus the map is, by
  construction, the transpose of the \(f_*\)-image of a unit — a composite of
  adjunction data and pushforward functors, with no ingredient beyond these.

  \medskip
  \emph{Step 2: restriction to \(U\) commutes with each building block.} Fix an affine
  open \(U \subseteq S'\) and write \(\rho_U\) for restriction to \(U\) —
  equivalently, the section functor \(\Gamma(U, -)\) on \(\mathcal{O}_{S'}\)-modules,
  or pullback along the open immersion \(U \hookrightarrow S'\). Restriction along
  \(U\) is a functor, and every ingredient of Step 1 is a natural transformation or a
  natural bijection; so applying \(\rho_U\) to the transpose-of-a-unit of Step 1
  returns the transpose-of-the-unit built from the restricted data. Concretely:
  \begin{itemize}
    \item[(a)] the unit \(\eta_{\mathcal{F}}\) of the \(((g')^*, (g')_*)\)-adjunction
      restricts to the unit of the corresponding adjunction on the restricted square,
      because units are natural in their object and pulling the cartesian square back
      along \(U \hookrightarrow S'\) yields again a cartesian square (with total space
      the open \((f')^{-1}(U)\) of \(X'\));
    \item[(b)] \(f_*\) of a restricted morphism is the restriction of \(f_*\) of the
      morphism, because pushforward commutes with restriction to an open of the base:
      for an open \(U\) of the base one has
      \((f_*\mathcal{G})|_U = (f|_{f^{-1}U})_*\bigl(\mathcal{G}|_{f^{-1}U}\bigr)\),
      naturally in \(\mathcal{G}\) — and likewise for \((g')_*\) and \(f'_*\);
    \item[(c)] the \((g^*, g_*)\)-adjunction transpose is a bijection natural in both
      variables, so it commutes with restricting its source and target along \(U\).
  \end{itemize}
  Chaining (a), (b), (c) — i.e.\ transporting the transpose-of-unit composite of
  Step 1 through the restriction functor \(\rho_U\) and replacing each block by its
  restricted counterpart — shows that the restriction of the base-change map to
  sections over \(U\) is itself the base-change map formed from the cartesian square
  restricted over \(U\). Choosing an affine open
  \(\operatorname{Spec}(R) \subseteq S\) containing the image \(g(U)\), affineness of
  \(f\) (local on the base and stable under base change) makes the total space
  restrict to an affine \(\operatorname{Spec}(A)\), and quasi-coherence of
  \(\mathcal{F}\) makes it restrict to a tilde-module \(\widetilde{M}\) there; so this
  restricted map is exactly the base-change map of the affine--affine square cut out
  over \(U = \operatorname{Spec}(R')\) and \(\operatorname{Spec}(R)\). The whole
  identification is thus definitional plus the naturalities (a)--(c); it introduces no
  obligation beyond them.

  \medskip
  \emph{Step 3: conclude the reduction.} By Step 2 the per-affine-open hypothesis
  required by the locality criterion Lemma~\ref{lem:modules_isIso_iff_affineOpens} —
  that the base-change map restricts to an isomorphism on sections over each affine
  open \(U \subseteq S'\) — is, under the identification of Step 2, precisely the
  affine--affine section assertion. That assertion is supplied, with \emph{no} flatness
  hypothesis, by the section-level computation
  Lemma~\ref{lem:pushforward_base_change_mate_cancelBaseChange}, which identifies the
  affine--affine base-change map on global sections with the cancellation isomorphism
  \(\operatorname{cancelBaseChange}\). Applying the locality criterion
  Lemma~\ref{lem:modules_isIso_iff_affineOpens} to this family of per-open
  isomorphisms shows that the base-change map is an isomorphism, as asserted.
\end{proof}

\subsection{The section-level base-change identity}

The section-level identification
Lemma~\ref{lem:pushforward_base_change_mate_cancelBaseChange} below is the single
highest-leverage obligation of the affine close. In the affine--affine model the
base-change map \(\theta\) is, after its domain and codomain are read through the two
proved tilde dictionaries, a concrete \(R'\)-linear map
\(R' \otimes_R M \to (R' \otimes_R A) \otimes_A M\), and the content of the lemma is
that this map is the inverse of the cancellation isomorphism
\(\operatorname{cancelBaseChange}\). Because the base-change map is \emph{defined} as the
transpose of the inner pushforward comparison \(\theta_{\mathrm{in}}\) under the
\((g^* \dashv g_*)\)-adjunction (\(g = \operatorname{Spec}\psi\)), its section-level value
is read off by the \emph{conjugate-counit calculus}: unfolding the transpose through the
counit form of the hom-set adjunction
(\(\operatorname{Adjunction.homEquiv\_counit}\)) factors \(\theta\) as
\(g^*(\theta_{\mathrm{in}})\) followed by the geometric counit \(\varepsilon_g\), and
conjugating this composite by the two tilde dictionaries
\(\Theta_{\mathrm{src}}\) (Lemma~\ref{lem:base_change_mate_domain_read}) and
\(\Theta_{\mathrm{tgt}}\) (Lemma~\ref{lem:base_change_mate_codomain_read}) replaces the
geometric counit by the \emph{algebraic} extension-of-scalars counit along \(\psi\). This
is the route the live formalization takes: the cancellation identity
Lemma~\ref{lem:pushforward_base_change_mate_cancelBaseChange} routes through the generator
trace Lemma~\ref{lem:base_change_mate_generator_trace}, which is a one-line corollary of the
section identity Lemma~\ref{lem:base_change_mate_section_identity}, whose single residual
obligation is the counit-conjugate coherence Lemma~\ref{lem:base_change_mate_gstar_transpose}.
The per-generator/element evaluation is \emph{not} available as a shortcut: evaluating the
opaque geometric composite on \(1 \otimes m\) produces no element-level normal form, so the
abstract conjugate calculus is the only tractable route. Two Mathlib results are recorded
first as dependency anchors.

\begin{lemma}[Pullback of a pair of affine \(\operatorname{Spec}\)-maps]
  \label{lem:pullbackSpecIso_mathlib}
  \lean{AlgebraicGeometry.pullbackSpecIso}
  \mathlibok
  \textit{Provided by Mathlib.}
  Let \(R \to S\) and \(R \to T\) be \(R\)-algebras. Then the fibre product of the
  induced morphisms \(\operatorname{Spec}(S) \to \operatorname{Spec}(R)\) and
  \(\operatorname{Spec}(T) \to \operatorname{Spec}(R)\) is canonically the spectrum of
  the tensor product,
  \[
    \operatorname{Limits.pullback}\bigl(\operatorname{Spec}(R\to S),\,
      \operatorname{Spec}(R\to T)\bigr)
    \;\cong\; \operatorname{Spec}\bigl(S \otimes_R T\bigr)
    \qquad (\texttt{AlgebraicGeometry.pullbackSpecIso}),
  \]
  and under this isomorphism the two cone legs are the \(\operatorname{Spec}\)-maps of
  the canonical tensor inclusions: \texttt{pullbackSpecIso\_hom\_fst} identifies the
  first projection with \(\operatorname{Spec}\) of
  \(\texttt{Algebra.TensorProduct.includeLeftRingHom} : S \to S \otimes_R T\), and
  \texttt{pullbackSpecIso\_hom\_snd} identifies the second projection with
  \(\operatorname{Spec}\) of
  \(\texttt{Algebra.TensorProduct.includeRight} : T \to S \otimes_R T\) (with companion
  forms \texttt{pullbackSpecIso\_inv\_fst}, \texttt{\_inv\_snd}).
\end{lemma}

\begin{lemma}[Cancellation of base change for the tensor tower]
  \label{lem:cancelBaseChange_mathlib}
  \lean{TensorProduct.AlgebraTensorModule.cancelBaseChange}
  \mathlibok
  \textit{Provided by Mathlib.}
  Let \(R \to A \to B\) be a tower of algebras, \(M\) a \(B\)-module that is also an
  \(A\)- and \(R\)-module compatibly, and \(N\) an \(R\)-module. Then there is a
  \(B\)-linear equivalence, carrying \emph{no} flatness hypothesis,
  \[
    \operatorname{cancelBaseChange} :
    M \otimes_A (A \otimes_R N) \;\xrightarrow{\ \sim\ }\; M \otimes_R N,
  \]
  with \(\operatorname{cancelBaseChange}\bigl(m \otimes_A (a \otimes_R n)\bigr)
  = (a \cdot m) \otimes_R n\) (\texttt{cancelBaseChange\_tmul}) and inverse
  \(\operatorname{cancelBaseChange}^{-1}(m \otimes_R n) = m \otimes_A (1 \otimes_R n)\)
  (\texttt{cancelBaseChange\_symm\_tmul}). Instantiated at the affine--affine model
  (\(M = R'\), \(B = R'\), \(N = M\)) this is the cancellation isomorphism
  \((R' \otimes_R A) \otimes_A M \cong R' \otimes_R M\) used below, read in either
  orientation.
\end{lemma}

\begin{lemma}[Pushout cancellation of base change (natively \(R'\)-linear)]
  \label{lem:isPushout_cancelBaseChange_mathlib}
  \lean{Algebra.IsPushout.cancelBaseChange}
  \mathlibok
  \textit{Provided by Mathlib.}
  Let \(R \to S\) and \(R \to A\) be commutative-ring maps and let \(B\) be an
  \(A\)-algebra and \(S\)-algebra making the square \(R \to S,\ R \to A,\ A \to B,\
  S \to B\) a pushout of commutative rings, i.e.\ \([\operatorname{Algebra.IsPushout} R S
  A B]\). Then for any \(A\)-module \(M\) there is a bundled \(S\)-linear equivalence
  \[
    \operatorname{cancelBaseChange} R\,S\,A\,B\,M :
    B \otimes_A M \;\xrightarrow{\ \sim\ }\; S \otimes_R M ,
  \]
  with \(\operatorname{cancelBaseChange}\) and its inverse computed on elementary tensors
  by \texttt{cancelBaseChange\_tmul} / \texttt{cancelBaseChange\_symm\_tmul}. The
  \(S\)-linearity is part of the bundled \(\operatorname{LinearEquiv}\): there is
  \emph{no} separate \texttt{map\_smul'} obligation to discharge. Instantiated at
  \(S = R'\), \(B = A \otimes_R R'\) — for which the pushout instance
  \(\operatorname{Algebra.IsPushout} R\,R'\,A\,(A \otimes_R R')\) is supplied for free by
  \(\operatorname{TensorProduct.isPushout'}\) — it is the natively \(R'\)-linear
  cancellation \((A \otimes_R R') \otimes_A M \cong R' \otimes_R M\) used as the core of
  the regrouping equivalence below.
\end{lemma}

\begin{lemma}


\leanok
[The pullback cone legs are the \(\operatorname{Spec}\)-maps of the tensor inclusions]
  \label{lem:pullback_fst_snd_specMap_tensor}
  \lean{AlgebraicGeometry.pullback_fst_snd_specMap_tensor}
  \uses{lem:pullbackSpecIso_mathlib}
  Let \(\psi : R \to R'\) and \(\varphi : R \to A\) be homomorphisms of commutative
  rings, and regard \(R'\) and \(A\) as \(R\)-algebras via \(\psi\) and \(\varphi\).
  The canonical isomorphism of Lemma~\ref{lem:pullbackSpecIso_mathlib} identifies the
  fiber product \(\operatorname{Spec}(A) \times_{\operatorname{Spec}(R)}
  \operatorname{Spec}(R')\) with \(\operatorname{Spec}(R' \otimes_R A)\). Under this
  identification the two cone legs are:
  \begin{align*}
    g' &\;=\; \operatorname{Spec}\bigl(A \to R' \otimes_R A,\; a \mapsto 1 \otimes a\bigr),\\
    f' &\;=\; \operatorname{Spec}\bigl(R' \to R' \otimes_R A,\; r' \mapsto r' \otimes 1\bigr),
  \end{align*}
  composed with the inverse of the canonical isomorphism.
\end{lemma}

\begin{proof}
  \uses{lem:pullbackSpecIso_mathlib}
  The identification of each cone leg follows from the companion identities of
  Lemma~\ref{lem:pullbackSpecIso_mathlib}. The only content beyond citing those
  identities is to match the given morphisms \(\operatorname{Spec}\varphi\) and
  \(\operatorname{Spec}\psi\) with the \(\operatorname{Spec}\)-maps of the algebra
  structures that \(\varphi\) and \(\psi\) define on \(A\) and \(R'\) — a definitional
  identification — after which the Mathlib identities apply verbatim. This identification
  of the cone legs is the one Mathlib-absent step of the section-level computation.
\end{proof}

\begin{lemma}


\leanok
[Domain read of the section-level base-change map]
  \label{lem:base_change_mate_domain_read}
  \lean{AlgebraicGeometry.base_change_mate_domain_read}
  \uses{lem:pushforward_spec_tilde_iso, lem:pullback_spec_tilde_iso}
  % SOURCE: [Stacks Project], Cohomology of Schemes, Lemma
  % `lemma-affine-base-change' (\S\,Cohomology and base change, I), proof, the
  % "use Schemes, Lemma widetilde-pullback to describe pullbacks and pushforwards"
  % step (read from references/stacks-coherent.tex, L927--938).
  % SOURCE QUOTE: "We use Schemes, Lemma \ref{schemes-lemma-widetilde-pullback}
  % to describe pullbacks and pushforwards of $\mathcal{F}$."
  \textit{Source: Stacks Project, Cohomology of Schemes, Lemma ``Affine base
  change'' (proof, ``describe pullbacks and pushforwards'' step).}
  In the affine--affine model (\(g = \operatorname{Spec}\psi\),
  \(f = \operatorname{Spec}\varphi\), \(\mathcal{F} = \widetilde M\) for an
  \(A\)-module \(M\)), the global sections of the domain \(g^*(f_*\widetilde M)\) of the
  base-change map are canonically \(R' \otimes_R M\) as an \(R'\)-module. Explicitly,
  the assembled isomorphism
  \[
    \Theta_{\mathrm{src}} :
    \Gamma\bigl(g^*(f_*\widetilde M)\bigr) \xrightarrow{\ \sim\ } R' \otimes_R M
  \]
  is the composite of the pushforward dictionary
  Lemma~\ref{lem:pushforward_spec_tilde_iso} (reading \(f_*\widetilde M\) as
  restriction of scalars \(\operatorname{restr}_\varphi M\)) followed by the pullback
  dictionary Lemma~\ref{lem:pullback_spec_tilde_iso} (reading \(g^*\) of a tilde as
  extension of scalars \(R' \otimes_R (-)\) along \(\psi\)).
\end{lemma}

\begin{proof}
  \uses{lem:pushforward_spec_tilde_iso, lem:pullback_spec_tilde_iso}
  The domain \(g^*(f_*\widetilde M)\) involves only the cospan maps
  \(f = \operatorname{Spec}\varphi\) and \(g = \operatorname{Spec}\psi\), each a genuine
  \(\operatorname{Spec}\)-map, so the two affine dictionaries apply directly with no
  pullback-leg identification needed. The pushforward dictionary
  Lemma~\ref{lem:pushforward_spec_tilde_iso} gives
  \(f_*\widetilde M = \widetilde{\operatorname{restr}_\varphi M}\) on global sections,
  i.e.\ \(M\) viewed as an \(R\)-module; the pullback dictionary
  Lemma~\ref{lem:pullback_spec_tilde_iso} gives
  \(g^*\widetilde N = \widetilde{R' \otimes_R N}\), so
  \(\Gamma\bigl(g^*(f_*\widetilde M)\bigr) = R' \otimes_R M\) as an \(R'\)-module.
  Composing the two dictionary isomorphisms in this order is \(\Theta_{\mathrm{src}}\).
\end{proof}

\begin{lemma}


\leanok
[Codomain read of the section-level base-change map]
  \label{lem:base_change_mate_codomain_read}
  \lean{AlgebraicGeometry.base_change_mate_codomain_read}
  \uses{lem:pullback_fst_snd_specMap_tensor, lem:pullback_spec_tilde_iso,
        lem:pushforward_spec_tilde_iso, lem:pullback_isEquivalence_of_iso}
  % SOURCE: [Stacks Project], Cohomology of Schemes, Lemma
  % `lemma-affine-base-change' (\S\,Cohomology and base change, I), proof, the
  % "Namely X' = Spec(R' \otimes_R A) and F' is the qcoh sheaf associated to
  % (R' \otimes_R A) \otimes_A M" step (read from references/stacks-coherent.tex,
  % L929--932).
  % SOURCE QUOTE: "Namely, $X' = \Spec(R' \otimes_R A)$ and
  % $\mathcal{F}'$ is the quasi-coherent sheaf associated
  % to $(R' \otimes_R A) \otimes_A M$."
  \textit{Source: Stacks Project, Cohomology of Schemes, Lemma ``Affine base
  change'' (proof, ``\(X' = \operatorname{Spec}(R' \otimes_R A)\)'' step).}
  In the affine--affine model, with \(f'\) and \(g'\) the canonical projections of the
  pullback square (Lemma~\ref{lem:pullback_fst_snd_specMap_tensor}), the global
  sections of the codomain \(f'_*(g')^*\widetilde M\) are canonically
  \((R' \otimes_R A) \otimes_A M\) as an \(R'\)-module. Explicitly, the assembled
  isomorphism
  \[
    \Theta_{\mathrm{tgt}} :
    \Gamma\bigl(f'_*(g')^*\widetilde M\bigr)
    \xrightarrow{\ \sim\ } (R' \otimes_R A) \otimes_A M
  \]
  is the composite of the pullback dictionary Lemma~\ref{lem:pullback_spec_tilde_iso}
  (for \((g')^*\)) followed by the pushforward dictionary
  Lemma~\ref{lem:pushforward_spec_tilde_iso} (for \(f'_*\)).
\end{lemma}

\begin{proof}
  \uses{lem:pullback_fst_snd_specMap_tensor, lem:pullback_spec_tilde_iso,
        lem:pushforward_spec_tilde_iso}
  Here the two functors \((g')^*\) and \(f'_*\) are taken along the pullback cone legs,
  which are \emph{not} a priori \(\operatorname{Spec}\)-maps; this is where
  Lemma~\ref{lem:pullback_fst_snd_specMap_tensor} enters. By
  Lemma~\ref{lem:pullback_fst_snd_specMap_tensor}, the first projection \(g'\) is the
  \(\operatorname{Spec}\)-map of the inclusion \(A \to R' \otimes_R A\), and the second
  projection \(f'\) is the \(\operatorname{Spec}\)-map of \(R' \to R' \otimes_R A\).
  With the legs so identified the affine dictionaries apply:
  the pullback dictionary Lemma~\ref{lem:pullback_spec_tilde_iso} along
  \(A \to R' \otimes_R A\) reads \((g')^*\widetilde M\) as the extension of scalars
  \((R' \otimes_R A) \otimes_A M\), and the pushforward dictionary
  Lemma~\ref{lem:pushforward_spec_tilde_iso} along \(R' \to R' \otimes_R A\) reads
  \(f'_*\) of it as restriction of scalars back to \(R'\). The result is
  \((R' \otimes_R A) \otimes_A M\) as an \(R'\)-module, and composing the two dictionary
  isomorphisms is \(\Theta_{\mathrm{tgt}}\). This matches the Stacks description
  \(X' = \operatorname{Spec}(R' \otimes_R A)\), with \(\mathcal{F}'\) the tilde of
  \((R' \otimes_R A) \otimes_A M\).
\end{proof}

\begin{lemma}

\leanok
[Codomain read assembled from the composite-adjunction comparison]
  \label{lem:base_change_mate_codomain_read_legs}
  \lean{AlgebraicGeometry.base_change_mate_codomain_read_legs}
  \uses{lem:base_change_mate_codomain_read, lem:pullback_spec_tilde_iso,
        lem:pushforward_spec_tilde_iso}
  % NOTE: this is the variable-legs codomain read in the project's `pullbackComp` form: the pullback
  % factor is the project pseudofunctor coherence `(pullbackComp e.hom (Spec ιA)).symm.app`, NOT the
  % abstract `leftAdjointCompIso`. The `leftAdjointCompIso` (conjugate-native) form is the separate
  % decl base_change_mate_codomain_read_legs_conj (conj-1a), which carries the two Mathlib anchors
  % lem:leftAdjointCompIso_mathlib and lem:conjugateEquiv_leftAdjointCompIso_inv_mathlib.
  It is the variable-legs form of the codomain read
  Lemma~\ref{lem:base_change_mate_codomain_read}. Fix \(\psi : R \to R'\),
  \(\varphi : R \to A\) and an \(A\)-module \(M\); let \(e : X' \xrightarrow{\sim}
  \operatorname{Spec}(A \otimes_R R')\) be the canonical isomorphism and
  \(\iota_A : A \to A \otimes_R R'\), \(\iota_{R'} : R' \to A \otimes_R R'\) the tensor
  inclusions. Reading the codomain \(f'_*(g')^*\widetilde M\) along the \emph{explicit composite}
  legs \(g' = e \circ \operatorname{Spec}\iota_A\) and \(f' = e \circ \operatorname{Spec}\iota_{R'}\)
  (carried as free variables together with the leg-equality hypotheses
  \(\mathit{hfst}/\mathit{hsnd}\)), there is a canonical isomorphism of \(R'\)-modules
  \[
    \Theta_{\mathrm{tgt}}^{\flat} :
    \Gamma\bigl(f'_*(g')^*\widetilde M\bigr)
    \xrightarrow{\ \sim\ }
    \operatorname{restr}_{\iota_{R'}}\bigl(\operatorname{ext}_{\iota_A} M\bigr)
    \;=\; (R' \otimes_R A) \otimes_A M ,
  \]
  assembled with its pullback factor taken as the project pseudofunctor pullback-composition
  coherence \(\mathrm{pullbackComp}_{e,\operatorname{Spec}\iota_A}\) — the square's commutativity
  entering through the pushforward congruence \(\mathrm{pushforwardCongr}\) and the pushforward
  composition coherence \(\mathrm{pushforwardComp}\) — followed by the affine pullback dictionary
  (Lemma~\ref{lem:pullback_spec_tilde_iso}) along \(\iota_A\) and the affine pushforward dictionary
  (Lemma~\ref{lem:pushforward_spec_tilde_iso}) along \(\iota_{R'}\). Carrying the legs as free
  variables (rather than the literal projections) is the structural device that lets them be
  eliminated by substitution on a well-typed motive. Instantiating at the
  literal projection legs \(g' = \operatorname{pr}_1\), \(f' = \operatorname{pr}_2\) recovers
  Lemma~\ref{lem:base_change_mate_codomain_read}.
\end{lemma}

\begin{proof}
  \uses{lem:base_change_mate_codomain_read, lem:pullback_spec_tilde_iso,
        lem:pushforward_spec_tilde_iso}
  The construction is the verbatim assembly of
  Lemma~\ref{lem:base_change_mate_codomain_read}, with the cone legs presented as free variables
  \(g', f'\) carrying their factorizations \(\mathit{hfst}/\mathit{hsnd}\) through
  \(e\) as explicit hypotheses. The pullback factor is the project pseudofunctor coherence
  \(\mathrm{pullbackComp}_{e,\operatorname{Spec}\iota_A}\); composing it with the two affine
  dictionaries (Lemmas~\ref{lem:pullback_spec_tilde_iso}, \ref{lem:pushforward_spec_tilde_iso})
  yields \(\Theta_{\mathrm{tgt}}^{\flat}\). Its agreement with
  Lemma~\ref{lem:base_change_mate_codomain_read} at the literal projection legs is the cone-leg
  factorization.
\end{proof}

% --- Supporting lemmas for the codomain read ---
% These helpers establish that pullback along a scheme isomorphism is an equivalence
% of module categories; the adjunction unit of an equivalence is an isomorphism,
% used in the codomain-read computation.

\begin{lemma}\leanok
[Pullback along an isomorphism is an equivalence of module categories]
  \label{lem:pullbackIsoEquivalenceOfIso}
  \lean{AlgebraicGeometry.pullbackIsoEquivalenceOfIso}
  Let \(f : X \to Y\) be an isomorphism of schemes. Then the pullback functor on
  \(\mathcal{O}\)-modules \(f^* : Y\text{-}\mathbf{Mod} \to X\text{-}\mathbf{Mod}\) is
  one half of an equivalence of categories
  \(Y\text{-}\mathbf{Mod} \simeq X\text{-}\mathbf{Mod}\), with quasi-inverse
  \((f^{-1})^*\).
\end{lemma}
\begin{proof}

  Assemble the equivalence from the functor \(f^*\) and the candidate inverse
  \((f^{-1})^*\). The unit and counit isomorphisms are built from the pseudofunctor
  coherences of pullback: \((f^{-1})^* \circ f^* \cong (f \circ f^{-1})^* = \mathrm{id}^*
  \cong \mathrm{id}\) and symmetrically, using
  \(f \circ f^{-1} = \mathrm{id}_Y\), \(f^{-1} \circ f = \mathrm{id}_X\)
  (composition-of-pullbacks, congruence under equal maps, and pullback-of-identity).
  The triangle identities are forced by these coherences.
\end{proof}

\begin{lemma}\leanok
[The pullback functor along an isomorphism is an equivalence]
  \label{lem:pullback_isEquivalence_of_iso}
  \lean{AlgebraicGeometry.pullback_isEquivalence_of_iso}
  \uses{lem:pullbackIsoEquivalenceOfIso}
  Consequently, for an isomorphism \(f : X \to Y\) the functor
  \(f^* : Y\text{-}\mathbf{Mod} \to X\text{-}\mathbf{Mod}\) is an equivalence of
  categories. In particular its adjunction unit is a natural isomorphism, so
  \((\eta_{f})_{W}\) is invertible on every object \(W\).
\end{lemma}
\begin{proof}

  The functor underlying the equivalence of
  Lemma~\ref{lem:pullbackIsoEquivalenceOfIso} is exactly \(f^*\), so \(f^*\) is an
  equivalence; an equivalence's unit is a natural isomorphism.
\end{proof}

% --- Section-level identity, three pieces: regroup equiv (pure tensor algebra, native
%     R'-linear via Algebra.IsPushout.cancelBaseChange) ← section identity (Γ(θ) is the
%     R'-base change of the algebraic unit) ← IsIso corollary (the pinned leaf). ---

\begin{lemma}\leanok
[Regrouping isomorphism for the section-level mate]
  \label{lem:base_change_mate_regroupEquiv}
  \lean{AlgebraicGeometry.base_change_mate_regroupEquiv}
  % NOTE: tensor-order convention. The Lean uses Mathlib's `pullbackSpecIso`
  % orientation `A ⊗[R] R'` (so the codomain is `(A ⊗_R R') ⊗_A M` and the generator
  % goes `r' ⊗ m ↦ (1 ⊗ r') ⊗ m`); the prose below writes `R' ⊗_R A`. The two are the
  % same module up to `TensorProduct.comm` — a faithful re-orientation, no content change.
  \uses{lem:isPushout_cancelBaseChange_mathlib}
  % SOURCE: [Stacks Project], Cohomology of Schemes, Lemma
  % `lemma-affine-base-change' (\S\,Cohomology and base change, I), proof, the
  % "boils down to the equality" step (read from references/stacks-coherent.tex,
  % L933--938).
  % SOURCE QUOTE: "Thus we see that the lemma boils down to the
  % equality
  % $$
  % (R' \otimes_R A) \otimes_A M = R' \otimes_R M
  % $$
  % as $R'$-modules."
  \textit{Source: Stacks Project, Cohomology of Schemes, Lemma ``Affine base
  change'' (proof, ``boils down to the equality'' step).}
  Let \(\psi : R \to R'\) and \(\varphi : R \to A\) be ring homomorphisms and \(M\) an
  \(A\)-module. There is a bundled \(R'\)-linear isomorphism
  \[
    \operatorname{regroup} :
    (R' \otimes_R A) \otimes_A M \;\xrightarrow{\ \sim\ }\; R' \otimes_R M ,
  \]
  the ``equality as \(R'\)-modules'' of the source, sending the generator
  \((r' \otimes 1) \otimes m \mapsto r' \otimes m\); equivalently its inverse sends
  \(r' \otimes m \mapsto (r' \otimes 1) \otimes m\). It carries \emph{no} flatness
  hypothesis.
\end{lemma}

\begin{proof}
  \uses{lem:isPushout_cancelBaseChange_mathlib}
  The underlying \(R'\)-linear equivalence is supplied, already as a bundled
  \(\operatorname{LinearEquiv}\) over \(R'\), by the natively \(R'\)-linear
  pushout-cancellation equivalence of
  Lemma~\ref{lem:isPushout_cancelBaseChange_mathlib}. Crucially, that helper bundles
  \(R'\)-linearity \emph{internally}: there is no hand-written scalar-compatibility
  obligation, no induction over elementary tensors, and in particular \emph{no} residual
  \(r' \cdot 0 = 0\) bookkeeping branch. The only work left to this proof is therefore to
  promote that \(R'\)-linear equivalence to a morphism of \(R'\)-module \emph{objects} in
  the exact extension/restriction-of-scalars packaging that the codomain and domain reads
  produce. We work in the \(A \otimes_R R'\) orientation
  (the same module as \(R' \otimes_R A\) up to the canonical commutativity isomorphism,
  per the tensor-order convention note).

  \emph{The lone residual obstruction is the \(A\)-action diamond, resolved at the object
  level.} The pushout-cancellation helper takes its relative tensor product over the
  \emph{canonical} \(A\)-algebra structure on \(A \otimes_R R'\)
  (multiplication in the left factor), whereas the carrier object produced by the
  dictionaries,
  the extension of scalars of \(M\) along the left inclusion \(\iota : A \to A \otimes_R R'\),
  tensors over the \(A\)-module structure on \(A \otimes_R R'\) induced by \(\iota\). These
  two \(A\)-module structures on
  \(A \otimes_R R'\) agree on elements but are distinct as data, so the two relative
  tensor products are distinct types even though their underlying abelian groups and
  their \(R'\)-actions coincide. This diamond is the \emph{only} thing the helper does
  not already discharge, and — because the \(R'\)-linearity is now entirely internal to
  the Mathlib equivalence — it is resolved purely at the carrier level by an
  identity-on-elements \(R'\)-linear bridge
  \[
    eT : (A \otimes_R R') \otimes_A^{\,\operatorname{restr}} M
    \;\xrightarrow{\ \sim\ }\;
    (A \otimes_R R') \otimes_A^{\,\mathrm{can}} M ,
  \]
  the identity map on generators, which is \(R'\)-linear because both sides carry the
  same \(R'\)-action (restriction of scalars along \(R' \to A \otimes_R R'\)).

  Composing the bridge \(eT\) with the natively \(R'\)-linear helper equivalence and
  promoting to an isomorphism of module objects yields the desired bundled
  \(R'\)-linear isomorphism of \(\operatorname{ModuleCat}(R')\) objects. Being a composite
  of linear equivalences it is an \(R'\)-linear isomorphism, and none of its constituents
  carries a flatness hypothesis.
\end{proof}

% --- Section-level identity, decomposed into three linked steps ---
%     The LHS unwinding of the base-change map on global sections splits at:
%       (1) lem:base_change_mate_unit_value  — the affine unit IS the algebraic unit;
%       (2) lem:base_change_mate_fstar_reindex — the pushforward pseudofunctor reindex;
%       (3) lem:base_change_mate_gstar_transpose — the (g^* ⊣ g_*) transpose on sections.

\begin{lemma}

\leanok
[The affine pullback--pushforward unit is the algebraic unit]
  \label{lem:base_change_mate_unit_value}
  \lean{AlgebraicGeometry.base_change_mate_unit_value}
  \uses{lem:pullback_spec_tilde_iso, lem:pushforward_spec_tilde_iso}
  % SOURCE: [Stacks Project], Cohomology of Schemes, Lemma `lemma-affine-base-change'
  % (\S\,Cohomology and base change, I), proof, the "boils down to the equality" step
  % (read from references/stacks-coherent.tex, L932--937).
  % SOURCE QUOTE: "Thus we see that the lemma boils down to the equality
  % $$ (R' \otimes_R A) \otimes_A M = R' \otimes_R M $$ as $R'$-modules."
  \textit{Source: Stacks Project, Cohomology of Schemes, Lemma ``Affine base change''
  (proof, ``boils down to the equality'' step).}
  Let \(\iota_A : A \to A \otimes_R R'\) be the canonical inclusion
  (\(a \mapsto a \otimes 1\)). The unit of the
  \(\bigl((\operatorname{Spec}\iota_A)^*, (\operatorname{Spec}\iota_A)_*\bigr)\)-adjunction
  evaluated at the affine sheaf \(\widetilde M\), read on the global sections over
  \(\operatorname{Spec} A\) through the tilde dictionaries
  (Lemma~\ref{lem:pullback_spec_tilde_iso}, Lemma~\ref{lem:pushforward_spec_tilde_iso}) and
  the tilde--\(\Gamma\) unit \(\widetilde{(-)}\Gamma\)-iso, equals the algebraic unit of the
  extension/restriction-of-scalars adjunction along \(\iota_A\), namely
  \(\eta_M : M \to (A \otimes_R R') \otimes_A M\), \(m \mapsto (1 \otimes 1) \otimes m\).
\end{lemma}
\begin{proof}
  \uses{lem:pullback_spec_tilde_iso, lem:pushforward_spec_tilde_iso}
  This is the affine, square-free heart of the computation: there is no pullback square and
  no base-change map yet, only the comparison of two units of two adjunctions related by the
  two tilde dictionaries. The geometric adjunction
  \(\bigl((\operatorname{Spec}\iota_A)^* \dashv (\operatorname{Spec}\iota_A)_*\bigr)\) on
  sheaves of modules transports, under the tilde dictionaries which identify
  \((\operatorname{Spec}\iota_A)^*\) of a tilde with extension of scalars
  (Lemma~\ref{lem:pullback_spec_tilde_iso}) and \((\operatorname{Spec}\iota_A)_*\) of a tilde
  with restriction of scalars (Lemma~\ref{lem:pushforward_spec_tilde_iso}), to the algebraic
  extension/restriction adjunction \((\iota_A^{\,*} \dashv \iota_{A\,*})\) on
  \(\operatorname{ModuleCat}\). The dictionaries being the comparison isomorphisms of these
  adjunctions, naturality identifies the two units; evaluating at \(M\) gives the stated
  value \(m \mapsto (1 \otimes 1) \otimes m\).
\end{proof}

\begin{definition}

\leanok
[The inner comparison value \(\rho\)]
  \label{def:base_change_mate_inner_value}
  \lean{AlgebraicGeometry.base_change_mate_inner_value}
  \uses{lem:base_change_mate_unit_value, lem:pullback_fst_snd_specMap_tensor}
  Let \(\psi : R \to R'\), \(\varphi : R \to A\) be ring homomorphisms, \(M\) an
  \(A\)-module, and \(\iota_A : A \to A \otimes_R R'\), \(\iota_{R'} : R' \to A \otimes_R R'\)
  the canonical tensor inclusions (Lemma~\ref{lem:pullback_fst_snd_specMap_tensor}).
  The \emph{inner comparison value} is the \(R\)-linear map
  \[
    \rho : M \longrightarrow (A \otimes_R R') \otimes_A M,
    \qquad m \longmapsto (1 \otimes 1) \otimes m,
  \]
  obtained as the restriction of scalars along \(\varphi\) of the algebraic unit
  \(\eta_M\) of the extension/restriction adjunction along \(\iota_A\)
  (Lemma~\ref{lem:base_change_mate_unit_value}), transported across the equality
  \(\iota_A \circ \varphi = \iota_{R'} \circ \psi\).

  Proved directly in Lean.
\end{definition}

% --- Mathlib dependency anchors for the leg-reindex engine (mate / conjugate calculus) ---
% These three are the categorical tools the leg-reindex engine
% lem:pullbackPushforward_unit_comp rests on; they are supplied verbatim by Mathlib.

\begin{lemma}[Unit-across-conjugate coherence]
  \label{lem:unit_conjugateEquiv_mathlib}
  \lean{CategoryTheory.unit_conjugateEquiv}
  \mathlibok
  \textit{Provided by Mathlib.}
  Let \(\mathrm{adj}_1 : L_1 \dashv R_1\) and \(\mathrm{adj}_2 : L_2 \dashv R_2\) be two
  adjunctions between the same pair of categories, with units \(\eta_1\) and \(\eta_2\), and let
  \(\alpha : L_2 \Rightarrow L_1\) be a natural transformation of the left adjoints. Write
  \(\operatorname{conj}(\mathrm{adj}_1, \mathrm{adj}_2)\) for the conjugation (mate) bijection,
  which carries \(\alpha\) to a natural transformation \(R_1 \Rightarrow R_2\) of the right
  adjoints. Then for every object \(c\),
  \[
    (\eta_1)_c \circ \bigl(\operatorname{conj}(\mathrm{adj}_1, \mathrm{adj}_2)\,\alpha\bigr)_{L_1 c}
      \;=\; (\eta_2)_c \circ R_2(\alpha_c) ,
  \]
  the transposed form of the conjugation relating the two units across the mate of \(\alpha\).
\end{lemma}

\begin{lemma}[Unit of a composite adjunction]
  \label{lem:comp_unit_app_mathlib}
  \lean{CategoryTheory.Adjunction.comp_unit_app}
  \mathlibok
  \textit{Provided by Mathlib.}
  Let \(\mathrm{adj}_1 : F \dashv G\) (with \(F : C \to D\)) and \(\mathrm{adj}_2 : H \dashv I\)
  (with \(H : D \to E\)) be composable adjunctions, with units \(\eta_1\) and \(\eta_2\). Then the
  unit of the composite adjunction
  \(\mathrm{adj}_1 . \mathrm{adj}_2 : (H \circ F) \dashv (G \circ I)\) is, at every object \(X\),
  \[
    \bigl(\mathrm{adj}_1 . \mathrm{adj}_2\bigr)\text{-unit}_X
      \;=\; (\eta_1)_X \circ G\bigl((\eta_2)_{F X}\bigr).
  \]
\end{lemma}

\begin{lemma}[The pullback--pushforward conjugate sends the pullback coherence to the pushforward coherence]
  \label{lem:conjugateEquiv_pullbackComp_inv_mathlib}
  \lean{AlgebraicGeometry.Scheme.Modules.conjugateEquiv_pullbackComp_inv}
  \mathlibok
  \textit{Provided by Mathlib.}
  Let \(f : X \to Y\) and \(g : Y \to Z\) be composable morphisms of schemes, and write
  \((h^{*} \dashv h_{*})\) for the pullback--pushforward adjunction on \(\mathcal{O}\)-modules
  attached to a morphism \(h\). Under the conjugation (mate) bijection between the composite
  adjunction \((g^{*}\dashv g_{*}) . (f^{*}\dashv f_{*})\) and the adjunction
  \(\bigl((f\circ g)^{*} \dashv (f\circ g)_{*}\bigr)\), the inverse pullback-composition coherence
  \((\mathrm{pullbackComp}_{f,g})^{\mathrm{inv}}\) is carried to the pushforward-composition
  coherence \((\mathrm{pushforwardComp}_{f,g})^{\mathrm{hom}}\):
  \[
    \operatorname{conj}\!\Bigl((g^{*}\dashv g_{*}) . (f^{*}\dashv f_{*}),\;
      \bigl((f\circ g)^{*}\dashv (f\circ g)_{*}\bigr)\Bigr)
      \bigl((\mathrm{pullbackComp}_{f,g})^{\mathrm{inv}}\bigr)
    \;=\; (\mathrm{pushforwardComp}_{f,g})^{\mathrm{hom}}.
  \]
\end{lemma}

% --- Mathlib dependency anchors for the conjugate-side re-encoding of the leg-reindex (Open Q2) ---
% These five categorical results are the tools of the conjugate-side (mate / Beck--Chevalley)
% proof route that replaces the abandoned direct-on-sections telescoping; all are supplied verbatim
% by Mathlib (CategoryTheory/Adjunction/Mates.lean and .../CompositionIso.lean).

\begin{lemma}[Conjugation is an equivalence on isomorphisms]
  \label{lem:conjugateIsoEquiv_mathlib}
  \lean{CategoryTheory.conjugateIsoEquiv}
  \mathlibok
  \textit{Provided by Mathlib.}
  Let \(\mathrm{adj}_1 : L_1 \dashv R_1\) and \(\mathrm{adj}_2 : L_2 \dashv R_2\) be two adjunctions
  between the same pair of categories. The conjugation (mate) bijection on natural transformations
  restricts to an equivalence on natural \emph{isomorphisms},
  \[
    \operatorname{conjIso}(\mathrm{adj}_1,\mathrm{adj}_2) :
      (L_2 \cong L_1) \;\simeq\; (R_1 \cong R_2),
  \]
  carrying an isomorphism of left adjoints to the conjugate isomorphism of right adjoints, and back. In
  particular the conjugate of an isomorphism is an isomorphism: this is the certificate that makes the
  affine flat-base-change comparison invertible without inspecting its concrete sheaf assembly.
\end{lemma}

\begin{lemma}[The iterated mate of a left-adjoint square is the conjugate (Beck--Chevalley)]
  \label{lem:iterated_mateEquiv_conjugateEquiv_mathlib}
  \lean{CategoryTheory.iterated_mateEquiv_conjugateEquiv}
  \mathlibok
  \textit{Provided by Mathlib.}
  For a square of functors whose four boundary functors are all left adjoints, with adjunctions
  \(\mathrm{adj}_1, \dots, \mathrm{adj}_4\), the iterated mate (the Beck--Chevalley \(2\)-cell of the
  pasted square) of a \(2\)-cell \(\alpha\) equals the conjugate of \(\alpha\) for the two composite
  adjunctions:
  \[
    \operatorname{mate}\bigl(\operatorname{mate}(\alpha)\bigr) \;=\;
      \operatorname{conj}\bigl(\mathrm{adj}_1.\mathrm{adj}_4,\ \mathrm{adj}_3.\mathrm{adj}_2\bigr)\,\alpha .
  \]
  Together with Lemma~\ref{lem:conjugateIsoEquiv_mathlib} this is the abstract reason a Beck--Chevalley
  transformation is a natural isomorphism exactly when \(\alpha\) is — the conceptual home of the
  iso-ness of the affine comparison \(g^* f_*\widetilde M \to f'_*(g')^*\widetilde M\), which the
  project realises directly as a conjugate of the pushforward \(\Gamma\)-comparison.
\end{lemma}

\begin{lemma}[Composite-adjunction comparison from a right-adjoint isomorphism]
  \label{lem:leftAdjointCompIso_mathlib}
  \lean{CategoryTheory.Adjunction.leftAdjointCompIso}
  \mathlibok
  \textit{Provided by Mathlib.}
  Let \(\mathrm{adj}_{01} : F_{01} \dashv G_{10}\), \(\mathrm{adj}_{12} : F_{12} \dashv G_{21}\) and
  \(\mathrm{adj}_{02} : F_{02} \dashv G_{20}\) be adjunctions, and let
  \(e : G_{21} \circ G_{10} \cong G_{20}\) be a natural isomorphism comparing the composite right adjoint
  with \(G_{20}\). Then there is a canonical natural isomorphism of the corresponding left adjoints
  \[
    \operatorname{leftAdjointCompIso}(e) : F_{01} \circ F_{12} \;\cong\; F_{02},
  \]
  namely the conjugate of \(e^{-1}\). It builds the comparison for a composite of two adjunctions out of
  a single right-adjoint comparison, carrying no propositional equality data.
\end{lemma}

\begin{lemma}[The conjugate of the composite-comparison inverse is the right-adjoint isomorphism]
  \label{lem:conjugateEquiv_leftAdjointCompIso_inv_mathlib}
  \lean{CategoryTheory.Adjunction.conjugateEquiv_leftAdjointCompIso_inv}
  \mathlibok
  \textit{Provided by Mathlib.}
  In the situation of Lemma~\ref{lem:leftAdjointCompIso_mathlib}, the conjugation bijection between the
  composite adjunction \(\mathrm{adj}_{01}.\mathrm{adj}_{12}\) and \(\mathrm{adj}_{02}\) sends the inverse
  of the composite comparison \(\operatorname{leftAdjointCompIso}(e)\) back to the hom of the
  right-adjoint isomorphism \(e\):
  \[
    \operatorname{conj}\bigl(\mathrm{adj}_{01}.\mathrm{adj}_{12},\ \mathrm{adj}_{02}\bigr)
      \bigl(\operatorname{leftAdjointCompIso}(e)^{-1}\bigr) \;=\; e^{\mathrm{hom}} .
  \]
  This is the one-line bridge that turns a coherence among composite-adjunction comparisons into an
  equation on the right-adjoint side, which the standard conjugate identities close.
\end{lemma}

\begin{lemma}[Transposed unit-across-conjugate coherence]
  \label{lem:unit_conjugateEquiv_symm_mathlib}
  \lean{CategoryTheory.unit_conjugateEquiv_symm}
  \mathlibok
  \textit{Provided by Mathlib.}
  Let \(\mathrm{adj}_1 : L_1 \dashv R_1\) and \(\mathrm{adj}_2 : L_2 \dashv R_2\) be adjunctions with
  units \(\eta_1,\eta_2\), and let \(\alpha : R_1 \Rightarrow R_2\) be a natural transformation of the
  right adjoints. Then for every object \(c\),
  \[
    (\eta_1)_c \,\circ\, \alpha_{L_1 c}
      \;=\; (\eta_2)_c \,\circ\, R_2\bigl(\bigl(\operatorname{conj}(\mathrm{adj}_1,\mathrm{adj}_2)^{-1}\alpha\bigr)_c\bigr),
  \]
  the transposed (inverse-conjugation) form of the unit-across-conjugate coherence
  Lemma~\ref{lem:unit_conjugateEquiv_mathlib}. This is the Seam~1 tool used one functor layer up in the
  conjugate-side leg-reindex: the cross-layer naturality of the pushforward \(\Gamma\)-comparison enters
  through this identity, not through a positional naturality rewrite.
\end{lemma}

\begin{lemma}[Transposed counit-across-conjugate coherence]
  \label{lem:conjugateEquiv_counit_symm_mathlib}
  \lean{CategoryTheory.conjugateEquiv_counit_symm}
  \mathlibok
  \textit{Provided by Mathlib.}
  Let \(\mathrm{adj}_1 : L_1 \dashv R_1\) and \(\mathrm{adj}_2 : L_2 \dashv R_2\) be adjunctions with
  counits \(\varepsilon_1,\varepsilon_2\), and let \(\alpha : L_2 \Rightarrow L_1\) be a natural
  transformation of the left adjoints. Then for every object \(d\),
  \[
    L_2\bigl(\alpha_d\bigr) \,\circ\, (\varepsilon_2)_d
      \;=\; \bigl(\bigl(\operatorname{conj}(\mathrm{adj}_1,\mathrm{adj}_2)^{-1}\alpha\bigr)\bigr)_d
        \,\circ\, (\varepsilon_1)_d ,
  \]
  the counit dual of the transposed unit-across-conjugate coherence
  Lemma~\ref{lem:unit_conjugateEquiv_symm_mathlib}. This is the master tool of the conjugate-counit
  discharge: instantiated at the composite adjunctions
  \(\mathrm{adj}_L = (\widetilde{(-)} \dashv \Gamma)_R \cdot ((\operatorname{Spec}\psi)^* \dashv
  (\operatorname{Spec}\psi)_*)\) and
  \(\mathrm{adj}_R = (\operatorname{ext}_\psi \dashv \operatorname{restr}_\psi) \cdot
  (\widetilde{(-)} \dashv \Gamma)_{R'}\) with \(\alpha = \beta^{\mathrm{hom}}\) the pushforward
  \(\Gamma\)-comparison (Lemma~\ref{lem:gammaPushforwardNatIso}), it transports the geometric counit
  \(\varepsilon_g\) (conjugated by the pullback dictionary
  Lemma~\ref{lem:pullback_spec_tilde_iso} and the tilde counit) to the algebraic
  extension/restriction-of-scalars counit along \(\psi\); the resulting identity is the landed master
  identity \(\mathrm{huce}\) consumed by Lemma~\ref{lem:base_change_mate_gstar_transpose}.
\end{lemma}

\begin{lemma}\leanok
[(conj-0) $\mathrm{pullbackComp}^{-1}$ is the $\operatorname{leftAdjointCompIso}$ of $\mathrm{pushforwardComp}$]
  \label{lem:pullbackComp_inv_eq_leftAdjointCompIso_inv}
  \lean{AlgebraicGeometry.pullbackComp_inv_eq_leftAdjointCompIso_inv}
  \uses{lem:leftAdjointCompIso_mathlib, lem:conjugateEquiv_leftAdjointCompIso_inv_mathlib,
        lem:conjugateEquiv_pullbackComp_inv_mathlib}
  Let \(f : X \to Y\) and \(g : Y \to Z\) be composable morphisms of schemes. The inverse of the
  project's pullback-composition pseudofunctor coherence
  \((\mathrm{pullbackComp}_{f,g})^{\mathrm{inv}}\) agrees with the inverse of the abstract
  composite-adjunction comparison \(\operatorname{leftAdjointCompIso}\)
  (Lemma~\ref{lem:leftAdjointCompIso_mathlib}) built from the pushforward coherence
  \(\mathrm{pushforwardComp}_{f,g}\):
  \[
    (\mathrm{pullbackComp}_{f,g})^{\mathrm{inv}}
      \;=\; \bigl(\operatorname{leftAdjointCompIso}(\mathrm{pushforwardComp}_{f,g})\bigr)^{\mathrm{inv}}.
  \]
  \begin{proof}
    \uses{lem:leftAdjointCompIso_mathlib, lem:conjugateEquiv_leftAdjointCompIso_inv_mathlib,
          lem:conjugateEquiv_pullbackComp_inv_mathlib}
    Both sides have the same image under the conjugation (mate) bijection
    \(\operatorname{conj}\bigl((g^{*}\dashv g_{*}).(f^{*}\dashv f_{*}),\ (f\circ g)^{*}\dashv (f\circ g)_{*}\bigr)\):
    the left side is sent to \((\mathrm{pushforwardComp}_{f,g})^{\mathrm{hom}}\) by the project coherence
    Lemma~\ref{lem:conjugateEquiv_pullbackComp_inv_mathlib}, and the right side is sent to the same hom by
    the Mathlib identity Lemma~\ref{lem:conjugateEquiv_leftAdjointCompIso_inv_mathlib}. Since conjugation
    is injective, the two inverses are equal.
  \end{proof}
\end{lemma}

\begin{lemma}\leanok
[(conj-0$'$) Iso-level form: $\mathrm{pullbackComp}$ equals $\operatorname{leftAdjointCompIso}$]
  \label{lem:pullbackComp_eq_leftAdjointCompIso}
  \lean{AlgebraicGeometry.pullbackComp_eq_leftAdjointCompIso}
  \uses{lem:pullbackComp_inv_eq_leftAdjointCompIso_inv}
  In the situation of Lemma~\ref{lem:pullbackComp_inv_eq_leftAdjointCompIso_inv}, the inverse-level
  identity upgrades to a full isomorphism equality
  \[
    \mathrm{pullbackComp}_{f,g}
      \;=\; \operatorname{leftAdjointCompIso}(\mathrm{pushforwardComp}_{f,g}),
  \]
  so the conjugate re-encoding may rewrite \(\mathrm{pullbackComp}\) for
  \(\operatorname{leftAdjointCompIso}\) (and back) freely in the comparison object.
  \begin{proof}
    \uses{lem:pullbackComp_inv_eq_leftAdjointCompIso_inv}
    An isomorphism in a category is determined by its inverse: if two isomorphisms have equal inverses
    they are equal, and an equality of isomorphisms is detected on the underlying morphisms. Applying
    these to the inverse-level identity of
    Lemma~\ref{lem:pullbackComp_inv_eq_leftAdjointCompIso_inv} gives the stated equation.
  \end{proof}
\end{lemma}

\begin{lemma}

\leanok
[Pseudofunctoriality of the pullback--pushforward unit]
  \label{lem:pullbackPushforward_unit_comp}
  \lean{AlgebraicGeometry.pullbackPushforward_unit_comp}
  \uses{lem:unit_conjugateEquiv_mathlib, lem:comp_unit_app_mathlib,
        lem:conjugateEquiv_pullbackComp_inv_mathlib}
  Let \(a : X_1 \to X_2\) and \(b : X_2 \to X_3\) be composable morphisms of schemes and let
  \(N\) be an \(\mathcal{O}_{X_3}\)-module. Write \(\eta^{h}\) for the unit of the
  pullback--pushforward adjunction \((h^{*}\dashv h_{*})\) attached to a morphism \(h\), and
  \(\mathrm{pushforwardComp}_{a,b}\), \(\mathrm{pullbackComp}_{a,b}\) for the pseudofunctor
  composition coherences. Then
  \[
    \eta^{b}_{N} \;\circ\; b_{*}\bigl(\eta^{a}_{\,b^{*}N}\bigr)
      \;\circ\; (\mathrm{pushforwardComp}_{a,b})^{\mathrm{hom}}_{N}
    \;=\;
    \eta^{a\circ b}_{N} \;\circ\;
      (a\circ b)_{*}\bigl((\mathrm{pullbackComp}_{a,b})^{\mathrm{inv}}_{N}\bigr),
  \]
  where composites are read left to right. In words: the unit of the composite morphism
  \(a \circ b\), reindexed by the inverse pullback coherence, factors as the \(b\)-unit, followed
  by \(b_{*}\) of the \(a\)-unit, followed by the pushforward coherence. This is the leg-reindex
  engine consumed by Seam~2 (Lemma~\ref{lem:base_change_mate_fstar_reindex}): with \(a\) an
  isomorphism and \(b\) an affine \(\operatorname{Spec}\)-map it turns the generic projection-leg
  unit into the affine \(\operatorname{Spec}\)-unit modulo the transparent coherences.
\end{lemma}
\begin{proof}
  \uses{lem:unit_conjugateEquiv_mathlib, lem:comp_unit_app_mathlib,
        lem:conjugateEquiv_pullbackComp_inv_mathlib}
  Apply the unit-across-conjugate coherence Lemma~\ref{lem:unit_conjugateEquiv_mathlib} to the
  composite adjunction \((b^{*}\dashv b_{*}) . (a^{*}\dashv a_{*})\), the adjunction
  \(\bigl((a\circ b)^{*}\dashv (a\circ b)_{*}\bigr)\), and the left-adjoint transformation
  \((\mathrm{pullbackComp}_{a,b})^{\mathrm{inv}}\). This yields an identity whose right-hand factor
  is the conjugate (mate) of \((\mathrm{pullbackComp}_{a,b})^{\mathrm{inv}}\); by
  Lemma~\ref{lem:conjugateEquiv_pullbackComp_inv_mathlib} that mate is exactly the pushforward
  coherence \((\mathrm{pushforwardComp}_{a,b})^{\mathrm{hom}}\). Expanding the unit of the composite
  adjunction by Lemma~\ref{lem:comp_unit_app_mathlib} —
  \(\bigl((b^{*}\dashv b_{*}).(a^{*}\dashv a_{*})\bigr)\text{-unit}_N
    = \eta^{b}_N \circ b_{*}(\eta^{a}_{b^{*}N})\) — and reassociating the resulting composite gives
  the stated equation.
\end{proof}

% --- Seam 2, step (ii): Gamma-collapse of the transparent pushforward coherences ---
% These three project-bespoke atoms record that, on global sections over Spec R, the two
% pushforward-composition coherences and the pushforward-congruence coherence appearing in
% the inner composite are transparent (their section value is the identity / a presheaf
% transport), hence collapse under the global-sections functor.

\begin{lemma}

\leanok
[$\Gamma$-collapse of the pushforward-composition coherence (hom factor)]
  \label{lem:gammaMap_pushforwardComp_hom_eq_id}
  \lean{AlgebraicGeometry.gammaMap_pushforwardComp_hom_eq_id}
  \uses{lem:gammaPushforwardIso}
  Let \(a : X_1 \to X_2\) and \(b : X_2 \to \operatorname{Spec} R\) be composable morphisms of
  schemes and \(M\) an \(\mathcal{O}_{X_1}\)-module. Applying the global-sections functor
  \(\Gamma\) (over the base \(R\)) to the component
  \((\mathrm{pushforwardComp}_{a,b})^{\mathrm{hom}}_M\) of the pushforward-composition coherence
  isomorphism yields the identity morphism:
  \[
    \Gamma\bigl((\mathrm{pushforwardComp}_{a,b})^{\mathrm{hom}}_M\bigr) = \operatorname{id}.
  \]
  The pushforward-composition coherence is transparent at the level of sections.
\end{lemma}
\begin{proof}
  \uses{lem:gammaPushforwardIso}
  The component \((\mathrm{pushforwardComp}_{a,b})^{\mathrm{hom}}_M\) evaluates to the identity
  on every open, and the global-sections functor sends the identity to the identity.
\end{proof}

\begin{lemma}

\leanok
[$\Gamma$-collapse of the pushforward-composition coherence (inv factor)]
  \label{lem:gammaMap_pushforwardComp_inv_eq_id}
  \lean{AlgebraicGeometry.gammaMap_pushforwardComp_inv_eq_id}
  \uses{lem:gammaPushforwardIso}
  In the situation of Lemma~\ref{lem:gammaMap_pushforwardComp_hom_eq_id}, the same
  \(\Gamma\)-collapse holds for the inverse component:
  \[
    \Gamma\bigl((\mathrm{pushforwardComp}_{a,b})^{\mathrm{inv}}_M\bigr) = \operatorname{id}.
  \]
\end{lemma}
\begin{proof}
  \uses{lem:gammaPushforwardIso}
  The inverse component evaluates to the identity on every open, and the global-sections
  functor preserves the identity.
\end{proof}

\begin{lemma}

\leanok
[$\Gamma$-collapse of the pushforward-congruence coherence]
  \label{lem:gammaMap_pushforwardCongr_hom}
  \lean{AlgebraicGeometry.gammaMap_pushforwardCongr_hom}
  \uses{lem:gammaPushforwardIso}
  Let \(f, g : X \to \operatorname{Spec} R\) be morphisms of schemes with \(f = g\), and \(M\)
  an \(\mathcal{O}_X\)-module. Applying the global-sections functor \(\Gamma\) to the component
  \((\mathrm{pushforwardCongr}\,(f = g))^{\mathrm{hom}}_M\) of the congruence coherence yields the
  canonical \(\operatorname{eqToHom}\) transport along the object equality induced by \(f = g\):
  \[
    \Gamma\bigl((\mathrm{pushforwardCongr}\,(f = g))^{\mathrm{hom}}_M\bigr)
      = \operatorname{eqToHom}(h),
  \]
  where \(h\) is the equality of the two section modules forced by \(f = g\). This is the precise
  form: the coherence is not literally the identity but the canonical transport across the equality
  of domains, which is the identity only after \(f\) and \(g\) have been made definitionally equal.
\end{lemma}
\begin{proof}
  \uses{lem:gammaPushforwardIso}
  Substituting \(f = g\) makes the two pushforwards along \(f\) and \(g\) the same object, and the
  congruence coherence becomes the identity; the global-sections functor sends the identity to the
  identity, which the surrounding statement records as the canonical
  \(\operatorname{eqToHom}\) transport along the induced equality of section modules. It carries no
  substantive content beyond this transport.
\end{proof}

% --- Seam 2, step (iii): the mate-unwinding crux, cut into five atomic links. ---
% These five lemmas decompose the step-(iii) obligation of
% lem:base_change_mate_fstar_reindex_legs (the mate-unwinding crux) into one
% mathematical move each. They are adjoint-mate calculus over the proved
% change-of-rings dictionaries; there is no external source. Composites
% are read left to right, as elsewhere in this chapter.

\begin{lemma}

\leanok
[(iii-1) Unit expansion: inverting the comp-coherence]
  \label{lem:base_change_mate_fstar_reindex_legs_unitExpand}
  % NOTE: superseded by the conjugate-side re-encoding (Open Q2). This direct-on-sections wrapper no
  % longer feeds the live `_legs` proof (now a conjugate-component identity). Lean pin retained.
  \lean{AlgebraicGeometry.base_change_mate_fstar_reindex_legs_unitExpand}
  \uses{lem:pullbackPushforward_unit_comp}
  Let \(e : X' \xrightarrow{\sim} \operatorname{Spec}(A \otimes_R R')\) be the canonical
  isomorphism, \(\iota_A : A \to A \otimes_R R'\) the left tensor inclusion, and set
  \(g' = e \circ \operatorname{Spec}\iota_A\). For an \(A\)-module \(M\), the bare \((g')\)-unit
  \(\eta^{g'}_{\widetilde M}\) of the \((g')\)-adjunction expands as the four-factor composite
  \[
    \eta^{g'}_{\widetilde M}
      \;=\;
      \eta^{\operatorname{Spec}\iota_A}_{\widetilde M}
      \,\circ\, (\operatorname{Spec}\iota_A)_{*}\bigl(\eta^{e}_{(\operatorname{Spec}\iota_A)^{*}\widetilde M}\bigr)
      \,\circ\, (\mathrm{pushforwardComp}_{e,\operatorname{Spec}\iota_A})^{\mathrm{hom}}_{\widetilde M}
      \,\circ\, (g')_{*}\bigl((\mathrm{pullbackComp}_{e,\operatorname{Spec}\iota_A})^{\mathrm{hom}}_{\widetilde M}\bigr),
  \]
  where composites are read left to right.
\end{lemma}
\begin{proof}
  \uses{lem:pullbackPushforward_unit_comp}
  The leg-reindex engine Lemma~\ref{lem:pullbackPushforward_unit_comp}, instantiated at \(a = e\),
  \(b = \operatorname{Spec}\iota_A\), \(N = \widetilde M\), reads
  \[
    \eta^{\operatorname{Spec}\iota_A}_{\widetilde M}
      \circ (\operatorname{Spec}\iota_A)_{*}(\eta^{e})
      \circ (\mathrm{pushforwardComp}_{e,\operatorname{Spec}\iota_A})^{\mathrm{hom}}
    \;=\;
    \eta^{g'}_{\widetilde M}
      \circ (g')_{*}\bigl((\mathrm{pullbackComp}_{e,\operatorname{Spec}\iota_A})^{\mathrm{inv}}\bigr).
  \]
  Post-compose both sides with \((g')_{*}\bigl((\mathrm{pullbackComp}_{e,\operatorname{Spec}\iota_A})^{\mathrm{hom}}\bigr)\).
  Since \((g')_{*}\) is a functor it carries
  \((\mathrm{pullbackComp})^{\mathrm{inv}} \circ (\mathrm{pullbackComp})^{\mathrm{hom}} = \mathrm{id}\)
  to the identity, so the right-hand side collapses to \(\eta^{g'}_{\widetilde M}\), giving the stated
  expansion.
\end{proof}

\begin{lemma}

\leanok
[(iii-2) Distribute the expansion through $(\operatorname{Spec}\varphi)_*$ and $\Gamma$]
  \label{lem:base_change_mate_fstar_reindex_legs_gammaDistribute}
  % NOTE: superseded by the conjugate-side re-encoding (Open Q2). This direct-on-sections wrapper no
  % longer feeds the live `_legs` proof (now a conjugate-component identity). Lean pin retained.
  \lean{AlgebraicGeometry.base_change_mate_fstar_reindex_legs_gammaDistribute}
  \uses{lem:base_change_mate_fstar_reindex_legs_unitExpand}
  In the situation of Lemma~\ref{lem:base_change_mate_fstar_reindex_legs_unitExpand}, applying the
  functor \((\operatorname{Spec}\varphi)_{*}\) and then global sections \(\Gamma = \Gamma_R\) to the
  expanded unit distributes over the four factors:
  \[
    \Gamma\bigl((\operatorname{Spec}\varphi)_{*}(\eta^{g'}_{\widetilde M})\bigr)
    \;=\;
    \Gamma\bigl((\operatorname{Spec}\varphi)_{*}\eta^{\operatorname{Spec}\iota_A}_{\widetilde M}\bigr)
    \,\circ\,
    \Gamma\bigl((\operatorname{Spec}\varphi)_{*}(\operatorname{Spec}\iota_A)_{*}\eta^{e}\bigr)
    \,\circ\,
    \Gamma\bigl((\operatorname{Spec}\varphi)_{*}(\mathrm{pushforwardComp}_{e,\operatorname{Spec}\iota_A})^{\mathrm{hom}}\bigr)
    \,\circ\,
    \Gamma\bigl((\operatorname{Spec}\varphi)_{*}(g')_{*}(\mathrm{pullbackComp}_{e,\operatorname{Spec}\iota_A})^{\mathrm{hom}}\bigr).
  \]
\end{lemma}
\begin{proof}
  \uses{lem:base_change_mate_fstar_reindex_legs_unitExpand}
  Rewrite \(\eta^{g'}_{\widetilde M}\) by the expansion of
  Lemma~\ref{lem:base_change_mate_fstar_reindex_legs_unitExpand}. Both \((\operatorname{Spec}\varphi)_{*}\)
  and \(\Gamma\) are functors, so each preserves the four-fold composite
  (\(F(u \circ v) = F u \circ F v\)). No further content.
\end{proof}

% --- Seam 2, step (iii): the mate-unwinding crux as a chain of five clean-term links. ---
% NOTE: These five lemmas re-express the step-(iii) telescoping of
% lem:base_change_mate_fstar_reindex_legs as a chain of equalities, each between two
% freshly stated terms over the same functor instance. This eliminates a nested-image vs
% composed-functor object diamond (\(G(H(M))\) vs \((H\circ G)(M)\)) that arises when
% distributing \(\Gamma\) over the composite. Each link is stated on its own clean composite
% where only one instance is in scope; the target chains the links and crosses the diamond once
% at the closing step. Throughout, \(e\) denotes the canonical square isomorphism,
% \(\iota_A = \mathrm{includeLeft}\), \(\Theta_{\mathrm{tgt}}\) the leg-parametrised codomain
% read (base_change_mate_codomain_read_legs), and \(\rho = \mathrm{base\_change\_mate\_inner\_value}\).

\begin{lemma}

\leanok
[(iii-L1$+$L2) Distribute the expanded unit and collapse the transparent $\mathrm{pushforwardComp}$ factor]
  \label{lem:base_change_mate_fstar_reindex_legs_link_distributeCollapse}
  % NOTE: superseded by the conjugate-side re-encoding (Open Q2). This direct-on-sections link no
  % longer feeds the live `_legs` proof (now a conjugate-component identity). Lean pin retained.
  \lean{AlgebraicGeometry.base_change_mate_fstar_reindex_legs_link_distributeCollapse}
  % NOTE: sole surviving fragment of the abandoned direct-on-sections route (its siblings
  % link_cancelEUnit / link_cancelPullbackComp / link_survivor were removed); retained because the
  % current base_change_mate_fstar_reindex_legs body still references this Lean declaration.
  \uses{lem:base_change_mate_fstar_reindex_legs_unitExpand,
        lem:base_change_mate_fstar_reindex_legs_gammaDistribute,
        lem:base_change_mate_inner_eCancel_pushforwardComp,
        lem:gammaMap_pushforwardComp_hom_eq_id}
  Work at the \emph{single} composite functor
  \(F := (\operatorname{Spec}\varphi)_{*} \circ \Gamma_R\) (pushforward along
  \(\operatorname{Spec}\varphi\), then global sections over \(\operatorname{Spec} R\)). Let \(a, b\) be
  the two legs of the factorization \(g' = b \circ a\) of the inner pullback leg (in the application
  \(a = e\) is the square isomorphism and \(b = \operatorname{Spec}\iota_A\)), and \(N\) an
  \(\mathcal{O}_{\operatorname{Spec} A}\)-module (in the application \(N = \widetilde M\)). Read at the
  one functor \(F\), the \(F\)-image of the \((b \circ a)\)-pullback--pushforward unit at \(N\) equals
  the three-factor composite
  \[
    F\bigl(\eta^{b}_{N}\bigr)
    \,\circ\,
    F\bigl(b_{*}\,\eta^{a}_{(\mathrm{pull}\,b)N}\bigr)
    \,\circ\,
    F\bigl((b \circ a)_{*}\,(\mathrm{pullbackComp}_{a,b})^{\mathrm{hom}}_{N}\bigr).
  \]
  The point is that both sides are read through the same functor instance \(F\): only one
  \(\mathcal{O}_X\)-module category structure is in play, so functoriality applies with no instance
  diamond. The distribution of \(F\) over the four expanded-unit factors and the collapse of the
  transparent middle \(\mathrm{pushforwardComp}\) factor are thereby fused into one diamond-free
  identity, the clean building block spliced into the telescoping of
  Lemma~\ref{lem:base_change_mate_fstar_reindex_legs}.
\end{lemma}
\begin{proof}
  \uses{lem:base_change_mate_fstar_reindex_legs_unitExpand,
        lem:base_change_mate_fstar_reindex_legs_gammaDistribute,
        lem:base_change_mate_inner_eCancel_pushforwardComp,
        lem:gammaMap_pushforwardComp_hom_eq_id}
  Expand the bare \((b \circ a)\)-unit by
  Lemma~\ref{lem:base_change_mate_fstar_reindex_legs_unitExpand} into its four factors, then distribute
  \(F\) over the composite by
  Lemma~\ref{lem:base_change_mate_fstar_reindex_legs_gammaDistribute}
  (\(F(u \circ v) = F u \circ F v\)). Because both sides are expressed through the one functor \(F\),
  this is a direct application of functoriality with no re-elaboration of the underlying category
  structure. Of the four resulting \(F\)-image factors, the third is
  \(F\bigl((\mathrm{pushforwardComp}_{a,b})^{\mathrm{hom}}\bigr)\): its section value is the identity
  (Lemma~\ref{lem:gammaMap_pushforwardComp_hom_eq_id}, the content of
  Lemma~\ref{lem:base_change_mate_inner_eCancel_pushforwardComp}), and \(F\) carries the identity to
  the identity. Collapsing this factor by the unit law and reassociating yields the stated three-factor
  form.
\end{proof}

% --- Seam 2, step (iii) FINE re-break: the conjugate-side discharge of the
%     leg-reindex crux, atomised into one mathematical claim per lemma.
%     This chain cuts at the real conjugate-calculus seams. All blocks are adjoint-mate
%     calculus over the proved change-of-rings dictionaries and the conj-0 foundation
%     (lem:pullbackComp_eq_leftAdjointCompIso); there is no external source. Composites are
%     read left to right.

\begin{lemma}[Conjugation distributes over composition (reassoc simp lemma)]
  \label{lem:conjugateEquiv_comp_mathlib}
  \lean{CategoryTheory.conjugateEquiv_comp}
  \mathlibok
  \textit{Provided by Mathlib.}
  Let \(\mathrm{adj}_1 : L_1 \dashv R_1\), \(\mathrm{adj}_2 : L_2 \dashv R_2\),
  \(\mathrm{adj}_3 : L_3 \dashv R_3\) be adjunctions between the same pair of categories, and let
  \(\alpha : L_2 \Rightarrow L_1\), \(\beta : L_3 \Rightarrow L_2\) be transformations of left adjoints.
  Then the conjugates compose, with order reversed onto the right-adjoint side:
  \[
    \operatorname{conj}(\mathrm{adj}_1,\mathrm{adj}_2)\,\alpha \;\circ\;
    \operatorname{conj}(\mathrm{adj}_2,\mathrm{adj}_3)\,\beta
      \;=\; \operatorname{conj}(\mathrm{adj}_1,\mathrm{adj}_3)\,(\beta \circ \alpha) .
  \]
  Tagged \(\mathtt{@[reassoc (attr := simp)]}\), it is the lead member of the conjugate
  reassociation simp set: it fuses the per-adjunction conjugates of a composite adjunction into the
  conjugate of the composite (the conjugate-side form of the paste-of-mates).
\end{lemma}

\begin{lemma}[Inverse conjugation distributes over composition (reassoc simp lemma)]
  \label{lem:conjugateEquiv_symm_comp_mathlib}
  \lean{CategoryTheory.conjugateEquiv_symm_comp}
  \mathlibok
  \textit{Provided by Mathlib.}
  The inverse-direction companion of Lemma~\ref{lem:conjugateEquiv_comp_mathlib}: for transformations
  \(\alpha : R_1 \Rightarrow R_2\), \(\beta : R_2 \Rightarrow R_3\) of the right adjoints,
  \[
    \operatorname{conj}(\mathrm{adj}_2,\mathrm{adj}_3)^{-1}\beta \;\circ\;
    \operatorname{conj}(\mathrm{adj}_1,\mathrm{adj}_2)^{-1}\alpha
      \;=\; \operatorname{conj}(\mathrm{adj}_1,\mathrm{adj}_3)^{-1}(\alpha \circ \beta) ,
  \]
  also \(\mathtt{@[reassoc (attr := simp)]}\). Together these two close the reassociation bookkeeping of
  the conjugate-side discharge.
\end{lemma}

\begin{definition}

\leanok
[The codomain read with an abstracted pullback factor]
  \label{def:base_change_mate_codomain_read_legs_param}
  \lean{AlgebraicGeometry.base_change_mate_codomain_read_legs_param}
  \uses{lem:base_change_mate_codomain_read_legs, lem:pullback_spec_tilde_iso,
        lem:pushforward_spec_tilde_iso}
  The variable-legs codomain read Lemma~\ref{lem:base_change_mate_codomain_read_legs}, restated with
  its single pullback-composition factor \(\mathrm{pullbackComp}_{e,\operatorname{Spec}\iota_A}\)
  abstracted as an \emph{explicit isomorphism argument} \(P_{\mathrm{comp}}\). For any iso
  \(P_{\mathrm{comp}}\) of the appropriate type it assembles the comparison
  \(\Gamma\bigl(f'_*(g')^*\widetilde M\bigr) \xrightarrow{\sim} (R' \otimes_R A) \otimes_A M\) by the
  same recipe as Lemma~\ref{lem:base_change_mate_codomain_read_legs}, using \(P_{\mathrm{comp}}\) in
  place of the project coherence. Instantiating \(P_{\mathrm{comp}}\) at the project's
  \(\mathrm{pullbackComp}_{e,\operatorname{Spec}\iota_A}\) recovers the original read; instantiating it
  at the abstract \(\operatorname{leftAdjointCompIso}\) of the free legs
  (Definition~\ref{def:conjPullbackFactor}) gives the conjugate-native read
  Lemma~\ref{lem:base_change_mate_codomain_read_legs_conj} (conj-1a).
\end{definition}

\begin{lemma}

\leanok
[The original read is the abstracted read at the project coherence]
  \label{lem:base_change_mate_codomain_read_legs_eq_param}
  \lean{AlgebraicGeometry.base_change_mate_codomain_read_legs_eq_param}
  \uses{def:base_change_mate_codomain_read_legs_param, lem:base_change_mate_codomain_read_legs}
  The variable-legs codomain read Lemma~\ref{lem:base_change_mate_codomain_read_legs} equals the
  abstracted read Definition~\ref{def:base_change_mate_codomain_read_legs_param} instantiated at the
  project pullback-composition coherence \(\mathrm{pullbackComp}_{e,\operatorname{Spec}\iota_A}\).
\end{lemma}
\begin{proof}
  \uses{def:base_change_mate_codomain_read_legs_param}
  The parametrized body is, by construction, the original body with the pullback factor named; the two
  are therefore definitionally equal, so the equality holds by reflexivity.
\end{proof}

\begin{definition}

\leanok
[The free-leg pullback factor]
  \label{def:conjPullbackFactor}
  \lean{AlgebraicGeometry.conjPullbackFactor}
  \uses{lem:leftAdjointCompIso_mathlib}
  The abstract composite-adjunction comparison \(\operatorname{leftAdjointCompIso}\)
  (Lemma~\ref{lem:leftAdjointCompIso_mathlib}) of the two free legs
  \(e = \mathrm{pullbackSpecIso}^{\mathrm{hom}}\) and \(\operatorname{Spec}\iota_A\), assembled from the
  pushforward composition coherence \(\mathrm{pushforwardComp}_{e,\operatorname{Spec}\iota_A}\). It is
  the pullback factor used by the conjugate-native codomain read
  Lemma~\ref{lem:base_change_mate_codomain_read_legs_conj}, isolated here as a named abbreviation.
\end{definition}

\begin{lemma}

\leanok
[The free-leg pullback factor is the project coherence]
  \label{lem:conjPullbackFactor_eq_pullbackComp}
  \lean{AlgebraicGeometry.conjPullbackFactor_eq_pullbackComp}
  \uses{def:conjPullbackFactor, lem:pullbackComp_eq_leftAdjointCompIso}
  The free-leg pullback factor Definition~\ref{def:conjPullbackFactor} equals the project's
  pseudofunctor pullback-composition coherence \(\mathrm{pullbackComp}_{e,\operatorname{Spec}\iota_A}\).
  This is the iso-level identity \((\mathrm{conj}\text{-}0')\) of
  Lemma~\ref{lem:pullbackComp_eq_leftAdjointCompIso} packaged at the free legs.
\end{lemma}
\begin{proof}
  \uses{lem:pullbackComp_eq_leftAdjointCompIso}
  Both sides unfold to the abstract \(\operatorname{leftAdjointCompIso}\) of the pushforward coherence;
  the identity is the free-leg instance of Lemma~\ref{lem:pullbackComp_eq_leftAdjointCompIso}.
\end{proof}

\begin{lemma}

\leanok
[(conj-1a) Codomain read assembled natively in the composite-conjugate calculus]
  \label{lem:base_change_mate_codomain_read_legs_conj}
  \lean{AlgebraicGeometry.base_change_mate_codomain_read_legs_conj}
  \uses{lem:pullbackComp_eq_leftAdjointCompIso, lem:leftAdjointCompIso_mathlib,
        lem:conjugateEquiv_leftAdjointCompIso_inv_mathlib, lem:pullback_spec_tilde_iso,
        lem:pushforward_spec_tilde_iso, def:base_change_mate_codomain_read_legs_param,
        def:conjPullbackFactor}
  This is the codomain read of
  Lemma~\ref{lem:base_change_mate_codomain_read_legs} \emph{rebuilt proof-free, natively in the
  composite-conjugate calculus}. Fix \(\psi : R \to R'\), \(\varphi : R \to A\) and an \(A\)-module
  \(M\); let \(e : X' \xrightarrow{\sim} \operatorname{Spec}(A \otimes_R R')\) be the canonical square
  isomorphism (Lemma~\ref{lem:pullbackSpecIso_mathlib}) and \(\iota_A : A \to A \otimes_R R'\),
  \(\iota_{R'} : R' \to A \otimes_R R'\) the tensor inclusions. Reading the codomain
  \(f'_*(g')^*\widetilde M\) along the \emph{explicit composite} legs
  \(g' = e \circ \operatorname{Spec}\iota_A\), \(f' = e \circ \operatorname{Spec}\iota_{R'}\), the
  comparison isomorphism of \(R'\)-modules
  \[
    \Theta_{\mathrm{tgt}}^{\sharp} :
    \Gamma\bigl(f'_*(g')^*\widetilde M\bigr)
    \xrightarrow{\ \sim\ }
    (R' \otimes_R A) \otimes_A M
  \]
  is assembled with its pullback factor taken as the composite-adjunction comparison
  \(\operatorname{leftAdjointCompIso}\) (Lemma~\ref{lem:leftAdjointCompIso_mathlib}) of the two free
  morphisms \(e\), \(\operatorname{Spec}\iota_A\) — the pasted square entering through the pushforward
  congruence \(\mathrm{pushforwardCongr}\) and composition \(\mathrm{pushforwardComp}\) coherences —
  rather than as the project's \(\mathrm{pullbackComp}\); the two agree by
  Lemma~\ref{lem:pullbackComp_eq_leftAdjointCompIso}, so this read carries \emph{no} leg-equality data
  \(\mathit{hfst}/\mathit{hsnd}\). It is then composed with the affine pullback
  (Lemma~\ref{lem:pullback_spec_tilde_iso}) and pushforward (Lemma~\ref{lem:pushforward_spec_tilde_iso})
  dictionaries. Because the comparison object is built from the free morphisms, its conjugate inverts to
  the right-adjoint comparison in one step
  (Lemma~\ref{lem:conjugateEquiv_leftAdjointCompIso_inv_mathlib}).
\end{lemma}

\begin{proof}
  \uses{lem:pullbackComp_eq_leftAdjointCompIso, lem:leftAdjointCompIso_mathlib,
        lem:pullback_spec_tilde_iso, lem:pushforward_spec_tilde_iso}
  Copy the assembly of Lemma~\ref{lem:base_change_mate_codomain_read_legs} verbatim, with the single
  change that the pullback-composition factor \(\mathrm{pullbackComp}_{e,\operatorname{Spec}\iota_A}\) is
  replaced by the abstract composite-adjunction comparison
  \(\operatorname{leftAdjointCompIso}(\mathrm{pushforwardComp}_{e,\operatorname{Spec}\iota_A})\). The
  replacement is along the identity of isomorphisms
  Lemma~\ref{lem:pullbackComp_eq_leftAdjointCompIso}, so the resulting object isomorphism is the same
  data presented natively in the conjugate calculus and carries no propositional leg-equality.
\end{proof}

\begin{lemma}

\leanok
[(conj-1b) The conjugate-native read agrees with the concrete read]
  \label{lem:base_change_mate_codomain_read_legs_conj_eq}
  \lean{AlgebraicGeometry.base_change_mate_codomain_read_legs_conj_eq}
  \uses{lem:base_change_mate_codomain_read_legs_conj, lem:base_change_mate_codomain_read_legs,
        lem:base_change_mate_codomain_read}
  The conjugate-native codomain read \(\Theta_{\mathrm{tgt}}^{\sharp}\)
  (Lemma~\ref{lem:base_change_mate_codomain_read_legs_conj}) equals, as an isomorphism of
  \(R'\)-modules, the leg-parametrised read \(\Theta_{\mathrm{tgt}}^{\flat}\)
  (Lemma~\ref{lem:base_change_mate_codomain_read_legs}) and hence, at the literal projection legs, the
  concrete read \(\Theta_{\mathrm{tgt}}\) (Lemma~\ref{lem:base_change_mate_codomain_read}). This single
  bridge keeps the downstream consumers of the concrete read — in particular the green reduction of
  Lemma~\ref{lem:base_change_mate_fstar_reindex} — defeq-unchanged when the reindex crux is re-expressed
  against the conjugate-native object.
\end{lemma}

\begin{proof}
  \uses{lem:base_change_mate_codomain_read_legs_conj, lem:base_change_mate_codomain_read_legs}
  The only difference between the two assemblies is the pullback factor, which is
  \(\mathrm{pullbackComp}_{e,\operatorname{Spec}\iota_A}\) in the leg-parametrised read and
  \(\operatorname{leftAdjointCompIso}(\mathrm{pushforwardComp}_{e,\operatorname{Spec}\iota_A})\) in the
  conjugate-native read; these are equal by
  Lemma~\ref{lem:pullbackComp_eq_leftAdjointCompIso}. Substituting that equality identifies the two
  composites; the agreement with the concrete read at the projection legs is the cone-leg factorization
  already recorded for Lemma~\ref{lem:base_change_mate_codomain_read_legs}.
\end{proof}

\begin{lemma}

\leanok
[(conj-2b) The pullback-side leg of the conjugate identity]
  \label{lem:base_change_mate_reindex_conj_pullbackLeg}
  \lean{AlgebraicGeometry.base_change_mate_reindex_conj_pullbackLeg}
  \uses{lem:conjugateEquiv_pullbackComp_inv_mathlib,
        lem:conjugateEquiv_leftAdjointCompIso_inv_mathlib, lem:pullbackComp_eq_leftAdjointCompIso}
  In the conjugate-side discharge of the leg-reindex (after
  \(\operatorname{conj}^{-1}\) has been applied to both sides, Lemma~\ref{lem:conjugateIsoEquiv_mathlib}),
  the pullback factor is identified in one step. The conjugate of the inverse of the composite
  comparison \(\operatorname{leftAdjointCompIso}\) carried by the conjugate-native read
  (Lemma~\ref{lem:base_change_mate_codomain_read_legs_conj}) is, by
  Lemma~\ref{lem:conjugateEquiv_leftAdjointCompIso_inv_mathlib}, the right-adjoint square isomorphism
  \(e\); and the conjugate of the inverse pullback-composition coherence
  \((\mathrm{pullbackComp}_{e,\operatorname{Spec}\iota_A})^{\mathrm{inv}}\) is the pushforward
  composition coherence \((\mathrm{pushforwardComp}_{e,\operatorname{Spec}\iota_A})^{\mathrm{hom}}\) by
  Lemma~\ref{lem:conjugateEquiv_pullbackComp_inv_mathlib} (the two coherences being interchangeable on
  the comparison object by Lemma~\ref{lem:pullbackComp_eq_leftAdjointCompIso}). Thus the pullback-side
  leg of the conjugate identity is the pushforward coherence, with no propositional transport.
\end{lemma}

\begin{proof}
  \uses{lem:conjugateEquiv_pullbackComp_inv_mathlib,
        lem:conjugateEquiv_leftAdjointCompIso_inv_mathlib, lem:pullbackComp_eq_leftAdjointCompIso}
  Both identifications are direct applications of the cited conjugation identities to the free legs
  \(e\), \(\operatorname{Spec}\iota_A\); the interchange of \(\mathrm{pullbackComp}\) and
  \(\operatorname{leftAdjointCompIso}\) is Lemma~\ref{lem:pullbackComp_eq_leftAdjointCompIso}.
\end{proof}

\begin{lemma}

\leanok
[(conj-2c) The pushforward-side collapse of the conjugate identity]
  \label{lem:base_change_mate_reindex_conj_pushforwardCollapse}
  \lean{AlgebraicGeometry.base_change_mate_reindex_conj_pushforwardCollapse}
  \uses{lem:gammaMap_pushforwardComp_hom_eq_id, lem:gammaMap_pushforwardComp_inv_eq_id,
        lem:gammaMap_pushforwardCongr_hom}
  On the pushforward side of the conjugate identity, read on global sections over
  \(\operatorname{Spec} R\), the transparent coherences collapse: the two
  pushforward-composition coherences
  \((\mathrm{pushforwardComp})^{\mathrm{hom}}\), \((\mathrm{pushforwardComp})^{\mathrm{inv}}\) and the
  pushforward-congruence coherence \((\mathrm{pushforwardCongr}\,\mathit{comm})^{\mathrm{hom}}\) all have
  identity \(\Gamma\)-image (Lemmas~\ref{lem:gammaMap_pushforwardComp_hom_eq_id},
  \ref{lem:gammaMap_pushforwardComp_inv_eq_id}, \ref{lem:gammaMap_pushforwardCongr_hom}). Hence these
  three factors of the inner composite drop out, leaving only the unit factor to be evaluated.
\end{lemma}

\begin{proof}
  \uses{lem:gammaMap_pushforwardComp_hom_eq_id, lem:gammaMap_pushforwardComp_inv_eq_id,
        lem:gammaMap_pushforwardCongr_hom}
  Each of the three coherence components is transparent on sections — its value at every open is the
  identity or a presheaf transport — so its image under \(\Gamma = \Gamma_R\) is the identity by the
  three collapse lemmas cited. A composite is unchanged by deleting identity factors.
\end{proof}

\begin{lemma}

\leanok
[(conj-2d) Cross-layer transport of the affine unit through the pushforward $\Gamma$-comparison]
  \label{lem:base_change_mate_reindex_conj_crossLayer}
  \lean{AlgebraicGeometry.base_change_mate_reindex_conj_crossLayer}
  \uses{lem:base_change_mate_unit_value, lem:unit_conjugateEquiv_symm_mathlib,
        lem:conjugateEquiv_comp_mathlib, lem:gammaPushforwardIso}
  This is the historical bottom of the term-mode route, isolated here as a single conjugate-side claim.
  The surviving \((\operatorname{Spec}\iota_A)\)-unit factor of the inner composite must be transported
  across the codomain read's \(\mathrm{unit\_iso}^{-1}\) factor, which lives one functor layer up,
  through the pushforward \(\Gamma\)-comparison \(\gamma_\psi\) (Lemma~\ref{lem:gammaPushforwardIso}).
  The claim is that this cross-layer transport is effected by the transposed unit-across-conjugate
  coherence Lemma~\ref{lem:unit_conjugateEquiv_symm_mathlib}, composed one functor layer up via the
  conjugation-composition coherence Lemma~\ref{lem:conjugateEquiv_comp_mathlib} (the Seam~1 tool raised
  one layer), and that the lone surviving affine unit then evaluates by Seam~1
  (Lemma~\ref{lem:base_change_mate_unit_value}) to the inner value
  \(\rho : m \mapsto (1 \otimes 1) \otimes m\). The point is that the naturality of \(\gamma_\psi\)
  enters \emph{as a conjugate-component identity}, never as a positional naturality rewrite under the
  \(\mathcal{O}_X\)-module instance diamond.
\end{lemma}

\begin{proof}
  \uses{lem:base_change_mate_unit_value, lem:unit_conjugateEquiv_symm_mathlib,
        lem:conjugateEquiv_comp_mathlib}
  Apply the transposed unit-across-conjugate coherence
  Lemma~\ref{lem:unit_conjugateEquiv_symm_mathlib} to the right-adjoint transformation
  \(\gamma_\psi\)-naturality factor, composed through the conjugation-composition coherence
  Lemma~\ref{lem:conjugateEquiv_comp_mathlib} to raise it past the
  \((\operatorname{Spec}\varphi)_*\) layer (exactly the Seam~1 mechanism, one functor layer up). This
  cancels the codomain read's \(\mathrm{unit\_iso}^{-1}\), leaving the bare affine
  \((\operatorname{Spec}\iota_A)\)-unit, whose value through the tilde dictionaries is the algebraic unit
  \(\rho\) of Lemma~\ref{lem:base_change_mate_unit_value}.
\end{proof}

\begin{definition}

\leanok
[Keystone right adjunction of the conjugate pair]
  \label{def:keystone_adjR}
  \lean{AlgebraicGeometry.keystoneAdjR}
  \uses{lem:pullbackSpecIso_mathlib}
  The right member \(\mathrm{adjR}\) of the conjugate pair used to prove the keystone leg-reindex
  coherence (Lemma~\ref{lem:base_change_mate_fstar_reindex_legs_conj}). With
  \(\iota_A = \mathrm{includeLeftRingHom} : R \to A \otimes_R R'\) and \(e\) the canonical isomorphism
  \(\mathrm{pullbackSpecIso}\), it is the nested composite adjunction
  \[
    \mathrm{adjR} \;=\;
    (\mathrm{extendRestrictScalarsAdj}\,\iota_A)\,\circ\,
    \bigl((\widetilde{\;\cdot\;}\dashv \Gamma)_{A\otimes_R R'}\,\circ\,
    (\mathrm{pullbackPushforwardAdj}\,e)\bigr),
  \]
  an adjunction \(\mathrm{ModuleCat}\,A \rightleftarrows \mathrm{pullback}.\mathrm{Modules}\) whose left
  adjoint is \(\mathrm{extendScalars}\,\iota_A \;\Rrightarrow\; \widetilde{(\cdot)}_{A\otimes R'}
  \;\Rrightarrow\; \mathrm{pullback}\,e\) and whose right adjoint is
  \((\mathrm{pushforward}\,e \Rrightarrow \Gamma_{A\otimes R'}) \Rrightarrow \mathrm{restrictScalars}\,\iota_A\).
  It is conjugate-comparable with the left adjunction \(\mathrm{adjL}\) (same source
  \(\mathrm{ModuleCat}\,A\) and target \(\mathrm{pullback}.\mathrm{Modules}\)).
\end{definition}

\begin{definition}

\leanok
[Keystone comparison natural isomorphism]
  \label{def:keystone_beta}
  \lean{AlgebraicGeometry.keystoneBeta}
  \uses{def:keystone_adjR, lem:gammaPushforwardNatIso}
  The comparison natural isomorphism \(\beta : R_1 \cong R_2\) between the two right adjoints of the keystone
  conjugate pair, i.e.\ \(\mathrm{pushforward}\,g' \Rrightarrow \Gamma_A \;\cong\;
  (\mathrm{pushforward}\,e \Rrightarrow \Gamma_{A\otimes R'}) \Rrightarrow \mathrm{restrictScalars}\,\iota_A\),
  assembled layer by layer: the pushforward-composition collapse (via
  \(\mathrm{pushforwardComp}\)) whiskered on \(\Gamma_A\), reassociated, then whiskered against
  \(\mathrm{gammaPushforwardNatIso}\,\iota_A\)
  (Lemma~\ref{lem:gammaPushforwardNatIso}) and reassociated back. Together with \(\mathrm{adjL}\) and
  \(\mathrm{adjR}\) (Definition~\ref{def:keystone_adjR}), it witnesses the \(A\)-level conjugate-unit
  natural transformation \(\mathrm{unit\_conjugateEquiv\_symm}\,\mathrm{adjL}\,\mathrm{adjR}\,\beta\),
  which is the principal input to Lemma~\ref{lem:base_change_mate_fstar_reindex_legs_conj}.
\end{definition}

\begin{lemma}

\leanok
[(conj-2a) The leg-reindex coherence on the conjugate side]
  \label{lem:base_change_mate_fstar_reindex_legs_conj}
  \lean{AlgebraicGeometry.base_change_mate_fstar_reindex_legs_conj}
  % NOTE: all three supporting legs are axiom-clean (conj-2b
  % base_change_mate_reindex_conj_pullbackLeg, conj-2c ..._pushforwardCollapse, conj-2d
  % ..._crossLayer). Route: ALIGN_WITH_MATHLIB → adopt the leftAdjointCompIso /
  % conjugate-.injective discharge (Fallback B), DROP element-ext (Fallback A, reserved by Mathlib
  % for atomic dictionary leaves only). The fix is NOT a one-shot recognise but the Mathlib idiom of
  % peeling ONE adjunction-pair / functor-layer at a time via conjugateEquiv_symm_comp + whiskering,
  % exactly as Mathlib's leftAdjointCompNatTrans₀₂₃_eq_conjugateEquiv_symm
  % (CompositionIso.lean:140) does. The proof sketch below encodes this route.
  \uses{lem:base_change_mate_codomain_read_legs_conj, lem:conjugateIsoEquiv_mathlib,
        lem:iterated_mateEquiv_conjugateEquiv_mathlib, lem:conjugateEquiv_comp_mathlib,
        lem:conjugateEquiv_symm_comp_mathlib, lem:base_change_mate_reindex_conj_pullbackLeg,
        lem:base_change_mate_reindex_conj_pushforwardCollapse,
        lem:base_change_mate_reindex_conj_crossLayer, def:base_change_mate_inner_value,
        def:keystone_adjR, def:keystone_beta}
  This is the leg-reindex coherence of Lemma~\ref{lem:base_change_mate_fstar_reindex_legs}, restated at
  the explicit composite legs \(g' = e \circ \operatorname{Spec}\iota_A\),
  \(f' = e \circ \operatorname{Spec}\iota_{R'}\) and \emph{against the conjugate-native codomain read}
  \(\Theta_{\mathrm{tgt}}^{\sharp}\) (Lemma~\ref{lem:base_change_mate_codomain_read_legs_conj}): the
  inner composite \(\theta_{\mathrm{in}}\), read on global sections over \(\operatorname{Spec} R\) with
  that codomain, equals the canonical \(R\)-linear map
  \(\rho : m \mapsto (1 \otimes 1) \otimes m\) of Definition~\ref{def:base_change_mate_inner_value}.
  Both sides are values of the conjugation bijection (Lemma~\ref{lem:conjugateIsoEquiv_mathlib};
  abstractly, the Beck--Chevalley reformulation Lemma~\ref{lem:iterated_mateEquiv_conjugateEquiv_mathlib}
  certifies that the comparison is a conjugate), so it suffices to prove the equality of their conjugate
  preimages on the right-adjoint side.
\end{lemma}

\begin{proof}
  \uses{lem:base_change_mate_codomain_read_legs_conj, lem:conjugateIsoEquiv_mathlib,
        lem:conjugateEquiv_comp_mathlib, lem:conjugateEquiv_symm_comp_mathlib,
        lem:base_change_mate_reindex_conj_pullbackLeg,
        lem:base_change_mate_reindex_conj_pushforwardCollapse,
        lem:base_change_mate_reindex_conj_crossLayer}
  Apply \((\operatorname{conj})^{-1}\) to both sides (injectivity of the conjugation bijection,
  Lemma~\ref{lem:conjugateIsoEquiv_mathlib}), and lift every locked right-adjoint-side natural
  transformation to a free left-adjoint-side preimage by surjectivity of the same bijection, so that no
  diamond-locked literal remains. Crucially, do \emph{not} attempt to recognise the whole inner
  composite as a single conjugate value in one step (the one-shot approach); instead, mirror
  Mathlib's \texttt{leftAdjointCompNatTrans₀₂₃\_eq\_conjugateEquiv\_symm}
  (\texttt{CompositionIso.lean}) and \emph{peel one adjunction-pair / functor-layer at a time}: split
  the composite into single-pair conjugate factors with \(\operatorname{conjugateEquiv\_symm\_comp}\)
  (Lemma~\ref{lem:conjugateEquiv_symm_comp_mathlib}) and discharge each factor by its push-through law
  (whiskering coherences of the same family). The resulting per-layer identities are exactly the three
  isolated legs: the pullback-side identification
  Lemma~\ref{lem:base_change_mate_reindex_conj_pullbackLeg}, the pushforward-side collapse
  Lemma~\ref{lem:base_change_mate_reindex_conj_pushforwardCollapse}, and the cross-layer transport
  Lemma~\ref{lem:base_change_mate_reindex_conj_crossLayer} (the latter built by generalising the
  \(\operatorname{unit\_conjugateEquiv\_symm}\) value \texttt{huce} to absorb the remaining factors),
  with the reassociation handled by the conjugate simp set (Lemmas~\ref{lem:conjugateEquiv_comp_mathlib},
  \ref{lem:conjugateEquiv_symm_comp_mathlib}). The surviving affine unit evaluates to \(\rho\) inside the
  cross-layer leg.
\end{proof}

\begin{lemma}

\leanok
[The pushforward pseudofunctor reindex at the explicit composite leg]
  \label{lem:base_change_mate_fstar_reindex_legs}
  \lean{AlgebraicGeometry.base_change_mate_fstar_reindex_legs}
  \uses{lem:base_change_mate_fstar_reindex_legs_conj,
        lem:base_change_mate_codomain_read_legs_conj_eq,
        lem:base_change_mate_codomain_read_legs, def:base_change_mate_inner_value}
  % ATOMISED. The conjugate-side discharge is cut into the chain conj-1a..conj-2d + the assembled
  % conjugate identity conj-2a (lem:base_change_mate_fstar_reindex_legs_conj). This proof is the thin
  % wrapper: instantiate the conjugate identity at the explicit composite legs and bridge the codomain
  % read back to the concrete one by lem:base_change_mate_codomain_read_legs_conj_eq.
  % The coherence is proved as an identity of `Scheme.Modules.conjugateEquiv` COMPONENTS, never as a
  % positional rewrite under the X.Modules instance diamond.
  \emph{Why the legs are taken as the explicit composite.} The reindex coherence is read along the
  explicit composite leg \(g' = e \circ \operatorname{Spec}\iota_A\) (and
  \(f' = e \circ \operatorname{Spec}\iota_{R'}\)), with \(e\) the canonical square isomorphism
  (Lemma~\ref{lem:pullbackSpecIso_mathlib}) and \(\iota_A,\iota_{R'}\) the tensor inclusions. Taking
  the legs as this composite — rather than as the literal projections, which are only propositionally
  equal to it — is what lets the codomain read be assembled from the free morphisms
  (Lemma~\ref{lem:base_change_mate_codomain_read_legs}) and the coherence be discharged on the
  conjugate side, where the \(\mathcal{O}_X\)-module instance diamond never bites.

  Fix \(\psi : R \to R'\), \(\varphi : R \to A\) and an \(A\)-module \(M\). The inner composite
  \(\theta_{\mathrm{in}}\) built from the \((g')\)-unit and the pushforward coherences, read on global
  sections over \(\operatorname{Spec} R\) with the codomain supplied by
  Lemma~\ref{lem:base_change_mate_codomain_read_legs}, equals the canonical \(R\)-linear map
  \(\rho : m \mapsto (1 \otimes 1) \otimes m\) of Definition~\ref{def:base_change_mate_inner_value}.
  Instantiating at the literal projection legs \(g' = \operatorname{pr}_1\),
  \(f' = \operatorname{pr}_2\) recovers the concrete Lemma~\ref{lem:base_change_mate_fstar_reindex}.
\end{lemma}
\begin{proof}
  \uses{lem:base_change_mate_fstar_reindex_legs_conj,
        lem:base_change_mate_codomain_read_legs_conj_eq,
        lem:base_change_mate_codomain_read_legs, def:base_change_mate_inner_value}
  The whole conjugate-side discharge is now carried by the assembled conjugate identity
  Lemma~\ref{lem:base_change_mate_fstar_reindex_legs_conj}, which states exactly this equality but with
  the codomain supplied by the conjugate-native read \(\Theta_{\mathrm{tgt}}^{\sharp}\)
  (Lemma~\ref{lem:base_change_mate_codomain_read_legs_conj}). Bridging that read back to the
  leg-parametrised read \(\Theta_{\mathrm{tgt}}^{\flat}\)
  (Lemma~\ref{lem:base_change_mate_codomain_read_legs}) appearing in the present statement is the lone
  remaining step, supplied by Lemma~\ref{lem:base_change_mate_codomain_read_legs_conj_eq}: substituting
  the equality of the two reads turns Lemma~\ref{lem:base_change_mate_fstar_reindex_legs_conj} into the
  asserted identity, whose right-hand side is \(\rho\)
  (Definition~\ref{def:base_change_mate_inner_value}). No positional rewrite under the
  \(\mathcal{O}_X\)-module instance diamond is performed at this level; the entire conjugate-component
  argument lives in the cited sub-lemmas.
\end{proof}

% NOTE: the `(g^* ⊣ g_*)`-transpose crux lem:base_change_mate_gstar_transpose cascades off
% this `_legs` identity by the SAME conjugate-side mechanism (its `Θ_src`/regroup/`Θ_tgt` move past the
% geometric factors is again a `conjugateEquiv` component identity, closed by the reassoc conjugate simp
% set + conjugateEquiv_pullbackComp_inv). It is a SEPARATE target.

\begin{lemma}

\leanok
[The pushforward pseudofunctor reindex of the inner comparison]
  \label{lem:base_change_mate_fstar_reindex}
  \lean{AlgebraicGeometry.base_change_mate_fstar_reindex}
  \uses{lem:base_change_mate_fstar_reindex_legs, lem:base_change_mate_unit_value,
        lem:base_change_mate_codomain_read, lem:pullback_fst_snd_specMap_tensor,
        lem:pushforward_spec_tilde_iso, lem:pullbackPushforward_unit_comp,
        def:base_change_mate_inner_value}
  % SOURCE: [Stacks Project], Cohomology of Schemes, Lemma `lemma-affine-base-change'
  % (\S\,Cohomology and base change, I), proof, "We use Schemes, Lemma
  % \ref{schemes-lemma-widetilde-pullback} to describe pullbacks and pushforwards"
  % (read from references/stacks-coherent.tex, L927--932).
  % SOURCE QUOTE: "We use Schemes, Lemma \ref{schemes-lemma-widetilde-pullback}
  % to describe pullbacks and pushforwards of $\mathcal{F}$. Namely,
  % $X' = \Spec(R' \otimes_R A)$ and $\mathcal{F}'$ is the quasi-coherent sheaf
  % associated to $(R' \otimes_R A) \otimes_A M$."
  \textit{Source: Stacks Project, Cohomology of Schemes, Lemma ``Affine base change''
  (proof, ``describe pullbacks and pushforwards'' step).}
  This is the literal-leg instantiation of the leg-parametrised engine
  Lemma~\ref{lem:base_change_mate_fstar_reindex_legs}.
  In the generic pullback square \(X' = \operatorname{pullback}(\operatorname{Spec}\varphi,
  \operatorname{Spec}\psi)\), with \(g'\) the first projection and \(f'\) the second,
  write \(\theta_{\mathrm{in}}\) for the inner composite built from the \((g')\)-unit and
  the pushforward pseudofunctor coherences that appear in the base-change map
  (Definition~\ref{def:pushforward_base_change_map}) before the
  \((g^* \dashv g_*)\)-transpose.
  Read on the global sections over \(\operatorname{Spec} R\) through the \(\Gamma\)-pushforward
  dictionaries (Lemma~\ref{lem:pushforward_spec_tilde_iso},
  Lemma~\ref{lem:gammaPushforwardTildeIso}), with the codomain identified by
  Lemma~\ref{lem:base_change_mate_codomain_read}, the map \(\theta_{\mathrm{in}}\) is
  the canonical \(R\)-linear map \(\rho : M \to (A \otimes_R R') \otimes_A M\),
  \(m \mapsto (1 \otimes 1) \otimes m\) — i.e.\ the restriction of scalars along \(\psi\)
  of the unit value of Lemma~\ref{lem:base_change_mate_unit_value}.
\end{lemma}
\begin{proof}
  \uses{lem:base_change_mate_fstar_reindex_legs, lem:base_change_mate_unit_value,
        lem:base_change_mate_codomain_read, lem:pullback_fst_snd_specMap_tensor,
        lem:pushforward_spec_tilde_iso, lem:pullbackPushforward_unit_comp,
        def:base_change_mate_inner_value}
  The concrete identity is the instantiation of the abstract variable-legs reindex
  Lemma~\ref{lem:base_change_mate_fstar_reindex_legs} at the literal pullback projection legs
  \(g' = \operatorname{pr}_1\), \(f' = \operatorname{pr}_2\): with the codomain read
  Lemma~\ref{lem:base_change_mate_codomain_read} being definitionally the variable-legs read at
  those legs, the concrete statement follows directly. The remainder of this proof records the
  mathematical content of that abstract identity.

  Throughout, let \(e : X' \xrightarrow{\sim} \operatorname{Spec}(A \otimes_R R')\) be the
  canonical isomorphism of Lemma~\ref{lem:pullbackSpecIso_mathlib}, and
  \(\iota_A : A \to A \otimes_R R'\), \(\iota_{R'} : R' \to A \otimes_R R'\) the tensor
  inclusions, so that by Lemma~\ref{lem:pullback_fst_snd_specMap_tensor} the two projection legs
  factor as
  \[
    g' \;=\; e \circ \operatorname{Spec}\iota_A,
    \qquad
    f' \;=\; e \circ \operatorname{Spec}\iota_{R'} .
  \]
  The inner composite \(\theta_{\mathrm{in}}\) is the four-factor map
  \[
    f_{*}\bigl(\eta^{g'}_{\widetilde M}\bigr) \;\circ\;
    (\mathrm{pushforwardComp}_{g', f})^{\mathrm{hom}} \;\circ\;
    (\mathrm{pushforwardCongr}\,w)^{\mathrm{hom}} \;\circ\;
    (\mathrm{pushforwardComp}_{f', g})^{\mathrm{inv}},
  \]
  with \(f = \operatorname{Spec}\varphi\), \(g = \operatorname{Spec}\psi\),
  \(\eta^{g'}\) the unit of the \((g')\)-adjunction, and
  \(w : g' \circ f = f' \circ g\) the commutativity of the square.

  \emph{Why a direct substitution fails.} The identity cannot be obtained by substituting the
  leg factorizations directly, because the legs \(g'\) and \(f'\) appear in multiple positions
  simultaneously: as the index of the adjunction \(\eta^{g'}\), inside the equality proof
  \(w\) parametrizing \(\mathrm{pushforwardCongr}\,w\), inside the argument
  \(f'_{*}\bigl((g')^{*}\widetilde M\bigr)\) of the codomain comparison
  \(\gamma_{\psi}\) (Lemma~\ref{lem:gammaPushforwardIso}), and inside the opaque codomain
  identification \(\Theta_{\mathrm{tgt}} = \) Lemma~\ref{lem:base_change_mate_codomain_read}.
  The reindex is therefore restructured into three sub-steps so that the legs become explicit
  free variables; it is \emph{not} a single rewrite.

  \emph{(i) Abstract variable-legs restatement.} Restate the codomain read — and, with it, the
  whole four-factor chain — for \emph{generic} morphisms \(g', f'\) carrying the cone-leg
  equalities \(g' = e \circ \operatorname{Spec}\iota_A\) and
  \(f' = e \circ \operatorname{Spec}\iota_{R'}\) as explicit hypotheses, in place of the literal
  projections \(\operatorname{pr}_1, \operatorname{pr}_2\). In this abstract form the legs are
  free variables, not terms pinned by the pullback's universal property, so each can be
  eliminated by substituting its defining equation: the equality proof \(w\), the adjunction
  index, and the \(\gamma_{\psi}\)/\(\Theta_{\mathrm{tgt}}\) arguments are all expressed in the
  free leg variables and transform uniformly under the substitution. This restatement converts
  the blocked rewrite into a legitimate one; it is the structural heart of the step.

  \emph{(ii) \(\Gamma\)-collapse of the transparent coherences.} On global sections, read through
  the \(\Gamma\)-pushforward dictionaries (Lemma~\ref{lem:gammaPushforwardTildeIso},
  Lemma~\ref{lem:gammaPushforwardIso}), the two pushforward-composition factors and the
  congruence factor are transparent. The composition coherences collapse to the identity —
  \((\mathrm{pushforwardComp}_{g',f})^{\mathrm{hom}}\) and
  \((\mathrm{pushforwardComp}_{f',g})^{\mathrm{inv}}\) evaluate, on the top open, to the identity
  map on restriction of scalars (their section-level values are the identity) — and they
  serve only to rebracket the iterated pushforward along the two factorizations
  \((g' \circ f)_* = (f' \circ g)_*\) of the square's composite. (The
  \((\mathrm{pushforwardComp}_{g',f})^{\mathrm{hom}}\) collapse is a transparent re-association
  that is most naturally accounted for together with the leg-reindex of step~(iii) rather than at
  this point; the \(\Gamma\)-collapse identity itself is the same.) The congruence factor collapses
  to a canonical transport along the equality of opens: its section-level value
  \((\mathrm{pushforwardCongr}\,w)^{\mathrm{hom}}\) is the
  presheaf-restriction along the canonical equality of opens, i.e.\ a repackaging along
  that equality carrying no substantive content. Thus after the \(\Gamma\)-collapse the only
  surviving non-trivial factor is \(f_{*}\bigl(\eta^{g'}_{\widetilde M}\bigr)\) read on global
  sections.

  \emph{(iii) Reduction to Seam~1 via the leg-reindex engine.} With the legs substituted
  (step~(i)) to \(g' = e \circ \operatorname{Spec}\iota_A\), the leg-reindex engine
  Lemma~\ref{lem:pullbackPushforward_unit_comp}, instantiated at \(a = e\),
  \(b = \operatorname{Spec}\iota_A\), \(N = \widetilde M\), rewrites the unit factor
  \(\eta^{g'} = \eta^{e \circ \operatorname{Spec}\iota_A}\), together with the
  \((\mathrm{pushforwardComp}_{e,\,\operatorname{Spec}\iota_A})\) coherence, into the affine
  \((\operatorname{Spec}\iota_A)\)-unit followed by \((\operatorname{Spec}\iota_A)_{*}\) of the
  \(e\)-unit. Since \(e\) is an isomorphism of schemes, its pullback functor is an equivalence
  (Lemma~\ref{lem:pullback_isEquivalence_of_iso}) and the \(e\)-unit is invertible, hence an
  isomorphism absorbed into the codomain identification \(\Theta_{\mathrm{tgt}}\). The surviving
  \((\operatorname{Spec}\iota_A)\)-unit has section value over \(\operatorname{Spec} A\) equal to
  the algebraic unit \(\eta_M : m \mapsto (1 \otimes 1) \otimes m\) of Seam~1
  (Lemma~\ref{lem:base_change_mate_unit_value}). Composing the surviving factor with the
  codomain identification \(\Theta_{\mathrm{tgt}}\)
  (Lemma~\ref{lem:base_change_mate_codomain_read}) and reading over \(\operatorname{Spec} R\) by
  restriction of scalars along \(\psi\) identifies all four factors with the restriction of
  scalars along \(\psi\) of the codomain-read transport of the unit value, namely
  \(\rho : m \mapsto (1 \otimes 1) \otimes m\)
  (Definition~\ref{def:base_change_mate_inner_value}) — the
  ``\(X' = \operatorname{Spec}(R' \otimes_R A)\), \(\mathcal F'\) is
  \((R' \otimes_R A) \otimes_A M\)'' bookkeeping of the source. Each of (i)--(iii) extracts as a
  named lemma; together they close the reindex without ever rewriting the bare legs.
\end{proof}

% --- Seam A: the inner-value identity Γ_R(θ_in) = ρ. ---
% The (g^* ⊣ g_*)-transpose crux lem:base_change_mate_gstar_transpose needs the inner
% pushforward comparison θ_in, read on Spec R sections, to equal ρ : m ↦ (1 ⊗ 1) ⊗ m
% (lem:base_change_mate_inner_value_eq, Seam A). This identity is realised THROUGH the
% leg-parametrised reindex lem:base_change_mate_fstar_reindex_legs, instantiated at the literal
% projection legs via lem:base_change_mate_fstar_reindex: leg substitution → unit expansion
% (unitExpand) → four-factor Γ-distribution (gammaDistribute) → eCancel telescoping against the
% codomain read (the three one-cancellation atoms below, assembled in
% lem:base_change_mate_inner_eCancel_assemble) → Seam 1 → ρ.
% The inline alternative — distributing the bare (g')-unit AT the literal projection legs,
% without first passing to the leg-parametrised form — is WALLED: the leg pr_1 is only
% propositionally (not definitionally) equal to e ∘ Spec ιA, and the codomain read is a closed
% construction whose type bakes in the leg, so the unit expansion meets a dependent-motive
% obstruction at the literal leg. The proof must therefore pass to the leg-parametrised form
% first. The four-factor distribution (A-1) and the three eCancel atoms (A-2a/b/c) are correct
% and reused UNCHANGED; only the plumbing — via _legs, not inline — changes.
% The generator close (Seam B) follows as an independent lemma.

\begin{lemma}

\leanok
[(A-1) Distribute the inner comparison into its four $\Gamma$-image factors]
  \label{lem:base_change_mate_inner_unitReduce}
  % NOTE: superseded by the conjugate-side re-encoding (Open Q2). This assembly-narrative node for the
  % abandoned direct-on-sections route no longer feeds the live `_legs` proof. Lean pin (the subsuming
  % `_legs` theorem) retained.
  \lean{AlgebraicGeometry.base_change_mate_fstar_reindex_legs}
  % LEAN INTERNAL: this is an assembly-narrative node with no dedicated Lean declaration; its
  % content — the four-factor Γ-distribution — is realised inside the leg-parametrised reindex
  % `base_change_mate_fstar_reindex_legs` (step (iii), links (1)-(2)), so the `\lean` pin names
  % that subsuming theorem (mirrors the convention of the eCancel narrative nodes; keeps this
  % dependency-bearing node wired into the graph rather than orphaning the proved atoms it
  % `\uses`). The four-factor unit distribution stated here is the unit expansion
  % `lem:base_change_mate_fstar_reindex_legs_unitExpand` followed by its (Spec φ)_*–Γ distribution
  % `lem:base_change_mate_fstar_reindex_legs_gammaDistribute` — equivalently the PROVED engine
  % `lem:pullbackPushforward_unit_comp` (`AlgebraicGeometry.pullbackPushforward_unit_comp`)
  % specialised at `a := e.hom` (the `pullbackSpecIso`), `b := Spec ιA`, `N := tilde M` (so
  % `g' = e.hom ≫ Spec ιA`), then carried through the two functors `(Spec φ)_*` and `Γ` by
  % functoriality (`Functor.map_comp`). IMPORTANT: this distribution must be performed in the
  % leg-parametrised form, where `g'` is identified with the composite `e.hom ≫ Spec ιA`
  % DEFINITIONALLY; the inline version, distributing at the literal projection leg `pr_1` (only
  % propositionally equal to that composite), is walled by a dependent-motive obstruction. There
  % is no separate obligation and hence no `% LEAN SIGNATURE`.
  \uses{lem:base_change_mate_fstar_reindex_legs_unitExpand,
        lem:base_change_mate_fstar_reindex_legs_gammaDistribute,
        lem:gammaMap_pushforwardComp_inv_eq_id, lem:gammaMap_pushforwardCongr_hom,
        lem:pullback_fst_snd_specMap_tensor}
  Fix \(\psi : R \to R'\), \(\varphi : R \to A\) and an \(A\)-module \(M\), and work in the
  \emph{concrete} pullback square \(X' = \operatorname{pullback}(\operatorname{Spec}\varphi,
  \operatorname{Spec}\psi)\) with \(g' = \operatorname{pr}_1\), \(f' = \operatorname{pr}_2\) the
  literal projections, \(e : X' \xrightarrow{\sim} \operatorname{Spec}(A \otimes_R R')\) the
  canonical isomorphism and \(\iota_A : A \to A \otimes_R R'\),
  \(\iota_{R'} : R' \to A \otimes_R R'\) the tensor inclusions, so that the projections factor as
  \(g' = e \circ \operatorname{Spec}\iota_A\), \(f' = e \circ \operatorname{Spec}\iota_{R'}\)
  (Lemma~\ref{lem:pullback_fst_snd_specMap_tensor}). Write \(\theta_{\mathrm{in}}\) for the inner
  composite built from the \((g')\)-unit and the pushforward pseudofunctor coherences
  (\(\mathrm{pushforwardComp}\) twice, \(\mathrm{pushforwardCongr}\) once). Then on global sections
  over \(\operatorname{Spec} R\), after the two transparent coherences collapse to the identity
  (Lemma~\ref{lem:gammaMap_pushforwardComp_inv_eq_id},
  Lemma~\ref{lem:gammaMap_pushforwardCongr_hom}), the \(\Gamma\)-image
  \(\Gamma\bigl((\operatorname{Spec}\varphi)_{*}\theta_{\mathrm{in}}\bigr)\) distributes — by the
  unit expansion of Lemma~\ref{lem:base_change_mate_fstar_reindex_legs_unitExpand} and its
  \((\operatorname{Spec}\varphi)_{*}\)--\(\Gamma\) distribution of
  Lemma~\ref{lem:base_change_mate_fstar_reindex_legs_gammaDistribute} — into the ordered composite
  of the four \(\Gamma\)-image factors
  \[
    \Gamma\bigl((\operatorname{Spec}\varphi)_{*}\eta^{\operatorname{Spec}\iota_A}_{\widetilde M}\bigr)
    \,\circ\,
    \Gamma\bigl((\operatorname{Spec}\varphi)_{*}(\operatorname{Spec}\iota_A)_{*}\eta^{e}\bigr)
    \,\circ\,
    \Gamma\bigl((\operatorname{Spec}\varphi)_{*}(\mathrm{pushforwardComp}_{e,\operatorname{Spec}\iota_A})^{\mathrm{hom}}\bigr)
    \,\circ\,
    \Gamma\bigl((\operatorname{Spec}\varphi)_{*}(g')_{*}(\mathrm{pullbackComp}_{e,\operatorname{Spec}\iota_A})^{\mathrm{hom}}\bigr),
  \]
  read left to right.
\end{lemma}
\begin{proof}
  \uses{lem:base_change_mate_fstar_reindex_legs_unitExpand,
        lem:base_change_mate_fstar_reindex_legs_gammaDistribute,
        lem:gammaMap_pushforwardComp_inv_eq_id, lem:gammaMap_pushforwardCongr_hom,
        lem:pullback_fst_snd_specMap_tensor}
  Identify the projection legs with the affine \(\operatorname{Spec}\)-maps of the tensor
  inclusions by Lemma~\ref{lem:pullback_fst_snd_specMap_tensor}; this exposes the composite leg
  \(g' = e \circ \operatorname{Spec}\iota_A\) needed to engage the leg-reindex atoms. Applying the
  global-sections functor to the four-factor composite \(\theta_{\mathrm{in}}\) and collapsing the
  inverse pushforward-composition factor (Lemma~\ref{lem:gammaMap_pushforwardComp_inv_eq_id}) and
  the congruence factor (Lemma~\ref{lem:gammaMap_pushforwardCongr_hom}) to the identity leaves the
  unit factor \((\operatorname{Spec}\varphi)_{*}\eta^{g'}_{\widetilde M}\) and the surviving
  \((\mathrm{pushforwardComp}_{g',\operatorname{Spec}\varphi})^{\mathrm{hom}}\). Expanding the bare
  \((g')\)-unit by Lemma~\ref{lem:base_change_mate_fstar_reindex_legs_unitExpand} and distributing
  through \((\operatorname{Spec}\varphi)_{*}\) and \(\Gamma\) by
  Lemma~\ref{lem:base_change_mate_fstar_reindex_legs_gammaDistribute} — both functors, so each
  preserves composites — yields the stated four-factor form. No further content.
\end{proof}

% --- The three one-cancellation atoms of the eCancel telescoping. ---
% Each is an independently-formalizable lemma whose proof is essentially a single move; together
% with the proved unit-distribution engine `lem:pullbackPushforward_unit_comp` they telescope the
% inner composite against the codomain read down to the lone affine `(Spec ιA)`-unit.

\begin{lemma}

\leanok
[(A-2a) The $e$-unit is an isomorphism]
  \label{lem:base_change_mate_inner_eCancel_eUnit}
  % NOTE: superseded by the conjugate-side re-encoding (Open Q2). The e-unit invertibility it records
  % is now consumed inside the conjugate-component identity of `_legs`, not as a standalone eCancel
  % atom. Lean pin retained.
  \lean{AlgebraicGeometry.base_change_mate_inner_eCancel_eUnit}
  \uses{lem:pullback_isEquivalence_of_iso}
  Let \(e : X \to Y\) be an isomorphism of schemes and \(N\) a \(Y\)-module. Then the unit
  \(\eta^{e}_{N} : N \to e_{*}\,e^{*}N\) of the \((e^{*} \dashv e_{*})\)-adjunction is an
  isomorphism. This is the single fact enabling cancellation~(1): in the inner telescoping it is
  the \(e\)-unit factor \((\operatorname{Spec}\iota_A)_{*}\eta^{e}\) (here \(e\) is the canonical
  square isomorphism \(\mathrm{pullbackSpecIso}\)), which the codomain read
  \(\Theta_{\mathrm{tgt}}\) bakes in as its inverse \(\eta^{e}\)-piece, so the two cancel.
\end{lemma}
\begin{proof}
  \uses{lem:pullback_isEquivalence_of_iso}
  Since \(e\) is an isomorphism of schemes, its pullback functor \(e^{*} = \mathrm{pullback}\,e\)
  is an equivalence (Lemma~\ref{lem:pullback_isEquivalence_of_iso}). The unit of an adjunction whose
  left adjoint is an equivalence is an isomorphism; hence \(\eta^{e}_{N}\) is invertible.
\end{proof}

\begin{lemma}

\leanok
[(A-2b) The surviving $\mathrm{pushforwardComp}$ factor has identity $\Gamma$-image]
  \label{lem:base_change_mate_inner_eCancel_pushforwardComp}
  % NOTE: superseded by the conjugate-side re-encoding (Open Q2). Its content (the Γ-collapse of the
  % pushforward-composition coherence) survives only through the live
  % lem:gammaMap_pushforwardComp_hom_eq_id, which the conjugate-side `_legs` proof uses directly. Lean
  % pin retained.
  \lean{AlgebraicGeometry.base_change_mate_inner_eCancel_pushforwardComp}
  \uses{lem:gammaMap_pushforwardComp_hom_eq_id}
  For composable scheme morphisms \(a : X_1 \to X_2\), \(b : X_2 \to \operatorname{Spec} A\), a ring
  map \(\varphi : R \to A\), and a module \(M\) on \(X_1\), the \(\Gamma\)-image over
  \(\operatorname{Spec} R\) of the \((\operatorname{Spec}\varphi)_{*}\)-image of the
  \(\mathrm{pushforwardComp}_{a,b}\) hom-coherence is the identity. This discharges
  cancellation~(2): the \(\mathrm{pushforwardComp}_{e,\operatorname{Spec}\iota_A}\) factor surviving
  the unit distribution drops out under \(\Gamma\bigl((\operatorname{Spec}\varphi)_{*}(-)\bigr)\),
  exactly as the unwrapped \(\mathrm{pushforwardComp}_{g',\operatorname{Spec}\varphi}\) factor drops
  by Lemma~\ref{lem:gammaMap_pushforwardComp_hom_eq_id}.
\end{lemma}
\begin{proof}
  \uses{lem:gammaMap_pushforwardComp_hom_eq_id}
  The section value of the \(\mathrm{pushforwardComp}\) hom-coherence is the identity at every open,
  so \((\mathrm{pushforwardComp}_{a,b})^{\mathrm{hom}}\) is the identity morphism; its
  \((\operatorname{Spec}\varphi)_{*}\)-image is then the identity by functoriality, and its further
  \(\Gamma\)-image is the identity by functoriality again. Two applications of the functorial
  identity law.
\end{proof}

\begin{lemma}

\leanok
[(A-2c) The $\mathrm{pullbackComp}$ factor cancels its inverse in the codomain read]
  \label{lem:base_change_mate_inner_eCancel_pullbackComp}
  % NOTE: superseded by the conjugate-side re-encoding (Open Q2). The pullbackComp hom--inverse
  % cancellation is now absorbed by lem:conjugateEquiv_pullbackComp_inv_mathlib on the conjugate side.
  % Lean pin retained.
  \lean{AlgebraicGeometry.base_change_mate_inner_eCancel_pullbackComp}
  \uses{lem:base_change_mate_codomain_read, lem:pullback_fst_snd_specMap_tensor}
  In the concrete pullback square, with \(e = \mathrm{pullbackSpecIso}\) the canonical isomorphism
  \(X' \xrightarrow{\sim} \operatorname{Spec}(A \otimes_R R')\) and \(\iota_A\) the left tensor
  inclusion, the hom and inverse of the pseudofunctor coherence
  \(\mathrm{pullbackComp}_{e,\operatorname{Spec}\iota_A}\) on \(\widetilde M\) compose to the
  identity. This is cancellation~(3): the codomain read \(\Theta_{\mathrm{tgt}}\)
  (Lemma~\ref{lem:base_change_mate_codomain_read}) is, by construction, assembled with exactly
  \((\mathrm{pullbackComp}_{e,\operatorname{Spec}\iota_A})^{\mathrm{inv}}\) (inside its leg-rewrite
  isomorphism \(\mathrm{iso\_g}\)), against which the trailing
  \((\mathrm{pullbackComp}_{e,\operatorname{Spec}\iota_A})^{\mathrm{hom}}\) factor of the unit
  distribution cancels.
\end{lemma}
\begin{proof}
  \uses{lem:base_change_mate_codomain_read}
  Both factors are the hom and inverse of one and the same isomorphism
  \(\mathrm{pullbackComp}_{e,\operatorname{Spec}\iota_A}\) evaluated at \(\widetilde M\); their
  composite is the identity by the hom--inverse law for isomorphisms. The matching against
  \(\Theta_{\mathrm{tgt}}\) is definitional: the codomain read is built from the same
  \((e, \operatorname{Spec}\iota_A)\) leg factorization
  (Lemma~\ref{lem:pullback_fst_snd_specMap_tensor}), so the inverse it bakes in is literally this
  isomorphism's inverse. A single application of the hom--inverse identity.
\end{proof}

\begin{lemma}

\leanok
[(A-2) Cancel the $e$-factors against the codomain read]
  \label{lem:base_change_mate_inner_eCancel}
  \lean{AlgebraicGeometry.base_change_mate_fstar_reindex_legs}
  % NOTE: superseded by the conjugate-side re-encoding (Open Q2). This mid-level assembly-narrative
  % node for the abandoned direct-on-sections eCancel telescoping no longer feeds the live `_legs`
  % proof. Lean pin (the subsuming `_legs` theorem) retained.
  % LEAN INTERNAL: mid-level assembly-narrative node with no dedicated Lean declaration. It states
  % the e-factor cancellation in the CONCRETE square (codomain read = base_change_mate_codomain_read,
  % literal projection legs). The LIVE, closeable realization of this cancellation is
  % `lem:base_change_mate_inner_eCancel_assemble`, which performs the same three one-cancellation
  % atoms against the LEG-PARAMETRISED codomain read `base_change_mate_codomain_read_legs` inside
  % `base_change_mate_fstar_reindex_legs` (step (iii), link (3)); the `\lean` pin names that
  % subsuming theorem. The inline form stated here — cancelling at the literal projection legs — is
  % walled (the legs are only propositionally equal to their composite form, a dependent-motive
  % obstruction), which is why the work routes through the `_legs` form. The cancellation is the
  % ordered application of the three one-cancellation atoms below — `..._eUnit` (cancellation 1),
  % `..._pushforwardComp` (cancellation 2), `..._pullbackComp` (cancellation 3) — to the four-factor
  % unit distribution supplied by `lem:pullbackPushforward_unit_comp`. (Its post-cancellation type —
  % the surviving `Spec ιA`-unit conjugated by the tilde/Γ dictionaries — resisted a clean
  % standalone signature; the three-atom split removes the need to name it.)
  \uses{lem:base_change_mate_inner_eCancel_eUnit,
        lem:base_change_mate_inner_eCancel_pushforwardComp,
        lem:base_change_mate_inner_eCancel_pullbackComp,
        lem:pullbackPushforward_unit_comp, lem:gammaMap_pushforwardComp_hom_eq_id,
        lem:base_change_mate_codomain_read}
  In the situation of Lemma~\ref{lem:base_change_mate_inner_unitReduce}, post-composing the
  four-factor distribution with the codomain identification \(\Theta_{\mathrm{tgt}}\) of
  Lemma~\ref{lem:base_change_mate_codomain_read} cancels three of the four factors against the
  matching \(e\)-pieces baked into \(\Theta_{\mathrm{tgt}}\), in the order:
  \begin{enumerate}
    \item the \(e\)-unit factor \((\operatorname{Spec}\iota_A)_{*}\eta^{e}\) is an isomorphism
      (Lemma~\ref{lem:base_change_mate_inner_eCancel_eUnit}) and cancels the inverse \(e\)-unit
      built into \(\Theta_{\mathrm{tgt}}\);
    \item the surviving
      \((\mathrm{pushforwardComp}_{e,\operatorname{Spec}\iota_A})^{\mathrm{hom}}\) has identity
      \(\Gamma\)-image (Lemma~\ref{lem:base_change_mate_inner_eCancel_pushforwardComp}) and drops;
    \item the trailing
      \((\mathrm{pullbackComp}_{e,\operatorname{Spec}\iota_A})^{\mathrm{hom}}\) cancels its inverse
      inside \(\Theta_{\mathrm{tgt}}\)
      (Lemma~\ref{lem:base_change_mate_inner_eCancel_pullbackComp}).
  \end{enumerate}
  After cancellation, only the affine \((\operatorname{Spec}\iota_A)\)-unit
  \(\eta^{\operatorname{Spec}\iota_A}_{\widetilde M}\), conjugated by the affine tilde/\(\Gamma\)
  dictionaries over \(\operatorname{Spec} A\), survives.
\end{lemma}
\begin{proof}
  \uses{lem:base_change_mate_inner_eCancel_eUnit,
        lem:base_change_mate_inner_eCancel_pushforwardComp,
        lem:base_change_mate_inner_eCancel_pullbackComp,
        lem:pullbackPushforward_unit_comp, lem:base_change_mate_codomain_read}
  Distribute the bare \((g')\)-unit by the proved engine
  Lemma~\ref{lem:pullbackPushforward_unit_comp} (with \(a := e^{\mathrm{hom}}\),
  \(b := \operatorname{Spec}\iota_A\), \(N := \widetilde M\), so \(g' = e^{\mathrm{hom}} \circ
  \operatorname{Spec}\iota_A\)) and read through \(\Gamma\bigl((\operatorname{Spec}\varphi)_{*}(-)\bigr)\)
  by functoriality, giving the four ordered factors. Then apply the three atoms in turn:
  Lemma~\ref{lem:base_change_mate_inner_eCancel_eUnit} cancels the \(e\)-unit against the inverse
  \(e\)-unit of \(\Theta_{\mathrm{tgt}}\); Lemma~\ref{lem:base_change_mate_inner_eCancel_pushforwardComp}
  drops the \(\mathrm{pushforwardComp}\) factor; and
  Lemma~\ref{lem:base_change_mate_inner_eCancel_pullbackComp} cancels the \(\mathrm{pullbackComp}\)
  factor against its inverse in \(\Theta_{\mathrm{tgt}}\). The lone survivor is the affine
  \((\operatorname{Spec}\iota_A)\)-unit read through the affine dictionaries, which is the
  load-bearing telescoping of the reindex.

  \emph{Order of operations (load-bearing).} The unit distribution of
  Lemma~\ref{lem:pullbackPushforward_unit_comp} must be performed while \(g'\) is still presented as
  the \emph{free composite} \(e^{\mathrm{hom}} \circ \operatorname{Spec}\iota_A\) — i.e.\ before the
  generic-square legs \((g', f')\) are identified with the concrete pullback projections
  \((\operatorname{pr}_1, \operatorname{pr}_2)\) by Lemma~\ref{lem:base_change_mate_codomain_read}.
  Distributing first keeps each of the four factors in the free \(a\!\circ\!b\) shape against which the
  three atoms are stated, so they apply directly. If instead one first reads through the codomain
  identification (substituting the literal projection legs and collapsing the
  \(\mathrm{pushforwardComp}\) bookkeeping), the surviving \((g')\)-unit becomes the rigid literal
  \((\mathrm{pullbackSpecIso})^{\mathrm{hom}} \circ \operatorname{Spec}(\mathrm{includeLeft})\), whose
  shape no longer presents as an \((a\!\circ\!b)\)-unit and against which the atoms are inapplicable.
  Hence: distribute the unit on the free composite first; cancel with the three atoms; \emph{then}
  read through Lemma~\ref{lem:base_change_mate_codomain_read}.
\end{proof}

\begin{lemma}


\leanok
[(A-2$'$) Assemble the eCancel telescoping in the leg-parametrised form]
  \label{lem:base_change_mate_inner_eCancel_assemble}
  \lean{AlgebraicGeometry.base_change_mate_fstar_reindex_legs}
  % NOTE: superseded by the conjugate-side re-encoding (Open Q2). This was the LIVE realization of the
  % eCancel telescoping inside `_legs` under the abandoned direct-on-sections route; the live `_legs`
  % proof is now the conjugate-component identity, which does not route through this telescoping. The
  % consuming lemmas (gstar_transpose) no longer `\uses` it. Lean pin (the subsuming `_legs` theorem)
  % retained.
  % NOTE: This is the final telescoping inside `base_change_mate_fstar_reindex_legs` (step iii,
  % link 3); the \lean{} pin names that subsuming theorem. It is the leg-parametrised analogue of
  % the inline eCancel: the same three one-cancellation atoms, run against the leg-parametrised
  % codomain read, where the leg identification \(g' = e \circ \mathrm{Spec}\,\iota_A\) is
  % definitional, so each atom's factor matches and the cancellation proceeds.
  % NOTE (mechanism): making the legs definitional removes the dependent-motive wall of the
  % literal-leg form, but a SECOND diamond survives — the nested-image vs composed-functor object
  % diamond on the `X.Modules` `CategoryStruct.comp`/`Functor.map` instance. After Γ is distributed
  % over the composite the factor to collapse sits over the nested-image object `G.obj (H.obj M)`,
  % while the collapsing atom fixes the composed-functor object `(H ⋙ G).obj M`; these are defeq but
  % not syntactic. Each factor-equation is lifted into the ambient composite one neighbour at a time by
  % congruence in the composition operation, and the chain is closed on the definitional bridge.
  \uses{lem:base_change_mate_inner_eCancel_eUnit,
        lem:base_change_mate_inner_eCancel_pushforwardComp,
        lem:base_change_mate_inner_eCancel_pullbackComp,
        lem:base_change_mate_codomain_read_legs,
        lem:gammaMap_pushforwardComp_hom_eq_id,
        lem:base_change_mate_unit_value}
  In the leg-parametrised reindex Lemma~\ref{lem:base_change_mate_fstar_reindex_legs}, after the
  unit expansion and four-factor \(\Gamma\)-distribution (step~(iii), links~(1)--(2)), the four
  distributed \(\Gamma\)-image factors cancel against the unfolded leg-parametrised codomain read
  Lemma~\ref{lem:base_change_mate_codomain_read_legs}, in the order:
  \begin{enumerate}
    \item the \(e\)-unit factor \((\operatorname{Spec}\iota_A)_{*}\eta^{e}\) is an isomorphism
      (Lemma~\ref{lem:base_change_mate_inner_eCancel_eUnit}) and cancels the inverse \(e\)-unit
      built into the codomain read;
    \item the surviving
      \((\mathrm{pushforwardComp}_{e,\operatorname{Spec}\iota_A})^{\mathrm{hom}}\) has identity
      \(\Gamma\)-image (Lemma~\ref{lem:base_change_mate_inner_eCancel_pushforwardComp}; the
      unwrapped factor by Lemma~\ref{lem:gammaMap_pushforwardComp_hom_eq_id}) and drops;
    \item the trailing
      \((\mathrm{pullbackComp}_{e,\operatorname{Spec}\iota_A})^{\mathrm{hom}}\) cancels its inverse
      built into the codomain read
      (Lemma~\ref{lem:base_change_mate_inner_eCancel_pullbackComp}).
  \end{enumerate}
  The lone survivor is the affine \((\operatorname{Spec}\iota_A)\)-unit
  \(\eta^{\operatorname{Spec}\iota_A}_{\widetilde M}\), whose section value over
  \(\operatorname{Spec} A\) is the algebraic unit \(\eta_M : m \mapsto (1 \otimes 1) \otimes m\) of
  Seam~1 (Lemma~\ref{lem:base_change_mate_unit_value}). This is the residual telescoping that closes
  the leg-parametrised reindex; it is the single remaining target of Seam~A.
\end{lemma}
\begin{proof}
  \uses{lem:base_change_mate_inner_eCancel_eUnit,
        lem:base_change_mate_inner_eCancel_pushforwardComp,
        lem:base_change_mate_inner_eCancel_pullbackComp,
        lem:base_change_mate_codomain_read_legs,
        lem:gammaMap_pushforwardComp_hom_eq_id,
        lem:base_change_mate_unit_value}
  The four distributed \(\Gamma\)-image factors are cancelled against the unfolded leg-parametrised
  codomain read Lemma~\ref{lem:base_change_mate_codomain_read_legs} by the three one-cancellation
  atoms. Carrying the legs as free parameters and identifying them with
  \(e \circ \operatorname{Spec}\iota_A\) definitionally removes the dependent-motive obstruction of
  the literal-leg form.

  \emph{The residual object-form diamond.} After distributing \(\Gamma\) over the composite, the
  factor to be collapsed sits over the \emph{nested-image} object \(G(H(M))\) — the form produced
  by distributing a functor over a composite of morphisms — whereas the collapsing atom fixes the
  same object in its \emph{composed-functor} form \((H \circ G)(M)\), the form in which the
  pushforward-composition coherence isomorphism presents its (co)domain. These two object
  presentations are definitionally but not syntactically equal, creating a diamond that recurs at
  each of the three cancellations.

  \emph{The resolution: congruence lifting.} Each factor-equation is lifted into the ambient
  composite \emph{one neighbour at a time}: a factor identity \(u = u'\) is promoted to
  \(c \circ u = c \circ u'\), or to \(u \circ c = u' \circ c\), by congruence in the composition
  operation with the neighbouring factor \(c\) held fixed, and the promoted steps are chained in
  sequence. The chain is then closed on the definitional bridge between the nested-image and
  composed-functor presentations of the shared object, absorbing the object-form diamond at that
  single closing step. The three atoms
  (Lemmas~\ref{lem:base_change_mate_inner_eCancel_eUnit},
  \ref{lem:base_change_mate_inner_eCancel_pushforwardComp},
  \ref{lem:base_change_mate_inner_eCancel_pullbackComp}) are supplied in the
  \(\Gamma \circ (\operatorname{Spec}\varphi)_{*}\) form and are spliced directly into the composite
  by this congruence chaining.

  \emph{The three cancellations.} Cancellation~(1): the \(e\)-unit is invertible
  (Lemma~\ref{lem:base_change_mate_inner_eCancel_eUnit}) and meets its inverse inside the codomain
  read. Cancellation~(2): the surviving pushforward-composition coherence has identity
  \(\Gamma\)-image (Lemma~\ref{lem:base_change_mate_inner_eCancel_pushforwardComp}, whose
  section-level content is the collapse Lemma~\ref{lem:gammaMap_pushforwardComp_hom_eq_id}).
  Cancellation~(3): the trailing pullback-comparison factor cancels the inverse the codomain read
  bakes in (Lemma~\ref{lem:base_change_mate_inner_eCancel_pullbackComp}). The lone survivor is the
  affine \((\operatorname{Spec}\iota_A)\)-unit, evaluated by Seam~1
  (Lemma~\ref{lem:base_change_mate_unit_value}) to \(\rho\).
\end{proof}

\begin{lemma}

\leanok
[(A) The inner comparison reads as $\rho$ on $\operatorname{Spec} R$ sections]
  \label{lem:base_change_mate_inner_value_eq}
  \lean{AlgebraicGeometry.base_change_mate_inner_value_eq}
  \uses{lem:base_change_mate_fstar_reindex,
        lem:base_change_mate_codomain_read, lem:pullback_fst_snd_specMap_tensor,
        lem:base_change_mate_unit_value,
        lem:pushforward_spec_tilde_iso, def:base_change_mate_inner_value}
  % SOURCE: [Stacks Project], Cohomology of Schemes, Lemma `lemma-affine-base-change'
  % (\S\,Cohomology and base change, I), proof, the "X' = Spec(R' \otimes_R A) and F' is the
  % qcoh sheaf associated to (R' \otimes_R A) \otimes_A M" step (read from
  % references/stacks-coherent.tex, L928--932).
  % SOURCE QUOTE: "We use Schemes, Lemma \ref{schemes-lemma-widetilde-pullback}
  % to describe pullbacks and pushforwards of $\mathcal{F}$.
  % Namely, $X' = \Spec(R' \otimes_R A)$ and
  % $\mathcal{F}'$ is the quasi-coherent sheaf associated
  % to $(R' \otimes_R A) \otimes_A M$."
  \textit{Source: Stacks Project, Cohomology of Schemes, Lemma ``Affine base change''
  (proof, ``describe pullbacks and pushforwards'' step).}
  In the concrete pullback square (\(g' = \operatorname{pr}_1\), \(f' = \operatorname{pr}_2\)),
  the inner composite \(\theta_{\mathrm{in}}\) of
  Lemma~\ref{lem:base_change_mate_inner_unitReduce}, read on global sections over
  \(\operatorname{Spec} R\) through the \(\Gamma\)-pushforward dictionaries
  (Lemma~\ref{lem:pushforward_spec_tilde_iso}) with the codomain pinned by
  Lemma~\ref{lem:base_change_mate_codomain_read}, equals the canonical \(R\)-linear map
  \[
    \rho : M \longrightarrow (A \otimes_R R') \otimes_A M, \qquad m \longmapsto (1 \otimes 1)\otimes m
  \]
  of Definition~\ref{def:base_change_mate_inner_value}. This is the
  ``\(X' = \operatorname{Spec}(R' \otimes_R A)\), \(\mathcal{F}'\) is
  \((R' \otimes_R A) \otimes_A M\)'' bookkeeping of the source.
\end{lemma}
\begin{proof}
  \uses{lem:base_change_mate_fstar_reindex,
        lem:base_change_mate_codomain_read, lem:pullback_fst_snd_specMap_tensor,
        lem:base_change_mate_unit_value,
        lem:pushforward_spec_tilde_iso, def:base_change_mate_inner_value}
  The inner-value identity is the instantiation of the leg-parametrised reindex
  Lemma~\ref{lem:base_change_mate_fstar_reindex} (the literal-leg form of the engine
  Lemma~\ref{lem:base_change_mate_fstar_reindex_legs}) at the projection legs
  \(g' = \operatorname{pr}_1\), \(f' = \operatorname{pr}_2\). At those legs the concrete codomain
  read Lemma~\ref{lem:base_change_mate_codomain_read} is definitionally the leg-parametrised read,
  and the leg factorizations \(g' = e \circ \operatorname{Spec}\iota_A\),
  \(f' = e \circ \operatorname{Spec}\iota_{R'}\) hold by
  Lemma~\ref{lem:pullback_fst_snd_specMap_tensor}; the reindex therefore evaluates
  \(\theta_{\mathrm{in}}\), read on \(\operatorname{Spec} R\) sections through the
  \(\Gamma\)-pushforward dictionaries (Lemma~\ref{lem:pushforward_spec_tilde_iso}), to the affine
  \((\operatorname{Spec}\iota_A)\)-unit value of Seam~1
  (Lemma~\ref{lem:base_change_mate_unit_value}), read over \(\operatorname{Spec} R\) by restriction
  of scalars along \(\psi\) and transported across the ring equation
  \(\iota_A \circ \varphi = \iota_{R'} \circ \psi\). This is exactly \(\rho\)
  (Definition~\ref{def:base_change_mate_inner_value}).

  \emph{Why the proof routes through the leg-parametrised form rather than distributing inline.} At
  the literal projection legs the bare \((g')\)-unit cannot be expanded in place: the leg
  \(\operatorname{pr}_1\) is only propositionally — not definitionally — equal to the composite
  \(e \circ \operatorname{Spec}\iota_A\) (Lemma~\ref{lem:pullback_fst_snd_specMap_tensor}), while
  the codomain read is a closed construction whose type bakes in the leg, so transporting the unit
  expansion across that propositional equality meets a dependent-motive obstruction. Passing first
  to the explicit-composite form Lemma~\ref{lem:base_change_mate_fstar_reindex_legs}, where the leg is
  identified with its composite definitionally, removes the obstruction; the reindex coherence is then
  discharged there on the conjugate side (Open Q2).
\end{proof}

\begin{lemma}

\leanok
[(B) Generator close: the base change of $\rho$ is the inverse regrouping]
  \label{lem:base_change_mate_gstar_generator_close}
  \lean{AlgebraicGeometry.base_change_mate_gstar_generator_close}
  \uses{lem:base_change_mate_regroupEquiv, def:base_change_mate_inner_value}
  % SOURCE: [Stacks Project], Cohomology of Schemes, Lemma `lemma-affine-base-change'
  % (\S\,Cohomology and base change, I), proof, the "boils down to the equality" step
  % (read from references/stacks-coherent.tex, L933--938).
  % SOURCE QUOTE: "Thus we see that the lemma boils down to the equality
  % $$ (R' \otimes_R A) \otimes_A M = R' \otimes_R M $$ as $R'$-modules."
  \textit{Source: Stacks Project, Cohomology of Schemes, Lemma ``Affine base change''
  (proof, ``boils down to the equality'' step).}
  Let \(\rho : M \to (A \otimes_R R') \otimes_A M\), \(m \mapsto (1 \otimes 1) \otimes m\), be the
  inner comparison value of Definition~\ref{def:base_change_mate_inner_value}. The extension of
  scalars along \(\psi : R \to R'\) of \(\rho\), post-composed with the algebraic counit
  \(\varepsilon^{\mathrm{alg}}\) of the extension/restriction-of-scalars adjunction along \(\psi\),
  is an \(R'\)-linear map \(R' \otimes_R M \to (A \otimes_R R') \otimes_A M\); on the generator
  \(r' \otimes m\) it returns \(r' \cdot ((1 \otimes 1) \otimes m) = (1 \otimes r') \otimes m\),
  which is exactly the value of the inverse regrouping isomorphism
  (Lemma~\ref{lem:base_change_mate_regroupEquiv}) on \(r' \otimes m\). Both maps being
  \(R'\)-linear and agreeing on the generators \(r' \otimes m\), they coincide:
  \[
    \operatorname{ext}_\psi(\rho) \,\circ\, \varepsilon^{\mathrm{alg}}
      \;=\; \operatorname{regroup}^{-1}.
  \]
\end{lemma}
\begin{proof}
  \uses{lem:base_change_mate_regroupEquiv, def:base_change_mate_inner_value}
  Both sides are \(R'\)-linear maps out of \(R' \otimes_R M\), so by extensionality it suffices to
  check the value on the generators \(r' \otimes m\). The left side sends
  \(r' \otimes m \mapsto r' \cdot \rho(m) = r' \cdot ((1 \otimes 1) \otimes m) =
  (1 \otimes r') \otimes m\), using that \(\varepsilon^{\mathrm{alg}}\) reads the extension of
  scalars by acting through the \(R'\)-scalar. The inverse regrouping isomorphism of
  Lemma~\ref{lem:base_change_mate_regroupEquiv} sends, by definition,
  \(r' \otimes m \mapsto (1 \otimes r') \otimes m\) (in the \(A \otimes_R R'\) orientation). The two
  values agree on every generator, so the maps are equal. No flatness hypothesis is used.
\end{proof}

\begin{lemma}\leanok
[Counit reading of the extended inner value (coverage --- variant of (b))]
  \label{lem:base_change_mate_extendScalars_inner_value_counit}
  \lean{AlgebraicGeometry.base_change_mate_extendScalars_inner_value_counit}
  \uses{lem:base_change_mate_gstar_generator_close, def:base_change_mate_inner_value}
  A redundant variant of \cref{lem:base_change_mate_gstar_generator_close}: it records the same identity \(\operatorname{ext}_\psi(\rho)\circ\varepsilon^{\mathrm{alg}}
  = \operatorname{regroup}^{-1}\) read through \(\mathrm{ExtendScalars.map'}\) and
  \(\mathrm{Counit.map\_apply\_one\_tmul}\) (the affine extend/restrict counit on a single generator
  \(1 \otimes x\)). It is a coverage node only; the assembly uses \cref{lem:base_change_mate_gstar_generator_close}.
\end{lemma}
\begin{proof}
  \uses{lem:base_change_mate_gstar_generator_close}
  After \(\mathrm{ext}\,x\) and unfolding \(\mathrm{ExtendScalars.map'}\), the counit
  \(\mathrm{Counit.map\_apply\_one\_tmul}\) reduces the left side to the bare inner value \(\rho(x)\), and
  \(\rho(x) = \mathrm{regroup}^{-1}(1\otimes x)\) holds definitionally (\(\mathrm{congrArg}\;\_\;\mathrm{rfl}\)),
  matching \cref{lem:base_change_mate_gstar_generator_close}.
\end{proof}

\begin{lemma}

\leanok
[(C) The geometric counit, conjugated by the dictionaries, is the algebraic counit]
  \label{lem:base_change_mate_gstar_counit_transport}
  \lean{AlgebraicGeometry.base_change_mate_gstar_counit_transport}
  \uses{lem:pullback_spec_tilde_iso, lem:gammaPushforwardNatIso,
        lem:unit_conjugateEquiv_mathlib}
  Let \(g = \operatorname{Spec}\psi\). Write \(\varepsilon_g\) for the counit of the geometric
  \((g^* \dashv g_*)\)-adjunction, \(\Pi_{\psi}\) for the pullback dictionary of
  Lemma~\ref{lem:pullback_spec_tilde_iso} (reading \(g^*\) of a tilde as extension of scalars along
  \(\psi\), built from \(\gamma_{\psi}\), Lemma~\ref{lem:gammaPushforwardNatIso}), and
  \(\varepsilon^{\mathrm{alg}}\) for the counit of the extension/restriction-of-scalars adjunction
  along \(\psi\). Then the geometric counit, conjugated by \(\Pi_{\psi}\) and the tilde--\(\Gamma\)
  counit on either side, equals the algebraic counit: as a single transport identity,
  \[
    \Pi_{\psi}^{-1} \,\circ\, g^*(\varepsilon^{\Gamma}) \,\circ\, \varepsilon_g
    \;=\;
    \widetilde{(-)}_{R'}(\varepsilon^{\mathrm{alg}}) \,\circ\, \varepsilon^{\Gamma},
  \]
  the counit-triangle (zig-zag) coherence for the composed adjunction
  \((\widetilde{(-)}_{R'} \dashv \Gamma_{R'}) . (g^* \dashv g_*)\), where \(\varepsilon^{\Gamma}\)
  denotes the tilde--\(\Gamma\) counit. This is the counit dual of the unit-across-conjugate
  coherence Lemma~\ref{lem:unit_conjugateEquiv_mathlib} used in Seam~1.
\end{lemma}
\begin{proof}
  \uses{lem:pullback_spec_tilde_iso, lem:gammaPushforwardNatIso,
        lem:unit_conjugateEquiv_mathlib}
  Instantiate the counit-across-conjugate coherence (the counit dual of
  Lemma~\ref{lem:unit_conjugateEquiv_mathlib}) at the composite adjunctions
  \(\mathrm{adj}_L = (\widetilde{(-)}_R \dashv \Gamma_R) . (g^* \dashv g_*)\) and
  \(\mathrm{adj}_R = (\,{}^{\psi}\!(-) \dashv {}_{\psi}(-)) . (\widetilde{(-)}_{R'} \dashv
  \Gamma_{R'})\), with the right-adjoint comparison \(\gamma_{\psi}\)
  (Lemma~\ref{lem:gammaPushforwardNatIso}). The conjugate of \(\gamma_{\psi}\) is, at every base
  module, the inverse of \(\Pi_{\psi}\) (Lemma~\ref{lem:pullback_spec_tilde_iso}), as in Seam~1.
  Splitting each composite counit into its tilde--\(\Gamma\) factor and its geometric / algebraic
  factor (the counit form of the composite-adjunction counit) and fusing the two splits into the
  coherence yields the stated transport identity, the dual of Seam~1's unit transport.
\end{proof}

\begin{lemma}

\leanok
[The $(g^* \dashv g_*)$ transpose of the comparison on sections]
  \label{lem:base_change_mate_gstar_transpose}
  \lean{AlgebraicGeometry.base_change_mate_gstar_transpose}
  % NOTE: This is the live residual crux of obligation 2 (the conjugate-counit route). The base-change
  % map is DEFINED as the (g^* ⊣ g_*) transpose, so any section-level evaluation unfolds via
  % `Adjunction.homEquiv_counit` to this mate form; the master counit-transport identity `huce`
  % (conjugateEquiv_counit_symm of the comparison) is landed and compiling, and the proof block below
  % records the remaining ~150-LOC telescoping (steps (a)/(b)/(c)).
  \uses{lem:base_change_mate_domain_read,
        lem:base_change_mate_codomain_read, lem:base_change_mate_regroupEquiv,
        lem:pullback_spec_tilde_iso, lem:base_change_mate_unit_value,
        lem:conjugateEquiv_counit_symm_mathlib}
  % SOURCE: [Stacks Project], Cohomology of Schemes, Lemma `lemma-affine-base-change'
  % (\S\,Cohomology and base change, I), proof, the "boils down to the equality" step
  % (read from references/stacks-coherent.tex, L933--938).
  % SOURCE QUOTE: "Thus we see that the lemma boils down to the equality
  % $$ (R' \otimes_R A) \otimes_A M = R' \otimes_R M $$ as $R'$-modules."
  \textit{Source: Stacks Project, Cohomology of Schemes, Lemma ``Affine base change''
  (proof, ``boils down to the equality'' step).}
  % LEAN SIGNATURE (LHS elaboration-checked). The factorization is Mathlib's
  %   `Adjunction.homEquiv_counit`:
  %     pushforwardBaseChangeMap (Spec.map φ) (Spec.map ψ) f' g' h.w (tilde M)
  %       = (Scheme.Modules.pullback (Spec.map ψ)).map inner
  %          ≫ (pullbackPushforwardAdjunction (Spec.map ψ)).counit.app Z
  %   (`inner`, `Z`, `f'`, `g'` as in lem:base_change_mate_fstar_reindex). The section-level
  %   content, the genuine coherence, is the equality of `ModuleCat R'` morphisms:
  %     {R R' A : CommRingCat.{u}} (ψ : R ⟶ R') (φ : R ⟶ A) (M : ModuleCat.{u} A) :
  %       (base_change_mate_domain_read ψ φ M).inv ≫
  %         (moduleSpecΓFunctor (R := R')).map
  %           ((Scheme.Modules.pullback (Spec.map ψ)).map inner ≫
  %             (pullbackPushforwardAdjunction (Spec.map ψ)).counit.app Z) ≫
  %         (base_change_mate_codomain_read ψ φ M).hom
  %         = (base_change_mate_regroupEquiv ψ φ M).inv
  % RECIPE: Seam 3 is now the live remaining crux of the affine section-level computation;
  %   it feeds `base_change_mate_section_identity` directly. It is the `extendScalars ψ`
  %   (base change along ψ) of the affine inner value `ρ : m ↦ (1 ⊗ 1) ⊗ m`.
  %   Route:
  %   1. Counit split. The base-change map factors via `Adjunction.homEquiv_counit` into
  %      `Γ(g^*(θ_in)) ≫ Γ(ε_g)` (g = Spec ψ); no opaque transpose remains.
  %   2. Conjugate by `Θ_src = base_change_mate_domain_read` (domain) and
  %      `Θ_tgt = base_change_mate_codomain_read` (codomain). `Θ_src` already bakes in
  %      `pullback_spec_tilde_iso ψ`, which reads `g^* = (Spec ψ)^*` of a tilde as
  %      `extendScalars ψ`; the residual is the counit-triangle identity / dictionary
  %      naturality, yielding `extendScalars ψ ∘ ρ` (ρ : m ↦ (1 ⊗ 1) ⊗ m, the affine unit
  %      value `base_change_mate_unit_value`).
  %   3. Both sides are `R'`-linear, so the final identification with
  %      `(base_change_mate_regroupEquiv ψ φ M).inv` is `regroupEquiv.inv` checked on
  %      generators (`regroupEquiv` is fully proven, lem:base_change_mate_regroupEquiv).
  % NOTE (mechanism): the conjugate-counit factor is itself a `Scheme.Modules.conjugateEquiv`
  % component, so the move of `Θ_src`/regroup/`Θ_tgt` past the geometric factors is done on the
  % conjugate side — never as a positional rewrite under the X.Modules instance diamond. The master
  % counit-transport identity `huce` is landed (CategoryTheory.conjugateEquiv_counit_symm of the
  % comparison, with hpullinv/hcounitL/hcounitR fused). The inner value (step (a)) is REPROVEN INLINE
  % from the PROVED standalone `base_change_mate_fstar_reindex_legs_unitExpand` +
  % `…_gammaDistribute` + `gammaMap_pushforwardComp_*` + Seam-1 `base_change_mate_unit_value` +
  % `pullbackPushforward_unit_comp` — NOT cited from the sorry-backed
  % `base_change_mate_fstar_reindex_legs` / `base_change_mate_inner_value_eq`.
  This is the crux of the affine section-level computation: it feeds the
  section identity Lemma~\ref{lem:base_change_mate_section_identity} directly, the section
  identity being its immediate consequence via the counit factorization of the base-change map.
  By the counit factorization of the \((g^* \dashv g_*)\)-adjunction
  (\(g = \operatorname{Spec}\psi\)), the base-change map
  (Definition~\ref{def:pushforward_base_change_map}) factors as \(g^*(\theta_{\mathrm{in}})\)
  followed by the counit, with no opaque adjunction transpose remaining. On the global sections
  over \(\operatorname{Spec} R'\), conjugated by the identification \(\Theta_{\mathrm{src}}\)
  of Lemma~\ref{lem:base_change_mate_domain_read} and the identification \(\Theta_{\mathrm{tgt}}\)
  of Lemma~\ref{lem:base_change_mate_codomain_read}, this transpose is the extension of scalars
  along \(\psi : R \to R'\) of the \(\operatorname{Spec} R\)-section reading \(\rho\) of the inner
  pushforward comparison \(\theta_{\mathrm{in}}\), namely the canonical \(R\)-linear map
  \(\rho : m \mapsto (1 \otimes 1) \otimes m\) — the affine unit value of
  Lemma~\ref{lem:base_change_mate_unit_value}; on the generator \(r' \otimes m\) it returns
  \((1 \otimes r') \otimes m\), which is exactly the inverse of the regrouping isomorphism
  (Lemma~\ref{lem:base_change_mate_regroupEquiv}).
\end{lemma}
\begin{proof}
  \uses{lem:base_change_mate_fstar_reindex_legs_unitExpand,
        lem:base_change_mate_fstar_reindex_legs_gammaDistribute,
        lem:base_change_mate_unit_value, lem:pullbackPushforward_unit_comp,
        lem:base_change_mate_regroupEquiv, lem:conjugateEquiv_counit_symm_mathlib,
        lem:pullback_spec_tilde_iso, lem:gammaMap_pushforwardComp_hom_eq_id,
        lem:gammaMap_pushforwardComp_inv_eq_id}
  This is the closing step of the affine section-level computation; the section
  identity Lemma~\ref{lem:base_change_mate_section_identity} follows from it directly.

  \emph{Counit factorization and the master transport identity.} The base-change map
  (Definition~\ref{def:pushforward_base_change_map}) is the \((g^* \dashv g_*)\)-transpose of
  \(\theta_{\mathrm{in}}\) (\(g = \operatorname{Spec}\psi\)); by the counit form of the hom-set
  adjunction it factors as \(g^*(\theta_{\mathrm{in}})\) followed by the geometric counit
  \(\varepsilon_g\), with no opaque transpose remaining,
  \(\Gamma(\theta) = \Gamma_{R'}\bigl(g^*(\theta_{\mathrm{in}})\bigr) \circ \Gamma_{R'}(\varepsilon_g)\).
  Instantiating the transposed counit-across-conjugate coherence
  Lemma~\ref{lem:conjugateEquiv_counit_symm_mathlib} at the two composite adjunctions
  \(\mathrm{adj}_L = (\widetilde{(-)}_R \dashv \Gamma_R) . (g^* \dashv g_*)\) and
  \(\mathrm{adj}_R = (\operatorname{ext}_\psi \dashv \operatorname{restr}_\psi) .
  (\widetilde{(-)}_{R'} \dashv \Gamma_{R'})\) with \(\alpha = \beta^{\mathrm{hom}}\) the pushforward
  \(\Gamma\)-comparison, and fusing the two composite-counit splits of \(\mathrm{adj}_L\) and
  \(\mathrm{adj}_R\), yields the \emph{master counit-transport identity}
  \[
    \widetilde{(-)}_{R'}\!\bigl(\varepsilon^{\mathrm{alg}}\bigr) \circ \varepsilon^{\Gamma}
    \;=\;
    \Pi_{\psi}^{-1} \circ g^*\!\bigl(\varepsilon^{\Gamma}\bigr) \circ \varepsilon_g ,
  \]
  where \(\varepsilon^{\mathrm{alg}}\) is the algebraic extension/restriction-of-scalars counit
  along \(\psi\), \(\varepsilon^{\Gamma}\) the tilde--\(\Gamma\) counit, and \(\Pi_{\psi}\) the
  pullback dictionary (Lemma~\ref{lem:pullback_spec_tilde_iso}); the dictionary enters because the
  conjugate of \(\beta^{\mathrm{hom}}\) is \(\Pi_{\psi}^{-1}\). The trailing dictionary and
  tilde-counit factors are isomorphisms, so solving this identity for \(\varepsilon_g\) and applying
  \(\Gamma_{R'}\) rewrites \(\Gamma_{R'}(\varepsilon_g)\) in the goal into the algebraic counit
  conjugated by exactly the dictionaries that \(\Theta_{\mathrm{src}}\)
  (Lemma~\ref{lem:base_change_mate_domain_read}) and \(\Theta_{\mathrm{tgt}}\)
  (Lemma~\ref{lem:base_change_mate_codomain_read}) bake in. The remaining \(\sim\)150 lines of
  telescoping are three steps.

  \emph{(a) Inner reindex} (\(\Gamma_R(\theta_{\mathrm{in}}) = \rho\)). The \(\operatorname{Spec} R\)
  section reading of the inner comparison \(\theta_{\mathrm{in}}\) is the canonical \(R\)-linear map
  \(\rho : m \mapsto (1 \otimes 1) \otimes m\) (the affine inner value). This is reproven
  \emph{inline}, not cited from the leg-reindex wrapper: expand the buried \((g')\)-unit by the
  leg-unit expansion Lemma~\ref{lem:base_change_mate_fstar_reindex_legs_unitExpand} (whose engine is
  the proved unit-distribution Lemma~\ref{lem:pullbackPushforward_unit_comp}), distribute the
  resulting four factors through \((\operatorname{Spec}\varphi)_*\) and \(\Gamma\) by
  Lemma~\ref{lem:base_change_mate_fstar_reindex_legs_gammaDistribute}, collapse the
  pushforward-composition coherences to the identity \(\Gamma\)-image
  (Lemmas~\ref{lem:gammaMap_pushforwardComp_hom_eq_id},
  \ref{lem:gammaMap_pushforwardComp_inv_eq_id}), and evaluate the lone surviving affine unit by
  Seam~1 (Lemma~\ref{lem:base_change_mate_unit_value}); the result is \(\rho\). With the inner value
  available the residual \(g^*\)-leg, conjugated by \(\Pi_{\psi}\), is
  \(\operatorname{ext}_\psi(\rho)\).

  \emph{(b) One-generator close.} The composite
  \(\operatorname{ext}_\psi(\rho) \circ \varepsilon^{\mathrm{alg}}\) is an \(R'\)-linear map
  \(R' \otimes_R M \to (R' \otimes_R A) \otimes_A M\); by \(\operatorname{ModuleCat}\)
  extensionality it suffices to evaluate it on the generator \(r' \otimes m\), where it returns
  \(r' \cdot ((1 \otimes 1) \otimes m) = (1 \otimes r') \otimes m\). This is exactly the value of the
  inverse regrouping isomorphism \(\operatorname{regroup}^{-1}\) on \(r' \otimes m\), so a
  one-generator \(\operatorname{ext}\) against Lemma~\ref{lem:base_change_mate_regroupEquiv} gives
  \(\operatorname{ext}_\psi(\rho) \circ \varepsilon^{\mathrm{alg}} = \operatorname{regroup}^{-1}\).

  \emph{(c) Dictionary cancellation.} It remains to match the \(\Pi_{\psi}\) (pullback dictionary,
  Lemma~\ref{lem:pullback_spec_tilde_iso}) and tilde-counit factors appearing in the master
  transport identity against the \(\operatorname{pushforward}\)- and \(\operatorname{pullback}\)-tilde
  dictionaries baked into the domain read \(\Theta_{\mathrm{src}}\) and codomain read
  \(\Theta_{\mathrm{tgt}}\). These factors are mutually inverse isomorphisms, and cancelling them
  reduces the conjugated geometric transpose to the algebraic composite of step~(b), namely
  \(\operatorname{regroup}^{-1}\). This is the claimed equation
  \(\Theta_{\mathrm{src}}^{-1} \ggg \Gamma(\theta) \ggg \Theta_{\mathrm{tgt}}
  = \operatorname{regroup}^{-1}\).
\end{proof}

% --- Section-level identity (now a one-line chain of the three seams above). The base-change
%     map on sections is the R'-base change of the algebraic unit of extension of scalars;
%     reading it through the two tilde dictionaries pins it as the inverse of the regrouping
%     isomorphism. ---

\begin{lemma}

\leanok
[The section-level base-change map is the base change of the unit]
  \label{lem:base_change_mate_section_identity}
  \lean{AlgebraicGeometry.base_change_mate_section_identity}
  \uses{def:pushforward_base_change_map, lem:base_change_mate_domain_read,
        lem:base_change_mate_codomain_read, lem:base_change_mate_regroupEquiv}
  % SOURCE: [Stacks Project], Cohomology of Schemes, Lemma
  % `lemma-affine-base-change' (\S\,Cohomology and base change, I), proof, the
  % "boils down to the equality" step (read from references/stacks-coherent.tex,
  % L932--937).
  % SOURCE QUOTE: "Thus we see that the lemma boils down to the
  % equality
  % $$
  % (R' \otimes_R A) \otimes_A M = R' \otimes_R M
  % $$
  % as $R'$-modules."
  % SOURCE (cohomological form): [Stacks Project], Cohomology of Schemes, Tag 02KH
  % `lemma-flat-base-change-cohomology', part (2) (read from
  % references/stacks-coherent.tex, L967--968) — the H^0/direct-image-with-base-change
  % specialisation this affine identity underlies.
  % SOURCE QUOTE (cohomological form): "if $S = \Spec(A)$ and $S' = \Spec(B)$, then
  % $H^i(X, \mathcal{F}) \otimes_A B = H^i(X', \mathcal{F}')$."
  \textit{Source: Stacks Project, Cohomology of Schemes, Lemma ``Affine base
  change'' (proof, ``boils down to the equality'' step); cohomological form Tag 02KH(2).}
  % LEAN SIGNATURE (the scaffold creates the decl pinned above with this signature).
  %   The conjugate of `Γ(θ)` through the domain read `Θ_src` and codomain read `Θ_tgt`
  %   is the `R'`-base change of the algebraic unit `η_M : M → (A ⊗_R R') ⊗_A M`,
  %   `m ↦ (1 ⊗ₜ 1) ⊗ₜ m` (`LinearMap.lTensor R' η_M` post-composed with the `R'`-action),
  %   which on the generator `r' ⊗ m` returns `r' • ((1 ⊗ 1) ⊗ m) = (1 ⊗ r') ⊗ m`; this is
  %   exactly `(base_change_mate_regroupEquiv ψ φ M).inv`. The formalized conclusion is the
  %   propositional equality of `ModuleCat R'` morphisms:
  %     {R R' A : CommRingCat.{u}} (ψ : R ⟶ R') (φ : R ⟶ A) (M : ModuleCat.{u} A) :
  %       (base_change_mate_domain_read ψ φ M).inv ≫
  %         (moduleSpecΓFunctor (R := R')).map
  %           (pushforwardBaseChangeMap (Spec.map φ) (Spec.map ψ)
  %             (Limits.pullback.snd (Spec.map φ) (Spec.map ψ))
  %             (Limits.pullback.fst (Spec.map φ) (Spec.map ψ))
  %             (IsPullback.of_hasPullback (Spec.map φ) (Spec.map ψ)).w (tilde M)) ≫
  %         (base_change_mate_codomain_read ψ φ M).hom
  %         = (base_change_mate_regroupEquiv ψ φ M).inv
  In the affine--affine model (\(g = \operatorname{Spec}\psi\),
  \(f = \operatorname{Spec}\varphi\), \(\mathcal{F} = \widetilde M\) for an \(A\)-module
  \(M\)), write \(\Gamma(\theta)\) for the global-sections incarnation of the base-change
  map (Definition~\ref{def:pushforward_base_change_map}). Let
  \(\Theta_{\mathrm{src}} : \Gamma\bigl(g^*(f_*\widetilde M)\bigr) \xrightarrow{\sim}
  R' \otimes_R M\) be the identification of Lemma~\ref{lem:base_change_mate_domain_read}
  and \(\Theta_{\mathrm{tgt}} : \Gamma\bigl(f'_*(g')^*\widetilde M\bigr) \xrightarrow{\sim}
  (A \otimes_R R') \otimes_A M\) the identification of
  Lemma~\ref{lem:base_change_mate_codomain_read}. Then the conjugated map
  \[
    \Theta_{\mathrm{tgt}} \circ \Gamma(\theta) \circ \Theta_{\mathrm{src}}^{-1} :
    R' \otimes_R M \longrightarrow (A \otimes_R R') \otimes_A M
  \]
  is the \(R'\)-base change of the algebraic unit \(\eta_M : M \to (A \otimes_R R') \otimes_A M\),
  \(m \mapsto (1 \otimes 1) \otimes m\): explicitly, the unique \(R'\)-linear map sending
  \(r' \otimes m \mapsto r' \cdot \eta_M(m) = (1 \otimes r') \otimes m\). Equivalently,
  this conjugated map is the \emph{inverse} of the regrouping isomorphism
  Lemma~\ref{lem:base_change_mate_regroupEquiv},
  \(\Theta_{\mathrm{tgt}} \circ \Gamma(\theta) \circ \Theta_{\mathrm{src}}^{-1}
  = \operatorname{regroup}^{-1}\), and in particular an isomorphism with \emph{no}
  flatness hypothesis.
\end{lemma}

\begin{proof}
  \uses{lem:base_change_mate_domain_read, lem:base_change_mate_codomain_read,
        lem:base_change_mate_gstar_transpose}
  The section identity is an immediate consequence of the \(g^*\)-transpose crux
  Lemma~\ref{lem:base_change_mate_gstar_transpose}, with the domain and codomain pinned by
  the domain read Lemma~\ref{lem:base_change_mate_domain_read} (\(\Theta_{\mathrm{src}}\)) and
  the codomain read Lemma~\ref{lem:base_change_mate_codomain_read} (\(\Theta_{\mathrm{tgt}}\)).
  By definition the base-change map (Definition~\ref{def:pushforward_base_change_map}) is the
  \((g^* \dashv g_*)\)-adjunction transpose of the inner pushforward comparison
  \(\theta_{\mathrm{in}}\) (with \(g = \operatorname{Spec}\psi\)); by the counit form of the
  hom-set adjunction it factors as \(g^*(\theta_{\mathrm{in}})\) followed by the counit
  \(\varepsilon_g\), with no opaque transpose remaining,
  \(\Gamma(\theta) = \Gamma\bigl(g^*(\theta_{\mathrm{in}})\bigr) \circ \Gamma(\varepsilon_g)\).
  Substituting this factorization, the conjugated map
  \[
    \Theta_{\mathrm{src}}^{-1} \circ \Gamma(\theta) \circ \Theta_{\mathrm{tgt}}
      = \Theta_{\mathrm{src}}^{-1} \circ
        \Gamma\bigl(g^*(\theta_{\mathrm{in}})\bigr) \circ \Gamma(\varepsilon_g) \circ
        \Theta_{\mathrm{tgt}}
  \]
  is exactly the left-hand side of the section-level equation of
  Lemma~\ref{lem:base_change_mate_gstar_transpose}, which states it equals
  \(\operatorname{regroup}^{-1}\). On the generator \(r' \otimes m\) the conjugated map
  therefore returns \((1 \otimes r') \otimes m\), the value of the inverse regrouping
  isomorphism Lemma~\ref{lem:base_change_mate_regroupEquiv}; both maps are \(R'\)-linear and
  agree on generators \(1 \otimes m\), so they coincide. This is the ``boils down to the
  equality'' \((A \otimes_R R') \otimes_A M = R' \otimes_R M\) of the source, with no
  flatness hypothesis used anywhere.
\end{proof}

\begin{lemma}


\leanok
[Generator trace of the section-level base-change map]
  \label{lem:base_change_mate_generator_trace}
  \lean{AlgebraicGeometry.base_change_mate_generator_trace}
  % NOTE: the Lean decl formalizes the `IsIso (Θ_src⁻¹ ≫ Γ(α) ≫ Θ_tgt)`
  % *corollary* of the section identity, not the literal generator formula
  % `r' ⊗ m ↦ (1 ⊗ r') ⊗ m`. The IsIso form is what the parent
  % lem:pushforward_base_change_mate_cancelBaseChange consumes; it is faithful and
  % non-vacuous. By the section identity (lem:base_change_mate_section_identity) the
  % conjugate equals a LinearEquiv (regroupEquiv.inv), so this is a one-line
  % `rw` + `infer_instance`.
  \uses{lem:base_change_mate_regroupEquiv, lem:base_change_mate_section_identity}
  With the domain and codomain pinned by Lemma~\ref{lem:base_change_mate_domain_read}
  and Lemma~\ref{lem:base_change_mate_codomain_read}, the conjugated section-level
  base-change map
  \[
    \Theta_{\mathrm{src}}^{-1} \;\ggg\; \Gamma(\alpha) \;\ggg\; \Theta_{\mathrm{tgt}} :
    R' \otimes_R M \longrightarrow (R' \otimes_R A) \otimes_A M
  \]
  is an isomorphism: \(\mathrm{IsIso}\bigl(\Theta_{\mathrm{src}}^{-1} \ggg \Gamma(\alpha)
  \ggg \Theta_{\mathrm{tgt}}\bigr)\).
\end{lemma}

\begin{proof}
  \uses{lem:base_change_mate_regroupEquiv, lem:base_change_mate_section_identity}
  By the section identity Lemma~\ref{lem:base_change_mate_section_identity}
  the conjugated map equals \(\operatorname{regroup}^{-1}\), the inverse of the bundled
  \(R'\)-linear isomorphism of Lemma~\ref{lem:base_change_mate_regroupEquiv}. The inverse
  of a linear equivalence is itself an isomorphism, so
  \(\mathrm{IsIso}\bigl(\Theta_{\mathrm{src}}^{-1} \ggg \Gamma(\alpha)
  \ggg \Theta_{\mathrm{tgt}}\bigr)\) follows since \(\operatorname{regroup}^{-1}\)
  is a linear equivalence and hence an isomorphism.
\end{proof}

\subsection{Tilde-transport direct proof of the section-level value}
\label{sec:fbc_tilde_transport_direct}

% This subsection gives a direct affine computation of the section-level value of
% the base-change map. It routes ENTIRELY through the single-step tilde dictionaries
% (Lemma~\ref{lem:pullback_spec_tilde_iso}, Lemma~\ref{lem:pushforward_spec_tilde_iso})
% and the tensor-product universal property; it does NOT depend on
% lem:base_change_mate_section_identity or lem:base_change_mate_gstar_transpose
% (the conjugate-calculus crux). Those conjugate scaffolding blocks remain in the
% chapter as a dictionary/record but are no longer on the critical path.

% NOTE (review iter-044): the "direct" bypass claimed above is ILLUSORY. iter-043 verified that
% the concrete affine square's inner unit is the `g'=pullback.fst` unit (no element normal form),
% so Gamma(alpha) must transit the same tilde dictionaries as the conjugate intertwining — this
% lemma collapses onto the SAME open keystone lem:base_change_mate_fstar_reindex_legs_conj. The
% Lean target `pushforward_base_change_mate_sections_direct` was NOT added (correctly). Treat this
% block as an abandoned route record, not a live critical path; do NOT dispatch a prover at it.
\begin{lemma}
[Direct section-level value of the affine base-change map]
  \label{lem:pushforward_base_change_mate_sections_direct}
  \lean{AlgebraicGeometry.pushforward_base_change_mate_sections_direct}
  \uses{def:pushforward_base_change_map,
        lem:base_change_mate_domain_read,
        lem:base_change_mate_codomain_read,
        lem:pullback_spec_tilde_iso,
        lem:pushforward_spec_tilde_iso,
        lem:cancelBaseChange_mathlib}
  In the affine--affine model — \(S = \operatorname{Spec}(R)\),
  \(S' = \operatorname{Spec}(R')\), \(X = \operatorname{Spec}(A)\),
  \(\mathcal{F} = \widetilde M\) for an \(A\)-module \(M\), \(g = \operatorname{Spec}\psi\)
  for \(\psi : R \to R'\), \(f = \operatorname{Spec}\varphi\) for \(\varphi : R \to A\) —
  the composition
  \[
    \Theta_{\mathrm{src}}^{-1} \;\ggg\; \Gamma(\alpha) \;\ggg\; \Theta_{\mathrm{tgt}}
    \;:\; R' \otimes_R M
    \;\longrightarrow\; (R' \otimes_R A) \otimes_A M
  \]
  equals the inverse cancellation isomorphism
  \(\operatorname{cancelBaseChange}^{-1} : r' \otimes m \mapsto (r' \otimes 1) \otimes m\)
  of Lemma~\ref{lem:cancelBaseChange_mathlib}, where \(\alpha\) is the base-change map
  (Definition~\ref{def:pushforward_base_change_map}) and \(\Theta_{\mathrm{src}}\),
  \(\Theta_{\mathrm{tgt}}\) are the domain and codomain identifications of
  Lemmas~\ref{lem:base_change_mate_domain_read}
  and~\ref{lem:base_change_mate_codomain_read}.
\end{lemma}

% SOURCE: [Stacks Project], Cohomology of Schemes, Lemma
% `lemma-affine-base-change' (\S\,Cohomology and base change, I), proof, the
% "boils down to the equality" step (read from references/stacks-coherent.tex,
% L927--938).
% SOURCE QUOTE PROOF: "We use Schemes, Lemma \ref{schemes-lemma-widetilde-pullback}
% to describe pullbacks and pushforwards of $\mathcal{F}$.
% Namely, $X' = \Spec(R' \otimes_R A)$ and
% $\mathcal{F}'$ is the quasi-coherent sheaf associated
% to $(R' \otimes_R A) \otimes_A M$.
% Thus we see that the lemma boils down to the
% equality
% $$
% (R' \otimes_R A) \otimes_A M = R' \otimes_R M
% $$
% as $R'$-modules."
\begin{proof}
  \uses{def:pushforward_base_change_map,
        lem:base_change_mate_domain_read,
        lem:base_change_mate_codomain_read,
        lem:pullback_spec_tilde_iso,
        lem:pushforward_spec_tilde_iso,
        lem:cancelBaseChange_mathlib}
  \textit{Source: Stacks Project, Cohomology of Schemes, Lemma ``Affine base
  change'' (proof, ``boils down to the equality'' step).}

  The proof is a direct module computation: we unfold the \emph{definition} of
  \(\alpha = \mathtt{pushforwardBaseChangeMap}\) and evaluate each adjunction
  unit/counit through the single-step tilde dictionaries, never invoking the
  composite-mate recognition. By
  Definition~\ref{def:pushforward_base_change_map}, \(\alpha\) is the
  \((g^* \dashv g_*)\)-adjunct of the map
  \[
    \beta : f_*\widetilde M \longrightarrow g_*\,f'_*(g')^*\widetilde M
          = f_*\,(g')_*(g')^*\widetilde M
  \]
  obtained by applying \(f_*\) to the unit
  \(\eta : \widetilde M \to (g')_*(g')^*\widetilde M\) of the
  \(((g')^*, (g')_*)\)-adjunction, using the square's commutativity
  \(g \circ f' = f \circ g'\). Explicitly, as adjunct,
  \(\alpha = \varepsilon_{f'_*(g')^*\widetilde M} \circ g^*(\beta)\), where
  \(\varepsilon\) is the counit of \(g^* \dashv g_*\). We trace each factor on global
  sections over the affine model, on a generator \(r' \otimes m \in R' \otimes_R M\).

  \emph{Step 1 — the inner unit \(\eta\) is the extension-of-scalars unit.}
  On \(X = \operatorname{Spec}(A)\), the projection \(g'\) is the
  \(\operatorname{Spec}\)-map of the algebra inclusion \(A \to R' \otimes_R A\)
  (Lemma~\ref{lem:base_change_mate_codomain_read}). By the pullback dictionary
  Lemma~\ref{lem:pullback_spec_tilde_iso}, \((g')^*\widetilde M\) is the tilde of the
  extension of scalars \((R' \otimes_R A) \otimes_A M\), and that lemma identifies the
  adjunction unit with the algebraic unit of extension of scalars; hence on sections
  \(\eta\) is the \(A\)-linear map
  \[
    M \longrightarrow (R' \otimes_R A) \otimes_A M, \qquad
    m \longmapsto (1 \otimes 1) \otimes m .
  \]
  This is forced by the universal property of the tensor product and the algebra map
  \(A \to R' \otimes_R A\); no mate identity is used.

  \emph{Step 2 — apply \(f_* = \operatorname{restr}_\varphi\).} The pushforward
  dictionary Lemma~\ref{lem:pushforward_spec_tilde_iso} reads \(f_*\) of a tilde as
  restriction of scalars along \(\varphi\). Restriction of scalars does not move
  elements, so on sections over \(S = \operatorname{Spec}(R)\) the map
  \(\beta = f_*(\eta)\) (with codomain rewritten by \(g \circ f' = f \circ g'\)) is the
  same assignment \(m \mapsto (1 \otimes 1) \otimes m\), now read \(R\)-linearly into
  \((R' \otimes_R A) \otimes_A M\).

  \emph{Step 3 — apply \(g^* = R' \otimes_R (-)\).} The pullback dictionary
  Lemma~\ref{lem:pullback_spec_tilde_iso} for \(g = \operatorname{Spec}\psi\) reads
  \(g^*\) of a tilde as extension of scalars \(R' \otimes_R (-)\) along \(\psi\), and is
  functorial: \(g^*(\beta)\) on sections over \(S' = \operatorname{Spec}(R')\) is
  \(\operatorname{id}_{R'} \otimes \beta\),
  \[
    R' \otimes_R M \longrightarrow R' \otimes_R \bigl((R' \otimes_R A) \otimes_A M\bigr),
    \qquad
    r' \otimes m \longmapsto r' \otimes \bigl((1 \otimes 1) \otimes m\bigr).
  \]
  The domain \(R' \otimes_R M\) here is exactly \(\Theta_{\mathrm{src}}\)'s reading of
  \(\Gamma(g^*(f_*\widetilde M))\) (Lemma~\ref{lem:base_change_mate_domain_read}).

  \emph{Step 4 — the counit \(\varepsilon\) is the extension-of-scalars counit.}
  Writing \(W = (R' \otimes_R A) \otimes_A M\) for the target \(R'\)-module, the counit
  \(\varepsilon : g^* g_* \widetilde W \to \widetilde W\) of \(g^* \dashv g_*\) is, on
  sections over \(\operatorname{Spec}(R')\) and again by
  Lemma~\ref{lem:pullback_spec_tilde_iso} (which identifies the counit with the
  algebraic action map of extension of scalars), the \(R'\)-linear map
  \[
    R' \otimes_R \bigl(\operatorname{restr}_\psi W\bigr) \longrightarrow W,
    \qquad r' \otimes w \longmapsto r' \cdot w ,
  \]
  where \(r' \cdot w\) is the \(R'\)-action on \(W\). The codomain \(W\) is exactly
  \(\Theta_{\mathrm{tgt}}\)'s reading of \(\Gamma(f'_*(g')^*\widetilde M)\)
  (Lemma~\ref{lem:base_change_mate_codomain_read}).

  \emph{Composition.} Composing Steps 3 and 4, the conjugated section map sends
  \[
    r' \otimes m
    \;\overset{g^*(\beta)}{\longmapsto}\;
    r' \otimes \bigl((1 \otimes 1) \otimes m\bigr)
    \;\overset{\varepsilon}{\longmapsto}\;
    r' \cdot \bigl((1 \otimes 1) \otimes m\bigr)
    = (r' \otimes 1) \otimes m ,
  \]
  the last equality because the \(R'\)-action on \((R' \otimes_R A) \otimes_A M\)
  multiplies into the \(R'\) tensor slot: \(r' \cdot ((1 \otimes 1) \otimes m)
  = ((r' \cdot 1) \otimes 1) \otimes m = (r' \otimes 1) \otimes m\). Thus
  \(\Theta_{\mathrm{src}}^{-1} \ggg \Gamma(\alpha) \ggg \Theta_{\mathrm{tgt}}\) is the
  \(R'\)-linear map \(r' \otimes m \mapsto (r' \otimes 1) \otimes m\), which is
  precisely \(\operatorname{cancelBaseChange}^{-1}\) of
  Lemma~\ref{lem:cancelBaseChange_mathlib}
  (\(\mathtt{cancelBaseChange\_symm\_tmul}\)). No flatness hypothesis is used, and the
  value falls out of the tensor-product universal property together with the algebra
  map \(\varphi : R \to A\) — not from any conjugate/mate coherence identity.

  % LEAN HINT: prove by `ext` (or `TensorProduct.induction_on`) on a generator
  % `r' ⊗ m`, then unfold `pushforwardBaseChangeMap` to its adjunct form
  % `ε ∘ g^*(f_* η)` and rewrite each unit/counit by the section-level images of
  % `pullback_spec_tilde_iso` / `pushforward_spec_tilde_iso`. The closing step is
  % `cancelBaseChange_symm_tmul`. The two endpoints `Θ_src`, `Θ_tgt` are supplied by
  % `base_change_mate_domain_read` / `base_change_mate_codomain_read`. No flatness
  % hypothesis and no appeal to `base_change_mate_gstar_transpose`.
\end{proof}

\begin{lemma}

\leanok
[Section-level value of the base-change map equals the cancellation isomorphism]
  \label{lem:pushforward_base_change_mate_cancelBaseChange}
  \lean{AlgebraicGeometry.pushforward_base_change_mate_cancelBaseChange}
  % NOTE: the Lean decl formalizes the `IsIso (Γ(α))` *corollary* of
  % this lemma, not the literal equality `Θ_tgt ∘ Γ(α) ∘ Θ_src⁻¹ = cancelBaseChange⁻¹`.
  % The IsIso form is faithful and non-vacuous (it is exactly what the affine close
  % consumes via lem:base_change_map_affine_local); stating the equality requires first
  % identifying pullback.fst/snd over the generic square with the Spec-of-tensor
  % inclusions (pullbackSpecIso), which is the same plumbing that blocks the proof.
  % NOTE: the DIRECT proof body of this decl assembles the IsIso conclusion via the
  % tilde-transport route: the domain read Θ_src and codomain read Θ_tgt identify Γ(α)'s
  % domain/codomain with `R' ⊗_R M` and `(R' ⊗_R A) ⊗_A M`, and the direct section-value
  % lemma lem:pushforward_base_change_mate_sections_direct shows the conjugate
  % `Θ_src⁻¹ ≫ Γ(α) ≫ Θ_tgt` equals `cancelBaseChange⁻¹` (no flatness), hence is an
  % isomorphism; conjugating back, Γ(α) is itself an isomorphism. This route does NOT
  % depend on lem:base_change_mate_section_identity or lem:base_change_mate_gstar_transpose
  % (the conjugate-counit crux), which remain in the chapter only as a dead-end record.
  \uses{def:pushforward_base_change_map, lem:base_change_mate_domain_read,
  lem:base_change_mate_codomain_read,
  lem:pushforward_base_change_mate_sections_direct,
  lem:cancelBaseChange_mathlib}
  % SOURCE: [Stacks Project], Cohomology of Schemes, Lemma
  % `lemma-affine-base-change' (\S\,Cohomology and base change, I), proof, the
  % "boils down to the equality" step (read from references/stacks-coherent.tex,
  % L927--938).
  % SOURCE QUOTE: "We use Schemes, Lemma \ref{schemes-lemma-widetilde-pullback}
  % to describe pullbacks and pushforwards of $\mathcal{F}$.
  % Namely, $X' = \Spec(R' \otimes_R A)$ and
  % $\mathcal{F}'$ is the quasi-coherent sheaf associated
  % to $(R' \otimes_R A) \otimes_A M$.
  % Thus we see that the lemma boils down to the
  % equality
  % $$
  % (R' \otimes_R A) \otimes_A M = R' \otimes_R M
  % $$
  % as $R'$-modules."
  \textit{Source: Stacks Project, Cohomology of Schemes, Lemma ``Affine base
  change'' (proof, ``boils down to the equality'' step).}
  In the affine--affine case \(S = \operatorname{Spec}(R)\),
  \(S' = \operatorname{Spec}(R')\), \(X = \operatorname{Spec}(A)\),
  \(\mathcal{F} = \widetilde{M}\) — with \(g = \operatorname{Spec}(\psi)\) for
  \(\psi : R \to R'\), \(f = \operatorname{Spec}(\varphi)\) for \(\varphi : R \to A\),
  and \(M\) an \(A\)-module — write \(\Gamma(\alpha)\) for the global-sections
  incarnation of the base-change map (Definition~\ref{def:pushforward_base_change_map}).
  Recall \(\alpha\) runs
  \(g^*(f_*\widetilde{M}) \to f'_*(g')^*\widetilde{M}\). Transport \(\Gamma(\alpha)\)
  through the two proved affine dictionaries that identify its domain and codomain: the
  pushforward dictionary Lemma~\ref{lem:pushforward_spec_tilde_iso}, followed by the
  pullback dictionary Lemma~\ref{lem:pullback_spec_tilde_iso}, identifies the domain
  \(\Gamma\bigl(g^*(f_*\widetilde{M})\bigr)\) with \(R' \otimes_R M\); the pullback
  dictionary Lemma~\ref{lem:pullback_spec_tilde_iso}, followed by the pushforward
  dictionary Lemma~\ref{lem:pushforward_spec_tilde_iso}, identifies the codomain
  \(\Gamma\bigl(f'_*(g')^*\widetilde{M}\bigr)\) with \((R' \otimes_R A) \otimes_A M\),
  both read as \(R'\)-modules. Under these identifications \(\Gamma(\alpha)\) is
  precisely the canonical cancellation isomorphism
  \[
    \operatorname{cancelBaseChange} :
    (R' \otimes_R A) \otimes_A M \;\xrightarrow{\ \sim\ }\; R' \otimes_R M,
    \qquad (r' \otimes a) \otimes m \;\mapsto\; r' \otimes (a \cdot m)
  \]
  (Lemma~\ref{lem:cancelBaseChange_mathlib}; here \(a \cdot m\) is the \(A\)-action on \(M\)). Equivalently, in the direction in which
  \(\Gamma(\alpha)\) is presented — from the source \(R' \otimes_R M\) to the target
  \((R' \otimes_R A) \otimes_A M\) — it is \(\operatorname{cancelBaseChange}^{-1}\),
  \(r' \otimes m \mapsto (r' \otimes 1) \otimes m\); the two readings agree because
  \(\operatorname{cancelBaseChange}\) is an isomorphism. Since
  \(\operatorname{cancelBaseChange}\) is an isomorphism with \emph{no} flatness
  hypothesis, this identification closes the affine case: \(\Gamma(\alpha)\) is an
  isomorphism, hence (full faithfulness / counit isomorphism of \(\widetilde{(-)}\) on
  quasi-coherent modules) so is \(\alpha\) on the affine chart.

\end{lemma}

\begin{proof}
  \uses{def:pushforward_base_change_map, lem:base_change_mate_domain_read,
        lem:base_change_mate_codomain_read,
        lem:pushforward_base_change_mate_sections_direct,
        lem:cancelBaseChange_mathlib}
  By Lemma~\ref{lem:base_change_mate_domain_read} the
  source \(\Gamma\bigl(g^*(f_*\widetilde M)\bigr)\) is \(R' \otimes_R M\) via
  \(\Theta_{\mathrm{src}}\), and by Lemma~\ref{lem:base_change_mate_codomain_read} the
  target \(\Gamma\bigl(f'_*(g')^*\widetilde M\bigr)\) is \((R' \otimes_R A) \otimes_A M\)
  via \(\Theta_{\mathrm{tgt}}\). By the direct section-value computation
  Lemma~\ref{lem:pushforward_base_change_mate_sections_direct}, the conjugated map
  \[
    \theta = \Theta_{\mathrm{src}}^{-1} \ggg \Gamma(\alpha) \ggg \Theta_{\mathrm{tgt}}
  \]
  equals \(\operatorname{cancelBaseChange}^{-1}\), the inverse of the cancellation
  isomorphism of Lemma~\ref{lem:cancelBaseChange_mathlib}. Since
  \(\operatorname{cancelBaseChange}\) is an isomorphism with \emph{no} flatness
  hypothesis, so is its inverse \(\theta\). As \(\theta\) is the conjugate of
  \(\Gamma(\alpha)\) by the two dictionary isomorphisms
  \(\Theta_{\mathrm{src}}, \Theta_{\mathrm{tgt}}\), it follows that \(\Gamma(\alpha)\) is
  itself an isomorphism — assembled by writing
  \(\Gamma(\alpha) = \Theta_{\mathrm{src}} \ggg \theta \ggg \Theta_{\mathrm{tgt}}^{-1}\)
  and inferring the instance. Both the source and target of \(\alpha\) are
  quasi-coherent, so the tilde-equivalence \(\widetilde{(-)}\) is fully faithful on
  them and \(\mathrm{IsIso}\,\Gamma(\alpha)\) yields \(\mathrm{IsIso}\,\alpha\); the
  closing conclusion the affine close consumes is \(\mathrm{IsIso}\ \Gamma(\alpha)\).
  This routes entirely through the tilde dictionaries — it does not use the
  conjugate-calculus crux Lemma~\ref{lem:base_change_mate_gstar_transpose} or the
  section identity Lemma~\ref{lem:base_change_mate_section_identity} — and aligns the
  section-level map with the Stacks ``boils down to the equality'' step.
\end{proof}

\begin{lemma}\leanok
[Affine base change]
  \label{lem:affine_base_change_pushforward}
  \lean{AlgebraicGeometry.affineBaseChange_pushforward_iso}
  \uses{def:pushforward_base_change_map, lem:base_change_map_affine_local,
        lem:modules_isIso_iff_affineOpens, lem:pullback_spec_tilde_iso,
        lem:pushforward_base_change_mate_cancelBaseChange,
        lem:pushforward_spec_tilde_iso}
  % SOURCE: [Stacks Project], Cohomology of Schemes, Lemma
  % `lemma-affine-base-change' (\S\,Cohomology and base change, I)
  % (read from references/stacks-coherent.tex, L906--918).
  % SOURCE QUOTE: "Let $f : X \to S$ be a morphism of schemes.
  % Let $\mathcal{F}$ be a quasi-coherent $\mathcal{O}_X$-module.
  % Assume $f$ is affine.
  % In this case $f_*\mathcal{F} \cong Rf_*\mathcal{F}$ is
  % a quasi-coherent sheaf, and for every base change diagram
  % (\ref{equation-base-change-diagram})
  % we have
  % $$
  % g^*f_*\mathcal{F} = f'_*(g')^*\mathcal{F}.
  % $$"
  \textit{Source: Stacks Project, Cohomology of Schemes, Lemma ``Affine base
  change''.}
  Assume \(f : X \to S\) is an affine morphism and \(\mathcal{F}\) is a
  quasi-coherent \(\mathcal{O}_X\)-module. Then for every base-change square as
  above the canonical map of
  Definition~\ref{def:pushforward_base_change_map} is an isomorphism
  \[
    g^*\bigl(f_*\mathcal{F}\bigr) \;\cong\; f'_*\bigl((g')^*\mathcal{F}\bigr).
  \]
  The quasi-coherence hypothesis on \(\mathcal{F}\) is essential and is not a mere
  convenience: the proof works over an affine open by writing
  \(\mathcal{F} = \widetilde{M}\) for a module \(M\), and a general
  \(\mathcal{O}_X\)-module need not be of this form
  \(\widetilde{M}\); without it neither the affine description of the pushforward
  nor the tensor-associativity computation below is available, and the conclusion
  fails.
\end{lemma}

% SOURCE QUOTE PROOF: "The vanishing of higher direct images is
% Lemma \ref{lemma-relative-affine-vanishing}.
% The statement is local on $S$ and $S'$. Hence we may
% assume $X = \Spec(A)$, $S = \Spec(R)$,
% $S' = \Spec(R')$ and $\mathcal{F} = \widetilde{M}$
% for some $A$-module $M$.
% We use Schemes, Lemma \ref{schemes-lemma-widetilde-pullback}
% to describe pullbacks and pushforwards of $\mathcal{F}$.
% Namely, $X' = \Spec(R' \otimes_R A)$ and
% $\mathcal{F}'$ is the quasi-coherent sheaf associated
% to $(R' \otimes_R A) \otimes_A M$.
% Thus we see that the lemma boils down to the
% equality
% $$
% (R' \otimes_R A) \otimes_A M = R' \otimes_R M
% $$
% as $R'$-modules."
% NOTE: This affine lemma is closed directly on global sections, via the two named
% obligations: lem:base_change_map_affine_local (locality reduction on the affine
% opens of S') and lem:pushforward_base_change_mate_cancelBaseChange (the section-
% level value of the base-change map is cancelBaseChange). Both proved affine tilde
% dictionaries -- lem:pushforward_spec_tilde_iso and lem:pullback_spec_tilde_iso --
% feed the second obligation. cancelBaseChange is in Mathlib with no flatness; the
% affine case carries no flatness hypothesis. No Čech infrastructure is used here.
\begin{proof}
  \uses{def:pushforward_base_change_map, lem:modules_isIso_iff_affineOpens, lem:pushforward_spec_tilde_iso, lem:pullback_spec_tilde_iso, lem:base_change_map_affine_local, lem:pushforward_base_change_mate_cancelBaseChange}
  The proof is the composite of the two named obligations, taken \emph{directly on
  global sections}.

  \medskip
  \emph{First reduction (locality on \(S'\)).} By
  Lemma~\ref{lem:modules_isIso_iff_affineOpens} it suffices to prove that the
  base-change map of Definition~\ref{def:pushforward_base_change_map} is an
  isomorphism on sections over every affine open of \(S'\). This is exactly the
  locality reduction Lemma~\ref{lem:base_change_map_affine_local}: restricting the
  cartesian square to an affine open \(U = \operatorname{Spec}(R') \subseteq S'\) and
  a chosen affine \(\operatorname{Spec}(R) \subseteq S\) containing \(g(U)\), the
  hypothesis \([\mathrm{IsAffineHom}\ f]\) makes \(X\) restrict to
  \(\operatorname{Spec}(A)\) and quasi-coherence of \(\mathcal{F}\) makes it restrict
  to \(\widetilde{M}\) for an \(A\)-module \(M\). (It is here that quasi-coherence is
  used: it is what licenses writing \(\mathcal{F} = \widetilde{M}\).) Computing the
  fibre product gives \(X' = \operatorname{Spec}(R' \otimes_R A)\) and
  \(\mathcal{F}' = (g')^*\mathcal{F}\) the tilde of the \((R' \otimes_R A)\)-module
  \((R' \otimes_R A) \otimes_A M\). The matter is thus reduced to the affine--affine
  model, where the base-change map is a morphism \(\alpha\) of quasi-coherent
  \((\operatorname{Spec} R')\)-modules.

  \medskip
  \emph{Passage to global sections.} The functor
  \(\widetilde{(-)} : \operatorname{ModuleCat}(R') \to
  (\operatorname{Spec} R').\mathrm{Modules}\) is fully faithful with essential image
  the quasi-coherent modules; equivalently the counit
  \(\widetilde{\Gamma(\mathcal{N})} \to \mathcal{N}\) is an isomorphism exactly when
  \(\mathcal{N}\) is quasi-coherent. Source and target of \(\alpha\) are
  quasi-coherent, so both counits are isomorphisms and, by their naturality,
  \[
    \mathrm{IsIso}\,\alpha
    \;\Longleftrightarrow\;
    \mathrm{IsIso}\bigl(\Gamma(\alpha)\bigr),
  \]
  where \(\Gamma(\alpha)\) is a concrete \(R'\)-linear map in
  \(\operatorname{ModuleCat}(R')\). The whole computation therefore stays at the level
  of modules.

  \medskip
  \emph{Section-level identification (the second obligation).} By the section-level
  computation Lemma~\ref{lem:pushforward_base_change_mate_cancelBaseChange}, the map
  \(\Gamma(\alpha)\), with its domain \(\Gamma(g^*(f_*\widetilde{M}))\) identified with
  \(R' \otimes_R M\) and its codomain \(\Gamma(f'_*(g')^*\widetilde{M})\) identified
  with \((R' \otimes_R A) \otimes_A M\) through the two dictionaries
  Lemmas~\ref{lem:pullback_spec_tilde_iso} and~\ref{lem:pushforward_spec_tilde_iso},
  \emph{is} the canonical cancellation isomorphism
  \[
    \operatorname{cancelBaseChange} :
    (R' \otimes_R A) \otimes_A M \;\xrightarrow{\ \sim\ }\; R' \otimes_R M,
    \qquad (r' \otimes a) \otimes m \;\mapsto\; r' \otimes (a \cdot m)
  \]
  of \(R'\)-modules (Lemma~\ref{lem:cancelBaseChange_mathlib}). This is the entire
  content of the affine close, and it is performed \emph{directly on the module of
  global sections} — there is no separate abstract adjoint-mate identification at the
  level of sheaves. Because \(\operatorname{cancelBaseChange}\) is an isomorphism with
  no flatness hypothesis whatsoever, \(\Gamma(\alpha)\) is an isomorphism, hence so is
  \(\alpha\) on the affine chart, and by the locality reduction so is the global
  base-change map. Flatness of \(g\) plays no role in the affine lemma; it enters only
  in the global statement Theorem~\ref{thm:flat_base_change_pushforward}, where it is
  used to commute \(- \otimes_A B\) with the finite equalizer computing \(H^0\).
\end{proof}

\section{Flat base change for the pushforward}

The global statement is assembled from the affine lemma and the single homological
input that flatness preserves the finite equalizer computing \(H^0\). That input is
supplied verbatim by Mathlib, recorded here as a dependency anchor.

\begin{lemma}[Flat base change preserves finite equalizers]
  \label{lem:flat_preserves_equalizer_mathlib}
  \lean{LinearMap.tensorEqLocusEquiv}
  \mathlibok
  \textit{Provided by Mathlib.}
  Let \(A \to B\) be a flat ring homomorphism (more generally, let \(B\) be a flat
  \(A\)-module that is also an \(A\)-algebra, with the scalar-tower compatibility),
  and let \(f, g : N \rightrightarrows P\) be a pair of \(A\)-linear maps. Write
  \(\operatorname{eq}(f, g) = \{x \in N : f(x) = g(x)\}\) for their equalizer. Then the
  canonical \(B\)-linear comparison map
  \[
    B \otimes_A \operatorname{eq}(f, g)
    \;\xrightarrow{\ \sim\ }\;
    \operatorname{eq}\bigl(\mathbf{1}_B \otimes f,\ \mathbf{1}_B \otimes g\bigr)
  \]
  is an isomorphism; that is, base change along a flat algebra commutes with the
  formation of the finite equalizer (equivalently, of the kernel: the case
  \(g = 0\) is \texttt{LinearMap.tensorKerEquiv}, underlain by
  \texttt{Module.Flat.ker\_lTensor\_eq} and \texttt{Module.Flat.eqLocus\_lTensor\_eq}).
  Since \(- \otimes_A B\) also preserves finite products, it follows that
  \(- \otimes_A B\) carries the degree-\(0\)/\(1\) equalizer
  \(\prod_i \Gamma(U_i, \mathcal{F}) \rightrightarrows
    \prod_{i,j} \Gamma(U_{ij}, \mathcal{F})\) of a finite cover to the corresponding
  equalizer of the base-changed cover.
\end{lemma}

\begin{lemma}[Sheaf condition as an equalizer of products]
  \label{lem:sheaf_equalizer_products_mathlib}
  \lean{TopCat.Presheaf.isSheaf_iff_isSheafEqualizerProducts}
  \mathlibok
  \textit{Provided by Mathlib.}
  Let \(\mathcal{F}\) be a presheaf with values in a category with products on a
  topological space, and let \(\mathcal{U} = \{U_i\}_{i \in \iota}\) be a family of
  opens with union \(U = \bigcup_i U_i\). Then \(\mathcal{F}\) satisfies the sheaf
  condition on \(\mathcal{U}\) if and only if the fork
  \[
    \mathcal{F}(U)
    \;\xrightarrow{\;\mathrm{res}\;}\;
    \prod_i \mathcal{F}(U_i)
    \;\underset{\mathrm{rightRes}}{\overset{\mathrm{leftRes}}{\rightrightarrows}}\;
    \prod_{(i,j)} \mathcal{F}(U_i \cap U_j)
  \]
  is a limit (an equalizer), where the two parallel maps restrict the \(U_i\)- and
  \(U_j\)-components to \(U_i \cap U_j\). This is Mathlib's
  \texttt{isSheaf\_iff\_isSheafEqualizerProducts} together with the fork
  \texttt{SheafConditionEqualizerProducts.fork}; it is exactly Čech degree
  \(0 \to 1\) and requires no separatedness.
\end{lemma}

\begin{lemma}\leanok
[Finite affine cover with quasi-compact overlaps]
  \label{lem:finite_affine_cover_qcqs}
  \lean{AlgebraicGeometry.Scheme.exists_finite_affineCover_inter_isQuasiCompact}
  A quasi-compact scheme \(X\) admits a \emph{finite} affine open cover
  \(\mathcal{U} : X = U_1 \cup \dots \cup U_t\). If moreover \(X\) is
  quasi-separated, then each pairwise intersection \(U_{ij} = U_i \cap U_j\) is
  quasi-compact; and if \(X\) is separated, each \(U_{ij}\) is in fact affine.
  \begin{proof}
    The affine opens form a basis, so \(X\) is covered by affine opens; quasi-compactness
    extracts a finite subcover \(U_1, \dots, U_t\) (by Mathlib's affine cover and the
    finite-subcover property of a quasi-compact space). Quasi-separatedness says the
    intersection of two quasi-compact opens is quasi-compact, so each \(U_{ij}\) is
    quasi-compact. When \(X\) is separated the diagonal is a closed immersion, whence
    the intersection of two affine opens is affine (Mathlib, for separated schemes).
  \end{proof}
\end{lemma}

\begin{lemma}\leanok
[$\Gamma(X,\mathcal{F})$ as a finite sheaf-condition equalizer]
  \label{lem:gamma_finite_equalizer}
  \lean{AlgebraicGeometry.Modules.gammaIsLimitSheafConditionFork}
  % NOTE: the pinned Lean decl is strictly more general than this prose — it takes ANY
  % open family `U : ι → X.Opens` over ANY scheme (no QCQS / finite / affine hypotheses), giving
  % `IsLimit (fork M.presheaf U)`. The finite affine cover of a QCQS scheme is the intended
  % application; the consolidated existential `Modules.exists_finite_affineCover_isLimit_sheafConditionFork`
  % now has its own block, Lemma~\ref{lem:gamma_finite_equalizer_cover}.
  \uses{lem:sheaf_equalizer_products_mathlib, lem:finite_affine_cover_qcqs}
  Let \(X\) be a quasi-compact, quasi-separated scheme and \(\mathcal{F}\) a sheaf of
  \(\mathcal{O}_X\)-modules. Fix a finite affine cover \(\mathcal{U} = \{U_i\}_{1 \le i \le t}\)
  of \(X\) as in Lemma~\ref{lem:finite_affine_cover_qcqs}. Then the module of global
  sections is the equalizer of the two restriction maps of the cover,
  \[
    \Gamma(X, \mathcal{F})
    \;=\;
    \operatorname{eq}\Bigl(
      \textstyle\prod_{i} \Gamma(U_i, \mathcal{F})
      \;\rightrightarrows\;
      \textstyle\prod_{i,j} \Gamma(U_{ij}, \mathcal{F})
    \Bigr),
    \qquad
    s \mapsto (s_i|_{U_{ij}}), \quad s \mapsto (s_j|_{U_{ij}}),
  \]
  and the products are \emph{finite} (indexed by the finite sets \(\{i\}\) and
  \(\{(i,j)\}\)).
  \begin{proof}
    \uses{lem:sheaf_equalizer_products_mathlib, lem:finite_affine_cover_qcqs}
    Apply the equalizer-products form of the sheaf condition
    (Lemma~\ref{lem:sheaf_equalizer_products_mathlib}) to the underlying sheaf of
    \(\mathcal{F}\) and the cover \(\mathcal{U}\): since \(\bigcup_i U_i = X\), the
    fork displayed there computes \(\mathcal{F}(X) = \Gamma(X, \mathcal{F})\) as the
    equalizer of \(\prod_i \Gamma(U_i, \mathcal{F}) \rightrightarrows \prod_{i,j}
    \Gamma(U_{ij}, \mathcal{F})\). The cover is finite by
    Lemma~\ref{lem:finite_affine_cover_qcqs}, so both products are finite. The
    module structure is inherited objectwise: the forgetful functor to abelian
    groups preserves and reflects the limit, so the equalizer is computed in
    \(\operatorname{ModuleCat}\) over the ground ring.
  \end{proof}
\end{lemma}

\begin{lemma}\leanok
[Global sections as a finite equalizer over a finite affine cover]
  \label{lem:gamma_finite_equalizer_cover}
  \lean{AlgebraicGeometry.Modules.exists_finite_affineCover_isLimit_sheafConditionFork}
  \uses{lem:finite_affine_cover_qcqs, lem:gamma_finite_equalizer}
  Let \(X\) be a quasi-compact, quasi-separated scheme and \(M : X.\mathrm{Modules}\). Then there
  \emph{exists} a finite affine open cover \(\{U_i\}_{i \in \iota}\) of \(X\) — \(\iota\) finite, each
  \(U_i\) affine, \(\bigsqcup_i U_i = \top\), all pairwise intersections \(U_i \cap U_j\)
  quasi-compact — whose sheaf-condition fork of the underlying presheaf of \(M\) is a limit:
  \[
    \Gamma(X, M) \;=\; \operatorname{eq}\Bigl(
      \textstyle\prod_i \Gamma(U_i, M) \;\rightrightarrows\; \prod_{i,j} \Gamma(U_{ij}, M)
    \Bigr).
  \]
  This consolidates the finite affine cover Lemma~\ref{lem:finite_affine_cover_qcqs} with the
  sheaf-condition equalizer Lemma~\ref{lem:gamma_finite_equalizer} into a single existential, the
  directly-usable ``global sections as a finite equalizer'' input of the flat-base-change argument.
  \begin{proof}
    \uses{lem:finite_affine_cover_qcqs, lem:gamma_finite_equalizer}
    Extract a finite affine cover with quasi-compact overlaps by
    Lemma~\ref{lem:finite_affine_cover_qcqs}, and apply the sheaf-condition equalizer
    Lemma~\ref{lem:gamma_finite_equalizer} (in its form valid for any open family) to that cover; the
    pair witnesses the existential.
  \end{proof}
\end{lemma}

% NOTE: the former "A-linear equalizer presentation (FBC-B build-ahead)" subsection
% (3 blocks pinning gammaCoverRestrictScalars/gammaCoverResMapsALinear/gammaEqLocusIso) was RETIRED —
% replaced by the element-level route blocks (gammaModA/gammaResA/leftRes/rightRes/
% gammaTopEquivEqLocus/baseChangeGammaEquiv) in the FBC-B globalization subsection. The retired pins
% named never-built decls and created phantom frontier nodes.

\begin{lemma}[Base change of the finite equalizer diagram]
  \label{lem:base_changed_equalizer_diagram}
  \lean{AlgebraicGeometry.baseChange_sheafConditionFork_tensorIso}
  \uses{lem:gamma_finite_equalizer, lem:affine_base_change_pushforward}
  % SOURCE: [Stacks Project], Tag 02KH ($i=0$ case), separated paragraph
  % (read from references/stacks-coherent.md; quoted verbatim in the
  % SOURCE QUOTE PROOF of thm:flat_base_change_pushforward below).
  % SOURCE QUOTE: "If $\mathcal{U}_B : X_B = (U_1)_B \cup \ldots \cup (U_t)_B$
  % we obtain by base change, then it is still the case that each $(U_i)_B$
  % is affine and that $X_B$ is separated. ... We have the following relation
  % between the {\v C}ech complexes
  % $\check{\mathcal{C}}^\bullet(\mathcal{U}_B, \mathcal{F}_B) =
  % \check{\mathcal{C}}^\bullet(\mathcal{U}, \mathcal{F}) \otimes_A B$
  % as follows from Lemma \ref{lemma-affine-base-change}."
  \textit{Source: Stacks Project, Tag 02KH (``Flat base change''), $i=0$ case,
  the relation \(\check{\mathcal{C}}^\bullet(\mathcal{U}_B,\mathcal{F}_B) =
  \check{\mathcal{C}}^\bullet(\mathcal{U},\mathcal{F}) \otimes_A B\).}
  Let \(A \to B\) be a ring map, \(X\) a separated scheme over \(A\) with finite
  affine cover \(\mathcal{U} = \{U_i\}\) (so each \(U_{ij}\) is affine), and
  \(\mathcal{F}\) quasi-coherent. Set \(X_B = X \times_{\operatorname{Spec} A}
  \operatorname{Spec} B\), \(\mathcal{F}_B\) the pullback, and
  \(\mathcal{U}_B = \{(U_i)_B\}\). Then \(\mathcal{U}_B\) is a finite affine cover of
  \(X_B\), and the sheaf-condition fork of Lemma~\ref{lem:gamma_finite_equalizer} for
  \((X_B, \mathcal{F}_B)\) is obtained from that for \((X, \mathcal{F})\) by applying
  \(- \otimes_A B\): there are compatible isomorphisms
  \[
    \Gamma\bigl((U_i)_B, \mathcal{F}_B\bigr) \cong \Gamma(U_i, \mathcal{F}) \otimes_A B,
    \qquad
    \Gamma\bigl((U_{ij})_B, \mathcal{F}_B\bigr) \cong \Gamma(U_{ij}, \mathcal{F}) \otimes_A B,
  \]
  intertwining the two restriction maps, so that the whole diagram
  \(\prod_i \Gamma((U_i)_B,\mathcal{F}_B) \rightrightarrows \prod_{i,j}
  \Gamma((U_{ij})_B,\mathcal{F}_B)\) is \((- \otimes_A B)\) applied to
  \(\prod_i \Gamma(U_i,\mathcal{F}) \rightrightarrows \prod_{i,j}
  \Gamma(U_{ij},\mathcal{F})\).
  \begin{proof}
    \uses{lem:gamma_finite_equalizer, lem:affine_base_change_pushforward}
    Base change preserves affineness and open immersions, so \((U_i)_B\) is affine and
    \(\{(U_i)_B\}\) covers \(X_B\); since \(X\) is separated each \(U_{ij}\), hence each
    \((U_{ij})_B\), is affine, and the fork of
    Lemma~\ref{lem:gamma_finite_equalizer} applies to \((X_B,\mathcal{F}_B)\). On each
    affine term the global-sections form of the affine lemma
    Lemma~\ref{lem:affine_base_change_pushforward} supplies the isomorphism
    \(\Gamma((U_i)_B,\mathcal{F}_B) \cong \Gamma(U_i,\mathcal{F}) \otimes_A B\)
    (and likewise on the \(U_{ij}\)); these are natural in restriction, so they
    intertwine \(\mathrm{leftRes}\) and \(\mathrm{rightRes}\). Finally
    \(- \otimes_A B\) commutes with the finite products \(\prod_i\), \(\prod_{i,j}\)
    (Mathlib, tensor product preserves finite products), so the base-changed fork is
    literally \((- \otimes_A B)\) applied to the original fork.
  \end{proof}
\end{lemma}

\begin{lemma}[Flat base change, separated case]
  \label{lem:flat_base_change_separated}
  \lean{AlgebraicGeometry.flatBaseChange_pushforward_isIso_of_isSeparated}
  \uses{lem:gamma_finite_equalizer, lem:base_changed_equalizer_diagram,
        lem:flat_preserves_equalizer_mathlib}
  % SOURCE: [Stacks Project], Tag 02KH ($i=0$ case), separated paragraph
  % (read from references/stacks-coherent.md; quoted verbatim in the
  % SOURCE QUOTE PROOF of thm:flat_base_change_pushforward below).
  % SOURCE QUOTE: "Since $A \to B$ is flat, the same thing remains true on
  % taking cohomology."
  \textit{Source: Stacks Project, Tag 02KH (``Flat base change''), $i=0$ case,
  separated paragraph.}
  Let \(A \to B\) be flat, \(X\) a separated, quasi-compact scheme over \(A\), and
  \(\mathcal{F}\) quasi-coherent. Then the comparison map
  \[
    \Gamma(X, \mathcal{F}) \otimes_A B \;\longrightarrow\; \Gamma(X_B, \mathcal{F}_B)
  \]
  is an isomorphism.
  \begin{proof}
    \uses{lem:gamma_finite_equalizer, lem:base_changed_equalizer_diagram,
          lem:flat_preserves_equalizer_mathlib}
    By Lemma~\ref{lem:gamma_finite_equalizer}, \(\Gamma(X,\mathcal{F})\) and
    \(\Gamma(X_B,\mathcal{F}_B)\) are the equalizers of the finite forks of the cover
    \(\mathcal{U}\) and its base change \(\mathcal{U}_B\). By
    Lemma~\ref{lem:base_changed_equalizer_diagram} the \(X_B\)-fork is the
    \(X\)-fork with \(- \otimes_A B\) applied. Since \(A \to B\) is flat,
    \(- \otimes_A B\) commutes with the formation of this finite equalizer
    (Lemma~\ref{lem:flat_preserves_equalizer_mathlib}); hence
    \(\Gamma(X,\mathcal{F}) \otimes_A B = \operatorname{eq}(\cdots) \otimes_A B
    \xrightarrow{\sim} \operatorname{eq}(\cdots \otimes_A B) = \Gamma(X_B,\mathcal{F}_B)\),
    and this iso is the comparison map.
  \end{proof}
\end{lemma}

\begin{lemma}[Flat base change, Mayer--Vietoris reduction of the quasi-separated case]
  \label{lem:flat_base_change_mayer_vietoris}
  \lean{AlgebraicGeometry.flatBaseChange_pushforward_mayerVietoris}
  \uses{lem:flat_base_change_separated, lem:affine_base_change_pushforward,
        lem:flat_preserves_equalizer_mathlib, lem:finite_affine_cover_qcqs}
  % SOURCE: [Stacks Project], Tag 02KH ($i=0$ case), quasi-separated paragraph
  % (read from references/stacks-coherent.md; quoted verbatim in the
  % SOURCE QUOTE PROOF of thm:flat_base_change_pushforward below).
  % SOURCE QUOTE: "In case $X$ is quasi-separated, choose an affine open
  % covering $\mathcal{U} : X = U_1 \cup \ldots \cup U_t$. ... The reader who
  % wishes to avoid this spectral sequence can use Mayer-Vietoris and induction
  % on $t$ as in the proof of Lemma
  % \ref{lemma-quasi-coherence-higher-direct-images}."
  \textit{Source: Stacks Project, Tag 02KH (``Flat base change''), $i=0$ case,
  quasi-separated paragraph (Mayer--Vietoris induction on \(t\)).}
  Let \(A \to B\) be flat, \(X\) a quasi-compact, quasi-separated scheme over \(A\),
  and \(\mathcal{F}\) quasi-coherent. Then the comparison map
  \(\Gamma(X,\mathcal{F}) \otimes_A B \to \Gamma(X_B,\mathcal{F}_B)\) is an
  isomorphism.
  \begin{proof}
    \uses{lem:flat_base_change_separated, lem:affine_base_change_pushforward,
          lem:flat_preserves_equalizer_mathlib, lem:finite_affine_cover_qcqs}
    Choose a finite affine cover \(\mathcal{U} : X = U_1 \cup \dots \cup U_t\)
    (Lemma~\ref{lem:finite_affine_cover_qcqs}) and induct on \(t\). For \(t = 1\),
    \(X = U_1\) is affine and the claim is the affine lemma
    Lemma~\ref{lem:affine_base_change_pushforward}. For \(t > 1\) put \(V = U_t\)
    (affine) and \(W = U_1 \cup \dots \cup U_{t-1}\) (covered by \(t-1\) affines, hence
    quasi-compact and quasi-separated). Their intersection \(W \cap V =
    \bigcup_{i<t}(U_i \cap U_t)\) is covered by the \(t-1\) opens \(U_i \cap U_t\), each
    an open of the affine \(U_i\) and so quasi-compact and separated; thus \(W \cap V\)
    is separated and quasi-compact, and flat base change holds for it by the separated
    case Lemma~\ref{lem:flat_base_change_separated}. The two-member cover
    \(\{W, V\}\) of \(X\) gives the left-exact Mayer--Vietoris equalizer
    \[
      0 \to \Gamma(X,\mathcal{F}) \to
        \Gamma(W,\mathcal{F}) \oplus \Gamma(V,\mathcal{F}) \to \Gamma(W\cap V,\mathcal{F}),
    \]
    and likewise for \(X_B\) with \(\{W_B, V_B\}\). Flatness makes \(- \otimes_A B\)
    preserve this finite kernel/equalizer
    (Lemma~\ref{lem:flat_preserves_equalizer_mathlib}); since flat base change holds for
    \(V\) (affine), for \(W\) (inductive hypothesis, \(t-1\) members), and for
    \(W \cap V\) (separated), the five-lemma/equalizer comparison forces it for \(X\).
  \end{proof}
\end{lemma}

\begin{lemma}[Reduction of the base-change map to the global-sections comparison]
  \label{lem:flat_base_change_reduce_global_sections}
  \lean{AlgebraicGeometry.flatBaseChange_isIso_iff_gammaTensorComparison}
  \uses{def:pushforward_base_change_map}
  % SOURCE: [Stacks Project], Tag 02KH (the reduction "part (1) follows from
  % part (2)") (read from references/stacks-coherent.md; quoted verbatim in
  % the SOURCE QUOTE PROOF of thm:flat_base_change_pushforward below).
  % SOURCE QUOTE: "since $R^if_*\mathcal{F}$ is quasi-coherent ..., it is the
  % quasi-coherent $\mathcal{O}_S$-module associated to the $A$-module
  % $H^0(S, R^if_*\mathcal{F}) = H^i(X, \mathcal{F})$ ... Since pullback by $g$
  % corresponds to $- \otimes_A B$ on modules ... we see that it suffices to
  % prove (2)."
  \textit{Source: Stacks Project, Tag 02KH (``Flat base change''), the reduction
  ``part (1) follows from part (2)''.}
  In the cartesian square with \(g\) flat and \(f\) quasi-compact and
  quasi-separated, the sheaf-level base-change map
  \(g^*(f_*\mathcal{F}) \to f'_*(g')^*\mathcal{F}\) of
  Definition~\ref{def:pushforward_base_change_map} is an isomorphism if and only if,
  after reducing to the affine bases \(S = \operatorname{Spec} A\),
  \(S' = \operatorname{Spec} B\), the module comparison map
  \(\Gamma(X,\mathcal{F}) \otimes_A B \to \Gamma(X_B, \mathcal{F}_B)\) is an
  isomorphism.
  \begin{proof}
    \uses{def:pushforward_base_change_map}
    Being an isomorphism is local on \(S'\), so we may take \(S = \operatorname{Spec} A\)
    and \(S' = \operatorname{Spec} B\) with \(A \to B\) flat. Because \(f\) is
    quasi-compact and quasi-separated, \(f_*\mathcal{F}\) is quasi-coherent on \(S\)
    (Stacks Tag 01XJ / Mathlib's quasi-coherence of the pushforward under a
    quasi-compact quasi-separated morphism), hence is the sheaf associated to the
    \(A\)-module \(\Gamma(S, f_*\mathcal{F}) = \Gamma(X,\mathcal{F})\); likewise
    \(f'_*\mathcal{F}'\) is associated to \(\Gamma(X_B,\mathcal{F}_B)\). Pullback along
    \(g\) corresponds to \(- \otimes_A B\) on modules, so under the
    tilde-equivalence the sheaf base-change map is the tilde of the comparison map
    \(\Gamma(X,\mathcal{F}) \otimes_A B \to \Gamma(X_B,\mathcal{F}_B)\); since
    \(\widetilde{(-)}\) is fully faithful on quasi-coherent modules, one is an
    isomorphism iff the other is.
  \end{proof}
\end{lemma}

\begin{theorem}\leanok
[Flat base change, $i=0$ case]
  \label{thm:flat_base_change_pushforward}
  \lean{AlgebraicGeometry.flatBaseChange_pushforward_isIso}
  \uses{def:pushforward_base_change_map, lem:affine_base_change_pushforward,
        lem:flat_preserves_equalizer_mathlib,
        lem:flat_base_change_reduce_global_sections,
        lem:flat_base_change_mayer_vietoris}
  % SOURCE: [Stacks Project], Tag 02KH (Lemma "Flat base change", part for
  % $i = 0$) (read from references/stacks-coherent.tex, L947--970, and
  % references/stacks-coherent.md).
  % SOURCE QUOTE: "Consider a cartesian diagram of schemes
  % $$
  % \xymatrix{
  % X' \ar[d]_{f'} \ar[r]_{g'} & X \ar[d]^f \\
  % S' \ar[r]^g & S
  % }
  % $$
  % Let $\mathcal{F}$ be a quasi-coherent $\mathcal{O}_X$-module
  % with pullback $\mathcal{F}' = (g')^*\mathcal{F}$.
  % Assume that $g$ is flat and that $f$ is quasi-compact and quasi-separated.
  % For any $i \geq 0$
  % \begin{enumerate}
  % \item the base change map of
  % Cohomology, Lemma \ref{cohomology-lemma-base-change-map-flat-case}
  % is an isomorphism
  % $$
  % g^*R^if_*\mathcal{F} \longrightarrow R^if'_*\mathcal{F}',
  % $$
  % \item if $S = \Spec(A)$ and $S' = \Spec(B)$, then
  % $H^i(X, \mathcal{F}) \otimes_A B = H^i(X', \mathcal{F}')$.
  % \end{enumerate}"
  \textit{Source: Stacks Project, Tag 02KH (``Flat base change''), the
  \(i = 0\) case.}
  Consider the cartesian square
  \[
    \begin{array}{ccc}
      X' & \xrightarrow{\;g'\;} & X \\[2pt]
      {\scriptstyle f'}\big\downarrow & & \big\downarrow{\scriptstyle f} \\[2pt]
      S' & \xrightarrow{\;g\;} & S
    \end{array}
  \]
  with \(\mathcal{F}\) a quasi-coherent \(\mathcal{O}_X\)-module and
  \(\mathcal{F}' = (g')^*\mathcal{F}\). Assume that \(g\) is flat and that \(f\)
  is quasi-compact and quasi-separated. Then the base-change map of
  Definition~\ref{def:pushforward_base_change_map} is an isomorphism
  \[
    g^*\bigl(f_*\mathcal{F}\bigr) \;\xrightarrow{\;\sim\;}\;
    f'_*\bigl((g')^*\mathcal{F}\bigr) = f'_*\mathcal{F}' .
  \]
  Equivalently, in the affine situation \(S = \operatorname{Spec}(A)\),
  \(S' = \operatorname{Spec}(B)\) with \(A \to B\) flat, the comparison map
  \[
    H^0(X, \mathcal{F}) \otimes_A B \;\xrightarrow{\;\sim\;}\;
    H^0(X_B, \mathcal{F}_B)
  \]
  is an isomorphism, where \(X_B = X \times_{\operatorname{Spec}(A)}
  \operatorname{Spec}(B)\) and \(\mathcal{F}_B\) is the pullback of
  \(\mathcal{F}\).

  The hypothesis that \(\mathcal{F}\) is quasi-coherent is essential: it is what
  makes \(f_*\mathcal{F}\) quasi-coherent on the affine base and thus the tilde of
  a module, and it is what lets the affine reduction
  (Lemma~\ref{lem:affine_base_change_pushforward}) write \(\mathcal{F}\) as
  \(\widetilde{M}\) over each affine open. Without it the comparison map need not be
  an isomorphism.
\end{theorem}

% SOURCE QUOTE PROOF: "Having said this, part (1) follows from part (2). Namely,
% part (1) is local on $S'$ and hence we may assume $S$
% and $S'$ are affine. In other words, we have $S = \Spec(A)$
% and $S' = \Spec(B)$ as in (2).
% Then since $R^if_*\mathcal{F}$ is quasi-coherent
% (Lemma \ref{lemma-quasi-coherence-higher-direct-images}),
% it is the quasi-coherent $\mathcal{O}_S$-module associated to the
% $A$-module $H^0(S, R^if_*\mathcal{F}) = H^i(X, \mathcal{F})$
% (equality by
% Lemma \ref{lemma-quasi-coherence-higher-direct-images-application}).
% Similarly, $R^if'_*\mathcal{F}'$ is the quasi-coherent
% $\mathcal{O}_{S'}$-module associated to the $B$-module
% $H^i(X', \mathcal{F}')$. Since pullback by $g$ corresponds
% to $- \otimes_A B$ on modules
% (Schemes, Lemma \ref{schemes-lemma-widetilde-pullback})
% we see that it suffices to prove (2).
%
% Let $A \to B$ be a flat ring homomorphism.
% Let $X$ be a quasi-compact and quasi-separated scheme over $A$.
% Let $\mathcal{F}$ be a quasi-coherent $\mathcal{O}_X$-module.
% Set $X_B = X \times_{\Spec(A)} \Spec(B)$ and denote
% $\mathcal{F}_B$ the pullback of $\mathcal{F}$.
% We are trying to show that the map
% $$
% H^i(X, \mathcal{F}) \otimes_A B \longrightarrow H^i(X_B, \mathcal{F}_B)
% $$
% (given by the reference in the statement of the lemma)
% is an isomorphism.
%
% In case $X$ is separated, choose an affine open covering
% $\mathcal{U} : X = U_1 \cup \ldots \cup U_t$ and recall that
% $$
% \check{H}^p(\mathcal{U}, \mathcal{F}) = H^p(X, \mathcal{F}),
% $$
% see
% Lemma \ref{lemma-cech-cohomology-quasi-coherent}.
% If $\mathcal{U}_B : X_B = (U_1)_B \cup \ldots \cup (U_t)_B$ we obtain
% by base change, then it is still the case that each $(U_i)_B$ is affine
% and that $X_B$ is separated. Thus we obtain
% $$
% \check{H}^p(\mathcal{U}_B, \mathcal{F}_B) = H^p(X_B, \mathcal{F}_B).
% $$
% We have the following relation between the {\v C}ech complexes
% $$
% \check{\mathcal{C}}^\bullet(\mathcal{U}_B, \mathcal{F}_B) =
% \check{\mathcal{C}}^\bullet(\mathcal{U}, \mathcal{F}) \otimes_A B
% $$
% as follows from
% Lemma \ref{lemma-affine-base-change}.
% Since $A \to B$ is flat, the same thing remains true on taking cohomology.
%
% In case $X$ is quasi-separated, choose an affine open covering
% $\mathcal{U} : X = U_1 \cup \ldots \cup U_t$. We will use the
% {\v C}ech-to-cohomology spectral sequence
% Cohomology, Lemma \ref{cohomology-lemma-cech-spectral-sequence}.
% The reader who wishes to avoid this spectral sequence
% can use Mayer-Vietoris and induction on $t$ as in the proof of
% Lemma \ref{lemma-quasi-coherence-higher-direct-images}."
% NOTE: For the $i = 0$ case the argument is Čech-free: $H^0(X,\mathcal{F})$ is the
% sheaf-condition equalizer of a finite affine cover (Čech degree 0/1, NOT Čech
% cohomology -- no Leray, no spectral sequence), each term base-changes by the affine
% lemma lem:affine_base_change_pushforward, and flatness of $A \to B$ commutes
% $-\otimes_A B$ with that finite equalizer via lem:flat_preserves_equalizer_mathlib
% (Mathlib's LinearMap.tensorEqLocusEquiv). The quasi-separated case is the
% Mayer--Vietoris induction on the number of cover members named in the source quote.
\begin{proof}
  \uses{def:pushforward_base_change_map, lem:affine_base_change_pushforward,
        lem:flat_preserves_equalizer_mathlib,
        lem:flat_base_change_reduce_global_sections,
        lem:flat_base_change_mayer_vietoris}
  For \(i = 0\) the argument is {\v C}ech-\emph{free}: it uses only the sheaf
  condition (the degree-\(0\)/\(1\) part of {\v C}ech, i.e.\ an equalizer of a finite
  cover), the affine lemma, and the fact that flat base change commutes with that
  finite equalizer. There is no Leray sequence and no {\v C}ech-to-cohomology spectral
  sequence. The proof is the composite of the decomposition chain.

  \medskip
  \emph{Reduction to the global-sections comparison.} By
  Lemma~\ref{lem:flat_base_change_reduce_global_sections}, the base-change map of
  Definition~\ref{def:pushforward_base_change_map} is an isomorphism if and only if,
  after reducing to affine bases \(S = \operatorname{Spec}(A)\),
  \(S' = \operatorname{Spec}(B)\) with \(A \to B\) flat, the module comparison map
  \[
    \Gamma(X, \mathcal{F}) \otimes_A B \longrightarrow \Gamma(X_B, \mathcal{F}_B),
    \qquad X_B = X \times_{\operatorname{Spec}(A)} \operatorname{Spec}(B),
  \]
  is an isomorphism (this uses that \(f_*\mathcal{F}\) is quasi-coherent — \(f\) is
  quasi-compact and quasi-separated — so each side is the tilde of its module of
  global sections, and that pullback along \(g\) is \(- \otimes_A B\) on modules).

  \medskip
  \emph{The comparison map is an isomorphism.} Since \(f\) is quasi-compact and
  quasi-separated, \(X\) is quasi-compact and quasi-separated over \(A\), so
  Lemma~\ref{lem:flat_base_change_mayer_vietoris} applies and gives
  \[
    \Gamma(X, \mathcal{F}) \otimes_A B \;\xrightarrow{\ \sim\ }\; \Gamma(X_B, \mathcal{F}_B).
  \]
  Internally that lemma presents \(\Gamma(X,\mathcal{F})\) as the finite
  sheaf-condition equalizer of a finite affine cover
  (Lemma~\ref{lem:gamma_finite_equalizer}), base changes the equalizer diagram term by
  term via the affine lemma Lemma~\ref{lem:affine_base_change_pushforward}
  (Lemma~\ref{lem:base_changed_equalizer_diagram}), and commutes \(- \otimes_A B\)
  past that finite equalizer by flatness
  (Lemma~\ref{lem:flat_preserves_equalizer_mathlib}), reducing the general
  quasi-separated case to the separated case
  (Lemma~\ref{lem:flat_base_change_separated}) by Mayer--Vietoris induction on the
  number of cover members. Combined with the reduction above, the sheaf-level
  base-change map is an isomorphism.
\end{proof}

\subsection{The \texorpdfstring{$A$}{A}-module \texorpdfstring{$\operatorname{eqLocus}$}{eqLocus} globalization (FBC-B)}

The separated-case argument is repackaged here at the level of \(A\)-modules, where
\(A = \Gamma(X,\mathcal{O}_X)\) is the ring of global sections of the structure sheaf, so that flat
base change becomes the Mathlib equivalence \(\operatorname{tensorEqLocusEquiv}\)
(Lemma~\ref{lem:flat_preserves_equalizer_mathlib}) applied to an honest
\(\operatorname{LinearMap.eqLocus}\). The blocks below construct that presentation concretely: every
module of sections \(\Gamma(U, M)\) is viewed as an \(A\)-module by restriction of scalars along the
structure-sheaf restriction \(A \to \Gamma(U,\mathcal{O}_X)\), the restriction maps of \(M\) become
\(A\)-linear, and the two product legs of the sheaf condition assemble into a pair of \(A\)-linear maps
whose equalizer locus is \(\Gamma(X, M)\). Throughout, \(M : X.\mathrm{Modules}\) is a sheaf of
\(\mathcal{O}_X\)-modules and \(U : \iota \to \mathrm{Opens}(X)\) is an open cover with
\(\bigvee_i U_i = \top\).

The bijectivity of the keystone presentation rests on two element-level sheaf axioms of the underlying
abelian-group sheaf \(M.\mathrm{presheaf}\), recorded as Mathlib dependency anchors.

\begin{lemma}[Sheaf separatedness: local equality forces global equality]
  \label{lem:sheaf_eq_of_locally_eq_mathlib}
  \lean{TopCat.Sheaf.eq_of_locally_eq'}
  \mathlibok
  \textit{Provided by Mathlib.}
  Let \(\mathcal{G}\) be a sheaf on a topological space \(X\), \(U : \iota \to \mathrm{Opens}(X)\) a
  family covering an open \(V\) (i.e.\ \(V \le \bigvee_i U_i\)), and \(s, t \in \mathcal{G}(V)\) two
  sections. If the restrictions of \(s\) and \(t\) to every \(U_i\) agree, then \(s = t\). This is the
  separatedness half of the sheaf condition, stated element-wise.
\end{lemma}

\begin{lemma}[Sheaf gluing: a compatible family has a unique amalgamation]
  \label{lem:sheaf_existsUnique_gluing_mathlib}
  \lean{TopCat.Sheaf.existsUnique_gluing'}
  \mathlibok
  \textit{Provided by Mathlib.}
  Let \(\mathcal{G}\) be a sheaf on \(X\), \(U : \iota \to \mathrm{Opens}(X)\) a family covering an open
  \(V\), and \((s_i)_i\) a family of sections \(s_i \in \mathcal{G}(U_i)\) that is \emph{compatible}
  (the restrictions of \(s_i\) and \(s_j\) to \(U_i \cap U_j\) agree for all \(i, j\)). Then there is a
  unique \(s \in \mathcal{G}(V)\) whose restriction to each \(U_i\) is \(s_i\). This is the gluing half
  of the sheaf condition, stated element-wise.
\end{lemma}

\begin{definition}\leanok
[The ground ring of global sections]
  \label{def:fbcb_groundRing}
  \lean{AlgebraicGeometry.groundRing}
  For a scheme \(X\), the \emph{ground ring} is the commutative ring
  \(A = \Gamma(X, \mathcal{O}_X) = \mathcal{O}_X(\top)\) of global sections of the structure sheaf,
  taken from the \(\mathbf{CommRing}\)-valued structure presheaf so its ring and algebra structures
  resolve canonically.
\end{definition}

\begin{definition}\leanok
[The structure-sheaf restriction ring map \texorpdfstring{$A \to \Gamma(U,\mathcal{O}_X)$}{A to sections over U}]
  \label{def:fbcb_rhoU}
  \lean{AlgebraicGeometry.rhoU}
  \uses{def:fbcb_groundRing}
  For an open \(U \subseteq X\), the structure-sheaf restriction along \(U \subseteq \top\) is a ring
  homomorphism \(\rho_U : A = \Gamma(X,\mathcal{O}_X) \to \Gamma(U,\mathcal{O}_X)\). It is the ring map
  along which sections over \(U\) are regarded as \(A\)-modules.
\end{definition}

\begin{lemma}\leanok
[Transitivity of the restriction ring maps]
  \label{lem:fbcb_rhoU_comp}
  \lean{AlgebraicGeometry.rhoU_comp}
  \uses{def:fbcb_rhoU}
  For opens \(V \le U \subseteq X\), the structure-sheaf restriction
  \(\Gamma(U,\mathcal{O}_X) \to \Gamma(V,\mathcal{O}_X)\) composed with \(\rho_U\) equals \(\rho_V\):
  \[
    \bigl(\Gamma(U,\mathcal{O}_X) \to \Gamma(V,\mathcal{O}_X)\bigr) \circ \rho_U = \rho_V .
  \]
  \begin{proof}
    \uses{def:fbcb_rhoU}
    Both sides are the structure-sheaf restriction \(\Gamma(X,\mathcal{O}_X) \to \Gamma(V,\mathcal{O}_X)\)
    along \(V \subseteq \top\), factored as \(V \subseteq U \subseteq \top\); functoriality of the
    structure presheaf collapses the two-step restriction to the one-step restriction.
  \end{proof}
\end{lemma}

\begin{definition}\leanok
[Sections over \texorpdfstring{$U$}{U} as an \texorpdfstring{$A$}{A}-module]
  \label{def:fbcb_gammaModA}
  \lean{AlgebraicGeometry.gammaModA}
  \uses{def:fbcb_rhoU}
  The sections \(\Gamma(U, M)\), a module over \(\Gamma(U,\mathcal{O}_X)\), are regarded as an
  \(A\)-module by restriction of scalars along \(\rho_U\) (Definition~\ref{def:fbcb_rhoU}). This is the
  common \(A\)-module home in which the global-sections equalizer is formed.
\end{definition}

\begin{definition}\leanok
[The restriction \texorpdfstring{$\Gamma(U,M) \to \Gamma(V,M)$}{Gamma(U,M) to Gamma(V,M)} as an \texorpdfstring{$A$}{A}-module morphism]
  \label{def:fbcb_gammaResAHom}
  \lean{AlgebraicGeometry.gammaResAHom}
  \uses{def:fbcb_gammaModA, lem:fbcb_rhoU_comp}
  For \(V \le U\), the structure-sheaf restriction \(M.\mathrm{map} : \Gamma(U,M) \to \Gamma(V,M)\) is
  promoted to a morphism of \(A\)-modules \(\Gamma(U,M) \to \Gamma(V,M)\) (Definition~\ref{def:fbcb_gammaModA}).
  It is built from \(M.\mathrm{map}\) by restriction of scalars, using the transitivity
  \(\rho\)-coherence Lemma~\ref{lem:fbcb_rhoU_comp} to identify the two scalar actions.
\end{definition}

\begin{definition}\leanok
[The \texorpdfstring{$A$}{A}-linear restriction map]
  \label{def:fbcb_gammaResA}
  \lean{AlgebraicGeometry.gammaResA}
  \uses{def:fbcb_gammaResAHom}
  The underlying \(A\)-linear map \(\Gamma(U,M) \to_A \Gamma(V,M)\) of the \(A\)-module morphism of
  Definition~\ref{def:fbcb_gammaResAHom}. It is the building block of the two product legs of the
  equalizer.
\end{definition}

\begin{lemma}\leanok
[The \texorpdfstring{$A$}{A}-linear restriction is the underlying presheaf restriction]
  \label{lem:fbcb_gammaResA_apply}
  \lean{AlgebraicGeometry.gammaResA_apply}
  \uses{def:fbcb_gammaResA}
  On underlying elements, the \(A\)-linear restriction map of Definition~\ref{def:fbcb_gammaResA} acts
  as the presheaf restriction of \(M\): for \(x \in \Gamma(U,M)\) and \(V \le U\),
  \(\mathrm{gammaResA}(x) = (M.\mathrm{map})\,x\).
  \begin{proof}
    \uses{def:fbcb_gammaResA}
    Restriction of scalars and the composition with the comparison isomorphism of
    Definition~\ref{def:fbcb_gammaResAHom} leave the underlying function unchanged, so the map agrees
    with \(M.\mathrm{map}\) on elements.
  \end{proof}
\end{lemma}

\begin{lemma}\leanok
[Functoriality of the \texorpdfstring{$A$}{A}-linear restriction maps]
  \label{lem:fbcb_gammaResA_comp}
  \lean{AlgebraicGeometry.gammaResA_comp}
  \uses{def:fbcb_gammaResA, lem:fbcb_gammaResA_apply}
  For opens \(W \le V \le U\) and \(x \in \Gamma(U,M)\), the \(A\)-linear restrictions compose as the
  single restriction:
  \[
    \mathrm{gammaResA}_{V \le U}\!\bigl(x\bigr) \text{ then } \mathrm{gammaResA}_{W \le V}
      \;=\; \mathrm{gammaResA}_{W \le U}\!\bigl(x\bigr).
  \]
  \begin{proof}
    \uses{def:fbcb_gammaResA, lem:fbcb_gammaResA_apply}
    By Lemma~\ref{lem:fbcb_gammaResA_apply} each \(A\)-linear restriction is the underlying presheaf
    restriction of \(M\); functoriality of the presheaf \(M\) collapses the composite of the two
    restrictions \(W \subseteq V \subseteq U\) to the single restriction \(W \subseteq U\).
  \end{proof}
\end{lemma}

\begin{definition}\leanok
[The left restriction leg]
  \label{def:fbcb_leftRes}
  \lean{AlgebraicGeometry.leftRes}
  \uses{def:fbcb_gammaResA}
  The \(A\)-linear map
  \[
    \mathrm{leftRes} : \prod_i \Gamma(U_i, M) \longrightarrow \prod_{(i,j)} \Gamma(U_i \cap U_j, M),
  \]
  whose \((i,j)\)-component is the \(A\)-linear restriction \(\Gamma(U_i, M) \to \Gamma(U_i \cap U_j, M)\)
  of Definition~\ref{def:fbcb_gammaResA} applied to the \(i\)-th factor (restriction along
  \(U_i \cap U_j \subseteq U_i\)).
\end{definition}

\begin{definition}\leanok
[The right restriction leg]
  \label{def:fbcb_rightRes}
  \lean{AlgebraicGeometry.rightRes}
  \uses{def:fbcb_gammaResA}
  The \(A\)-linear map \(\mathrm{rightRes} : \prod_i \Gamma(U_i, M) \to \prod_{(i,j)} \Gamma(U_i \cap U_j, M)\)
  whose \((i,j)\)-component is the \(A\)-linear restriction of Definition~\ref{def:fbcb_gammaResA} applied
  to the \(j\)-th factor (restriction along \(U_i \cap U_j \subseteq U_j\)).
\end{definition}

\begin{definition}\leanok
[Restriction of a global section to the cover]
  \label{def:fbcb_toCover}
  \lean{AlgebraicGeometry.toCover}
  \uses{def:fbcb_gammaResA}
  The \(A\)-linear map \(\mathrm{toCover} : \Gamma(\top, M) \to \prod_i \Gamma(U_i, M)\) restricting a
  global section to each member of the cover, with \(i\)-th component the \(A\)-linear restriction
  \(\Gamma(\top, M) \to \Gamma(U_i, M)\) of Definition~\ref{def:fbcb_gammaResA}.
\end{definition}

\begin{lemma}\leanok
[The restricted family is compatible]
  \label{lem:fbcb_leftRes_toCover}
  \lean{AlgebraicGeometry.leftRes_toCover}
  \uses{def:fbcb_leftRes, def:fbcb_rightRes, def:fbcb_toCover, lem:fbcb_gammaResA_comp}
  For every global section \(s \in \Gamma(\top, M)\) the two legs agree on the restricted family:
  \[
    \mathrm{leftRes}\bigl(\mathrm{toCover}(s)\bigr) = \mathrm{rightRes}\bigl(\mathrm{toCover}(s)\bigr).
  \]
  \begin{proof}
    \uses{def:fbcb_leftRes, def:fbcb_rightRes, def:fbcb_toCover, lem:fbcb_gammaResA_comp}
    Both sides are, in the \((i,j)\)-component, the restriction of \(s\) from \(\top\) to
    \(U_i \cap U_j\): the left leg restricts \(\top \to U_i \to U_i \cap U_j\) and the right leg
    restricts \(\top \to U_j \to U_i \cap U_j\), and by the functoriality
    Lemma~\ref{lem:fbcb_gammaResA_comp} each two-step restriction collapses to the one-step restriction
    \(\top \to U_i \cap U_j\). Hence the two components coincide.
  \end{proof}
\end{lemma}

\begin{definition}\leanok
[The corestriction into the equalizer locus]
  \label{def:fbcb_toCoverEqLocus}
  \lean{AlgebraicGeometry.toCoverEqLocus}
  \uses{def:fbcb_toCover, lem:fbcb_leftRes_toCover}
  By the compatibility Lemma~\ref{lem:fbcb_leftRes_toCover}, the map \(\mathrm{toCover}\) corestricts to
  an \(A\)-linear map
  \[
    \mathrm{toCoverEqLocus} : \Gamma(\top, M)
      \longrightarrow \operatorname{eqLocus}(\mathrm{leftRes}, \mathrm{rightRes})
      = \{\, t \in \textstyle\prod_i \Gamma(U_i,M) : \mathrm{leftRes}(t) = \mathrm{rightRes}(t)\,\}.
  \]
\end{definition}

\begin{lemma}\leanok
[Global sections as an \texorpdfstring{$A$}{A}-module equalizer locus]
  \label{lem:fbcb_gammaTopEquivEqLocus}
  \lean{AlgebraicGeometry.gammaTopEquivEqLocus}
  \uses{def:fbcb_toCoverEqLocus, def:fbcb_gammaModA, def:fbcb_gammaResA, def:fbcb_leftRes,
        def:fbcb_rightRes, lem:sheaf_eq_of_locally_eq_mathlib, lem:sheaf_existsUnique_gluing_mathlib}
  % NOTE: bijectivity is established by the ELEMENT-LEVEL sheaf axioms
  % (lem:sheaf_eq_of_locally_eq_mathlib for injectivity, lem:sheaf_existsUnique_gluing_mathlib for
  % surjectivity) on the underlying Ab-sheaf M.presheaf; it does NOT route through the
  % sheaf-condition-fork limit (lem:sheaf_equalizer_products_mathlib / gammaIsLimitSheafConditionFork).
  For \(M : X.\mathrm{Modules}\) and any open cover \(U\) of \(X\) (\(\bigvee_i U_i = \top\)), the
  corestriction \(\mathrm{toCoverEqLocus}\) of Definition~\ref{def:fbcb_toCoverEqLocus} is an
  \(A\)-linear isomorphism
  \[
    \Gamma(X, M) = \Gamma(\top, M) \;\cong\; \operatorname{eqLocus}(\mathrm{leftRes}, \mathrm{rightRes}),
  \]
  identifying the global sections, as a module over the ground ring \(A = \Gamma(X,\mathcal{O}_X)\), with
  the submodule on which the two restriction legs agree.
  \begin{proof}
    \uses{def:fbcb_toCoverEqLocus, lem:sheaf_eq_of_locally_eq_mathlib,
          lem:sheaf_existsUnique_gluing_mathlib}
    The map \(\mathrm{toCoverEqLocus}\) is \(A\)-linear by construction, so it suffices to prove it
    bijective on underlying sets. \emph{Injectivity:} if two global sections \(s, t\) have equal images,
    their restrictions to every \(U_i\) agree, so by sheaf separatedness
    (Lemma~\ref{lem:sheaf_eq_of_locally_eq_mathlib}) applied to the underlying abelian-group sheaf
    \(M.\mathrm{presheaf}\), \(s = t\). \emph{Surjectivity:} an element of the locus is a family
    \((s_i)_i\) with \(\mathrm{leftRes} = \mathrm{rightRes}\), i.e.\ a compatible family on the cover; by
    the gluing axiom (Lemma~\ref{lem:sheaf_existsUnique_gluing_mathlib}) it amalgamates to a global
    section restricting to \((s_i)_i\). Hence \(\mathrm{toCoverEqLocus}\) is an \(A\)-linear bijection,
    so an isomorphism.
  \end{proof}
\end{lemma}

\begin{lemma}\leanok
[Flat base change commutes with the \texorpdfstring{$H^0$}{H0} equalizer]
  \label{lem:fbcb_baseChangeGammaEquiv}
  \lean{AlgebraicGeometry.baseChangeGammaEquiv}
  \uses{lem:fbcb_gammaTopEquivEqLocus, lem:flat_preserves_equalizer_mathlib}
  Let \(B\) be a flat \(A\)-algebra (\(A = \Gamma(X,\mathcal{O}_X)\)). Then base changing the global
  sections along \(A \to B\) is the equalizer of the base-changed restriction legs: there is a
  \(B\)-linear isomorphism
  \[
    B \otimes_A \Gamma(X, M) \;\cong\;
      \operatorname{eqLocus}\bigl(\mathrm{id}_B \otimes \mathrm{leftRes},\ \mathrm{id}_B \otimes \mathrm{rightRes}\bigr).
  \]
  \begin{proof}
    \uses{lem:fbcb_gammaTopEquivEqLocus, lem:flat_preserves_equalizer_mathlib}
    Tensoring the keystone isomorphism Lemma~\ref{lem:fbcb_gammaTopEquivEqLocus} by \(B\) over \(A\)
    gives \(B \otimes_A \Gamma(X,M) \cong B \otimes_A \operatorname{eqLocus}(\mathrm{leftRes},\mathrm{rightRes})\),
    and by flatness of \(B\) over \(A\) the Mathlib equivalence \(\operatorname{tensorEqLocusEquiv}\)
    (Lemma~\ref{lem:flat_preserves_equalizer_mathlib}) carries \(B \otimes_A (-)\) of the equalizer to the
    equalizer of the base-changed legs. Composing the two yields the stated \(B\)-linear isomorphism.
    By the same presentation applied to the pulled-back sheaf, the right-hand side is the global sections
    of \(M\) over the base-changed scheme, so this is the module-level core of the \(H^0\)
    flat-base-change comparison.
  \end{proof}
\end{lemma}

\section{FBC-B: the global \texorpdfstring{$H^0$}{H0} flat-base-change isomorphism via the finite affine-cover equalizer (direct route)}

The blocks of the preceding subsection assemble into a \emph{direct} proof of the named global
isomorphism that bypasses the adjoint-mate keystone of the affine/mate route. The route is
\emph{{\v C}ech-cohomology-free}: \(H^0\) is a finite limit — the equalizer of a finite affine cover —
and the flat functor \(-\otimes_A B\) preserves finite limits, so the global comparison isomorphism
follows from per-chart affine base change \emph{without} the abstract conjugate/mate identification and
without the Mayer--Vietoris induction.

Concretely, the module-level core
\(B \otimes_A \Gamma(X,\mathcal{F}) \cong \operatorname{eqLocus}(\mathrm{id}_B\otimes\mathrm{leftRes},
\mathrm{id}_B\otimes\mathrm{rightRes})\) (Lemma~\ref{lem:fbcb_baseChangeGammaEquiv}) already commutes flat
base change past the finite \(H^0\) equalizer of a cover \(\{U_i\}\) of \(X\). The single remaining step
is to recognise the right-hand base-changed equalizer, chart by chart, as the equalizer-locus
presentation (Lemma~\ref{lem:fbcb_gammaTopEquivEqLocus}) of the global sections of the pulled-back module
over the base-changed cover \(\{(U_i)_B\}\) of \(X' = X_B\). The theorem below states that assembly.

\begin{definition}[Ground-ring algebra map of the base-change pullback]
  \label{def:fbcb_pullbackGroundRingAlg}
  \lean{AlgebraicGeometry.Modules.pullbackGroundRingAlg}
  For a flat \(A\)-algebra \(B\) (\(A = \Gamma(X,\OO_X)\)) and the base-change pullback
  \(X' = X \times_{\operatorname{Spec} A} \operatorname{Spec} B\), the projection
  \(f' : X' \to \operatorname{Spec} B\) induces a ring homomorphism
  \(B \to \Gamma(X',\OO_{X'})\), obtained as
  \(f'^{\sharp}_{\top} \circ (\Gamma\!\operatorname{Spec}\text{-iso})^{-1}\). This is the
  map along which \(\Gamma(X',\mathcal{F}')\) is viewed as a \(B\)-module (restriction of
  scalars) in Theorem~\ref{thm:fbcb_global_direct}; chasing \(\Gamma(X',\OO_{X'}) = B\)
  directly is avoided (that equality is the theorem at \(\mathcal{F} = \OO_X\)).
\end{definition}

\begin{theorem}[Global \(H^0\) flat base change via the finite-cover equalizer, direct route]
  \label{thm:fbcb_global_direct}
  \lean{AlgebraicGeometry.Modules.baseChangeGammaPullbackEquiv}
  \uses{lem:finite_affine_cover_qcqs, lem:fbcb_gammaTopEquivEqLocus,
        lem:fbcb_baseChangeGammaEquiv, lem:pullback_spec_tilde_iso,
        lem:affine_base_change_pushforward, lem:flat_base_change_reduce_global_sections,
        def:fbcb_pullbackGroundRingAlg}
  % SOURCE: [Stacks Project], Tag 02KH (``Flat base change''), part (2), $i = 0$ case
  % (read from references/stacks-coherent.tex, L966--968).
  % SOURCE QUOTE: "\item if $S = \Spec(A)$ and $S' = \Spec(B)$, then
  % $H^i(X, \mathcal{F}) \otimes_A B = H^i(X', \mathcal{F}')$."
  \textit{Source: Stacks Project, Tag 02KH (``Flat base change''), part (2), \(i = 0\) case.}
  Let \(X\) be a quasi-compact, quasi-separated scheme with ground ring
  \(A = \Gamma(X,\mathcal{O}_X)\), let \(\mathcal{F}\) be a quasi-coherent \(\mathcal{O}_X\)-module, and
  let \(B\) be a flat \(A\)-algebra. Write \(X' = X_B = X \times_{\operatorname{Spec}(A)}
  \operatorname{Spec}(B)\) with projection \(g' : X' \to X\) and \(\mathcal{F}' = (g')^*\mathcal{F}\) the
  pullback. Then the comparison map
  \[
    \Gamma(X,\mathcal{F}) \otimes_A B \;\xrightarrow{\;\sim\;}\; \Gamma(X', \mathcal{F}')
  \]
  is a \(B\)-linear isomorphism, assembled directly from the finite affine-cover equalizer presentation;
  consequently the sheaf-level base-change map \(g^*(f_*\mathcal{F}) \to f'_*\mathcal{F}'\) is an
  isomorphism, recovering the named global target Theorem~\ref{thm:flat_base_change_pushforward}. The
  route uses neither the adjoint-mate keystone of the conjugate route nor the Mayer--Vietoris induction
  on the number of cover members.

  % SOURCE QUOTE PROOF: "\item We have $\psi^* \widetilde M = \widetilde{S \otimes_R M}$
  % functorially in the $R$-module $M$." (references/stacks-schemes.tex, Tag 01I9, part (1) —
  % the per-chart pullback-is-base-change identification feeding the assembly.)
  \begin{proof}
    \uses{lem:finite_affine_cover_qcqs, lem:fbcb_gammaTopEquivEqLocus,
          lem:fbcb_baseChangeGammaEquiv, lem:pullback_spec_tilde_iso,
          lem:affine_base_change_pushforward, lem:flat_base_change_reduce_global_sections}
    Choose a finite affine cover \(\{U_i\}_{i\in\iota}\) of \(X\) with quasi-compact pairwise overlaps
    (Lemma~\ref{lem:finite_affine_cover_qcqs}), so that \(\bigvee_i U_i = \top\).

    \emph{Flat base change past the finite equalizer.} By
    Lemma~\ref{lem:fbcb_baseChangeGammaEquiv}, flatness of \(B\) over \(A\) gives a \(B\)-linear
    isomorphism
    \[
      B \otimes_A \Gamma(X,\mathcal{F}) \;\cong\;
        \operatorname{eqLocus}\bigl(\mathrm{id}_B\otimes\mathrm{leftRes},\ \mathrm{id}_B\otimes\mathrm{rightRes}\bigr),
    \]
    where \(\mathrm{leftRes}, \mathrm{rightRes} : \prod_i \Gamma(U_i,\mathcal{F}) \rightrightarrows
    \prod_{i,j}\Gamma(U_i\cap U_j,\mathcal{F})\) are the two restriction legs of the cover.

    \emph{Per-chart identification of the base-changed legs.} The base-changed cover
    \(\{(U_i)_B\}_{i\in\iota}\) is again a finite affine cover, now of \(X'\): base change preserves
    affineness and open immersions and carries \(\bigvee_i U_i = \top\) to \(\bigvee_i (U_i)_B = \top\).
    On each chart the affine pullback dictionary (Lemma~\ref{lem:pullback_spec_tilde_iso}, Stacks
    Tag 01I9: \((\operatorname{Spec}\varphi)^*\widetilde M \cong \widetilde{B\otimes_A M}\)) reads the
    sections of \(\mathcal{F}'\) over \((U_i)_B\) as the base change of the sections of \(\mathcal{F}\),
    \[
      \Gamma\bigl((U_i)_B, \mathcal{F}'\bigr) \;\cong\; \Gamma(U_i,\mathcal{F}) \otimes_A B,
    \]
    and likewise on each overlap, where the affine base-change lemma
    (Lemma~\ref{lem:affine_base_change_pushforward}) supplies
    \(\Gamma((U_i\cap U_j)_B, \mathcal{F}') \cong \Gamma(U_i\cap U_j,\mathcal{F})\otimes_A B\). These
    identifications are natural in restriction, so they intertwine the base-changed legs
    \(\mathrm{id}_B\otimes\mathrm{leftRes}\), \(\mathrm{id}_B\otimes\mathrm{rightRes}\) with the
    restriction legs \(\mathrm{leftRes}'\), \(\mathrm{rightRes}'\) of the cover \(\{(U_i)_B\}\) of
    \(X'\); hence they restrict to an isomorphism of equalizer loci
    \[
      \operatorname{eqLocus}\bigl(\mathrm{id}_B\otimes\mathrm{leftRes},\ \mathrm{id}_B\otimes\mathrm{rightRes}\bigr)
      \;\cong\;
      \operatorname{eqLocus}\bigl(\mathrm{leftRes}',\ \mathrm{rightRes}'\bigr).
    \]

    \emph{Recognising the right-hand equalizer.} Applying the equalizer-locus presentation
    (Lemma~\ref{lem:fbcb_gammaTopEquivEqLocus}) to \(\mathcal{F}'\) and the cover \(\{(U_i)_B\}\) of
    \(X'\) identifies the latter equalizer locus with the global sections,
    \(\operatorname{eqLocus}(\mathrm{leftRes}',\mathrm{rightRes}') \cong \Gamma(X',\mathcal{F}')\).
    Composing the three isomorphisms yields the \(B\)-linear isomorphism
    \(\Gamma(X,\mathcal{F})\otimes_A B \cong \Gamma(X',\mathcal{F}')\), which is the comparison map.

    \emph{Back to the sheaf-level statement.} Finally, the reduction
    Lemma~\ref{lem:flat_base_change_reduce_global_sections} — each pushforward is the tilde of its module
    of global sections, and pullback along \(g\) is \(-\otimes_A B\) on modules — turns this module
    isomorphism into the isomorphism of the sheaf-level base-change map
    \(g^*(f_*\mathcal{F}) \to f'_*\mathcal{F}'\), the named global target
    Theorem~\ref{thm:flat_base_change_pushforward}. No conjugate/mate identification and no {\v C}ech
    cohomology beyond the degree \(0/1\) sheaf condition enters.
  \end{proof}
\end{theorem}
